% Rédiger une dissertation — Français Tle Tech — Cours signé OnBuch+
% Programme officiel MINESEC. Compiler avec : tectonic N20-rediger-dissertation.tex
\newcommand{\DOCMATIERE}{Français}
\newcommand{\DOCNIVEAU}{Tle Tech}
\newcommand{\DOCMODULE}{Français – classe de Terminale technique}
\newcommand{\DOCLECON}{N20}
\newcommand{\DOCTITRE}{Rédiger une dissertation}
% OnBuch+ — préambule « cours signés » ENRICHI (compilé avec tectonic / XeLaTeX)
% Système visuel « Pop » d'OnBuch+ : fond crème (jamais blanc), bordures encre
% épaisses, ombres « dures » décalées, titres Archivo Black en capitales.
% Version MATHÉMATIQUES : graphes (pgfplots/tikz), boîtes pédagogiques
% (définition, propriété, méthode, attention, le savais-tu, exercice résolu,
% exercice, TP, bilan), page de garde et signature OnBuch+ ancrée en bas.
%
% Macros attendues dans main.tex AVANT \input{preamble} :
%   \DOCMATIERE  \DOCNIVEAU  \DOCMODULE  \DOCLECON (numéro)  \DOCTITRE
\documentclass[11pt]{article}

\usepackage{fontspec}
\usepackage[french]{babel}
\usepackage[a4paper,margin=2cm,top=3.4cm,bottom=2.8cm,headheight=1.4cm,footskip=1.3cm]{geometry}
\usepackage{amsmath,amssymb,amsfonts}
\usepackage{mathtools} % \xmapsto, \coloneqq... souvent utilisés par les modèles
\usepackage{cancel} % simplifications barrées (bilans de forces)
% \abs{z} (module, valeur absolue) et \norm{u} (norme), utilisés par les modèles
\providecommand{\abs}{}\let\abs\relax\DeclarePairedDelimiter{\abs}{\lvert}{\rvert}
\providecommand{\norm}{}\let\norm\relax\DeclarePairedDelimiter{\norm}{\lVert}{\rVert}
\usepackage[table]{xcolor} % [table] : fournit aussi \rowcolors (alternance de lignes)
\usepackage{pagecolor}
\usepackage{tikz}
\usetikzlibrary{arrows.meta,positioning,calc,shapes.geometric,shapes.misc,shapes.arrows,decorations.pathmorphing,decorations.pathreplacing,decorations.markings,patterns,shadows,mindmap,trees,fit,backgrounds,matrix,intersections,angles,quotes}
\usepackage{pgfplots}
\usepackage[version=4]{mhchem} % équations nucléaires \ce{^{235}_{92}U}
\usepackage[european,siunitx]{circuitikz} % schémas électriques
\pgfplotsset{compat=1.18}
\usepgfplotslibrary{fillbetween,groupplots} % aires entre courbes (calcul intégral)
\usepackage{enumitem}
\usepackage{fancyhdr}
% [most] charge toutes les bibliothèques tcolorbox (listings, minted, raster,
% documentation...) et gonfle inutilement la mémoire TeX sur un document long
% (~300 boîtes) : on ne charge que ce qui est réellement utilisé.
\usepackage[skins,breakable]{tcolorbox}
\usepackage{graphicx}
\usepackage{colortbl}
\usepackage{booktabs}
\usepackage{tabularx}
\usepackage{multirow}
\usepackage{diagbox} % case d'en-tête diagonale des tableaux à double entrée
% Colonnes extensibles centrée / gauche / droite (C, L, R), souvent utilisées par les modèles
\newcolumntype{C}{>{\centering\arraybackslash}X}
\newcolumntype{L}{>{\raggedright\arraybackslash}X}
\newcolumntype{R}{>{\raggedleft\arraybackslash}X}
\usepackage{array}
\usepackage{multicol}
\usepackage{siunitx}
% parse-numbers=false : accepte aussi \SI{2,5 \times 10^{-3}}{...}
\sisetup{locale=FR,output-decimal-marker={,},group-minimum-digits=4,parse-numbers=false}
\DeclareSIUnit\ohm{\ensuremath{\symup{\Omega}}} % U+2126 absent de la police
\DeclareSIUnit\division{div} % réglages d'oscilloscope (ms/div, V/div)
\DeclareSIUnit\tr{tr} % tours (vitesse de rotation, ex. \SI{1500}{\tr\per\minute})
\DeclareSIUnit\molar{M} % concentration molaire (ex. \SI{3}{\molar}, \SI{1}{\micro\molar})
\DeclareSIUnit\an{an} % année (le modèle confond parfois avec \year, un compteur TeX) — ex. \SI{5}{\kilo\meter\per\an}
\DeclareSIUnit\FCFA{FCFA} % franc CFA (ex. \SI{500}{\FCFA\per\kilo\gram})
\DeclareSIUnit\degreeBrix{\textdegree Brix} % teneur en sucre d'un sirop/jus (réfractométrie)
\DeclareSIUnit\month{mois} % le modèle utilise parfois \month, un compteur TeX, comme unité de durée
\DeclareSIUnit\microliter{\micro\liter} % le modèle écrit parfois \microliter en un seul token
\DeclareSIUnit\million{millions} % dénombrement (ex. \SI{15}{\million\per\mL} spermatozoïdes)
\usepackage{qrcode}
% \degree : commande fréquemment utilisée par les modèles pour un angle ou une
% température en dehors de \SI{}{} (non fournie par les paquets chargés).
\providecommand{\degree}{^{\circ}}
% \qed (fin de démonstration), utilisé par les modèles sans amsthm
\providecommand{\qed}{\hfill$\square$}
% \barre : double barre des valeurs interdites dans un tableau de variations (tkz-tab)
\providecommand{\barre}{\ensuremath{\|}}
% environnement proof (amsthm non chargé) utilisé par les modèles
\ifdefined\proof\else\newenvironment{proof}[1][Démonstration]{\par\noindent\textit{#1.}\ }{\qed\par}\fi
\usepackage{lastpage}
\usepackage{hyperref}
\hypersetup{hidelinks}

% ============================================================
% Palette « Pop » OnBuch+ — mêmes hex que la classe P (plus_ui.dart)
% ============================================================
\definecolor{poporange}{HTML}{F59321}
\definecolor{popdark}{HTML}{E07A0C}
\definecolor{popcream}{HTML}{FFF6E8}
\definecolor{popink}{HTML}{1C1714}
\definecolor{poppurple}{HTML}{7A5AE0}
\definecolor{popgreen}{HTML}{2E9E6A}
\definecolor{poppink}{HTML}{E0457B}
\definecolor{popblue}{HTML}{2F7DE1}
\definecolor{popgold}{HTML}{C98A12}
\definecolor{popmuted}{HTML}{8A7F76}
\definecolor{popline}{HTML}{EADFD2}
\definecolor{popgray}{HTML}{8A7F76}
\colorlet{popgrey}{popgray}
% Alias tolérants (le modèle nomme parfois les couleurs différemment)
\colorlet{popwhite}{white}
\colorlet{popblack}{popink}
\colorlet{popred}{poppink}
\colorlet{popyellow}{popgold}
% Teintes claires (fonds de tableaux/encadrés) — le modèle nomme parfois
% les couleurs avec un suffixe L (ex. popblueL) sans les avoir définies.
\colorlet{poporangeL}{poporange!20}
\colorlet{popdarkL}{popdark!20}
\colorlet{poppurpleL}{poppurple!20}
\colorlet{popgreenL}{popgreen!20}
\colorlet{poppinkL}{poppink!20}
\colorlet{popblueL}{popblue!20}
\colorlet{popgoldL}{popgold!20}
\colorlet{popredL}{popred!20}
\colorlet{popgrayL}{popgray!20}
% Teintes claires pour fonds de tableaux / zones
\colorlet{popblueL}{popblue!10!white}
\colorlet{poporangeL}{poporange!14!white}
\colorlet{popgreenL}{popgreen!12!white}
\colorlet{poppurpleL}{poppurple!12!white}
\colorlet{poppinkL}{poppink!12!white}
\colorlet{popgoldL}{popgold!14!white}

\pagecolor{popcream}

% ============================================================
% Typographie
% ============================================================
\defaultfontfeatures{Ligatures=TeX}
\setmainfont{PlusJakartaSans-Medium.ttf}[Path=../../../fonts/,BoldFont=PlusJakartaSans-Bold.ttf]
% Maths en sans-serif (Fira Math) avec chiffres et lettres droites de Plus
% Jakarta, pour que les formules se fondent dans le texte.
\usepackage[math-style=ISO,bold-style=ISO]{unicode-math}
\setmathfont{FiraMath-Regular.otf}[Path=../../../fonts/]
\setmathfont{PlusJakartaSans-Medium.ttf}[Path=../../../fonts/,range={up/num,up/{Latin,latin},\mathup}]
\usepackage{icomma} % virgule décimale sans espace en mode math
% Caractères mathématiques Unicode absents des polices de texte (Plus Jakarta,
% Archivo Black) : rendus par leur équivalent mathématique, en texte comme en
% formule — sinon ils s'affichent en carré vide (ex. « Arithmétique dans ℤ »).
\usepackage{newunicodechar}
\newunicodechar{ℝ}{\ensuremath{\mathbb{R}}}
\newunicodechar{ℤ}{\ensuremath{\mathbb{Z}}}
\newunicodechar{ℂ}{\ensuremath{\mathbb{C}}}
\newunicodechar{ℕ}{\ensuremath{\mathbb{N}}}
\newunicodechar{ℚ}{\ensuremath{\mathbb{Q}}}
\newunicodechar{𝔻}{\ensuremath{\mathbb{D}}}
\newunicodechar{α}{\ensuremath{\alpha}}
\newunicodechar{β}{\ensuremath{\beta}}
\newunicodechar{γ}{\ensuremath{\gamma}}
\newunicodechar{δ}{\ensuremath{\delta}}
\newunicodechar{ε}{\ensuremath{\varepsilon}}
\newunicodechar{θ}{\ensuremath{\theta}}
\newunicodechar{λ}{\ensuremath{\lambda}}
\newunicodechar{μ}{\ensuremath{\mu}}
\newunicodechar{σ}{\ensuremath{\sigma}}
\newunicodechar{τ}{\ensuremath{\tau}}
\newunicodechar{φ}{\ensuremath{\varphi}}
\newunicodechar{ψ}{\ensuremath{\psi}}
\newunicodechar{ω}{\ensuremath{\omega}}
\newunicodechar{Σ}{\ensuremath{\Sigma}}
% majuscules produites par \MakeUppercase dans les titres (α-aminés -> Α)
\newunicodechar{Α}{\ensuremath{\alpha}}
\newunicodechar{Β}{\ensuremath{\beta}}
\newunicodechar{Γ}{\ensuremath{\Gamma}}
\newunicodechar{Δ}{\ensuremath{\Delta}}
\newunicodechar{Ε}{\ensuremath{\varepsilon}}
\newunicodechar{Θ}{\ensuremath{\Theta}}
\newunicodechar{Λ}{\ensuremath{\Lambda}}
\newunicodechar{Μ}{\ensuremath{\mu}}
\newunicodechar{Τ}{\ensuremath{\tau}}
\newunicodechar{Φ}{\ensuremath{\Phi}}
\newunicodechar{Ψ}{\ensuremath{\Psi}}
\newunicodechar{Ω}{\ensuremath{\Omega}}
\newunicodechar{↦}{\ensuremath{\mapsto}}
\newunicodechar{∈}{\ensuremath{\in}}
\newunicodechar{∉}{\ensuremath{\notin}}
\newunicodechar{∧}{\ensuremath{\wedge}}
\newunicodechar{⇔}{\ensuremath{\Leftrightarrow}}
\newunicodechar{⟺}{\ensuremath{\Longleftrightarrow}}
\newunicodechar{⇒}{\ensuremath{\Rightarrow}}
\newunicodechar{⟹}{\ensuremath{\Longrightarrow}}
\newunicodechar{∘}{\ensuremath{\circ}}
\newunicodechar{∪}{\ensuremath{\cup}}
\newunicodechar{∩}{\ensuremath{\cap}}
\newunicodechar{≡}{\ensuremath{\equiv}}
\newunicodechar{⊂}{\ensuremath{\subset}}
\newunicodechar{⊥}{\ensuremath{\perp}}
\newunicodechar{∃}{\ensuremath{\exists}}
\newunicodechar{∀}{\ensuremath{\forall}}
\newunicodechar{∛}{\ensuremath{\sqrt[3]{\,}}}
\newunicodechar{∜}{\ensuremath{\sqrt[4]{\,}}}
\newunicodechar{⃗}{\ensuremath{{}^{\to}}}
\newunicodechar{ⁿ}{\ifmmode^{n}\else\textsuperscript{\ensuremath{n}}\fi}
\newunicodechar{ˣ}{\ifmmode^{x}\else\textsuperscript{\ensuremath{x}}\fi}
\newunicodechar{ᵇ}{\ifmmode^{b}\else\textsuperscript{\ensuremath{b}}\fi}
\newunicodechar{ʳ}{\ifmmode^{r}\else\textsuperscript{\ensuremath{r}}\fi}
\newunicodechar{ᵘ}{\ifmmode^{u}\else\textsuperscript{\ensuremath{u}}\fi}
\newunicodechar{ʸ}{\ifmmode^{y}\else\textsuperscript{\ensuremath{y}}\fi}
\newunicodechar{ᵃ}{\ifmmode^{a}\else\textsuperscript{\ensuremath{a}}\fi}
\newunicodechar{ᵉ}{\ifmmode^{e}\else\textsuperscript{\ensuremath{e}}\fi}
\newunicodechar{ᵏ}{\ifmmode^{k}\else\textsuperscript{\ensuremath{k}}\fi}
\newunicodechar{ᵖ}{\ifmmode^{p}\else\textsuperscript{\ensuremath{p}}\fi}
\newunicodechar{ᴺ}{\ifmmode^{N}\else\textsuperscript{\ensuremath{N}}\fi}
\newunicodechar{ᶜ}{\ifmmode^{c}\else\textsuperscript{\ensuremath{c}}\fi}
\newunicodechar{ʰ}{\ifmmode^{h}\else\textsuperscript{\ensuremath{h}}\fi}
\newunicodechar{ᵠ}{\ifmmode^{q}\else\textsuperscript{\ensuremath{q}}\fi}
\newunicodechar{ₙ}{\ifmmode_{n}\else\textsubscript{\ensuremath{n}}\fi}
\newunicodechar{ₐ}{\ifmmode_{a}\else\textsubscript{\ensuremath{a}}\fi}
\newunicodechar{ᵢ}{\ifmmode_{i}\else\textsubscript{\ensuremath{i}}\fi}
\newunicodechar{ᵣ}{\ifmmode_{r}\else\textsubscript{\ensuremath{r}}\fi}
\newunicodechar{ₖ}{\ifmmode_{k}\else\textsubscript{\ensuremath{k}}\fi}
\newunicodechar{ᵦ}{\ifmmode_{\beta}\else\textsubscript{\ensuremath{\beta}}\fi}
\newunicodechar{ₓ}{\ifmmode_{x}\else\textsubscript{\ensuremath{x}}\fi}
\newfontfamily\popdisplay{ArchivoBlack-Regular.ttf}[Path=../../../fonts/]
\newfontfamily\poplogo{SpaceGrotesk-Bold.ttf}[Path=../../../fonts/]
\newfontfamily\popsemi{PlusJakartaSans-SemiBold.ttf}[Path=../../../fonts/]
\newfontfamily\popxbold{PlusJakartaSans-ExtraBold.ttf}[Path=../../../fonts/]
\newfontfamily\popxboldls{PlusJakartaSans-ExtraBold.ttf}[Path=../../../fonts/,LetterSpace=3.0]

\setlist[enumerate]{leftmargin=1.7em,itemsep=5pt,topsep=4pt}
\setlist[itemize]{leftmargin=1.7em,itemsep=4pt,topsep=4pt,label={\color{poporange}\textbullet}}
\setlength{\parindent}{0pt}
\setlength{\parskip}{5pt}
\renewcommand{\arraystretch}{1.25}

% pgfplots : style Pop par défaut
\pgfplotsset{
  every axis/.append style={axis line style={popink,line width=1pt},tick style={popink},
    grid style={popline,line width=0.5pt},label style={font=\small},tick label style={font=\footnotesize}},
  popcurve/.style={poporange,line width=1.8pt,no markers,smooth},
}
% Même style disponible dans un \draw TikZ simple (hors pgfplots)
\tikzset{popcurve/.style={poporange,line width=1.8pt}}
\tikzset{popgrid/.style={popline,very thin}}
\colorlet{popcurve}{poporange} % le nom du style est parfois utilisé comme couleur

% ============================================================
% Pill / squiggle
% ============================================================
\newcommand{\poppill}[3]{%
  \tikz[baseline=-0.55ex]\node[draw=popink,line width=1.2pt,rounded corners=2.6pt,%
    fill=#2,inner xsep=6.5pt,inner ysep=3pt]{%
    \popxboldls\fontsize{7}{7}\selectfont\color{#3}\MakeUppercase{#1}};%
}
\newcommand{\popsquiggle}[1]{%
  \tikz[baseline=-0.4ex]{
    \draw[#1,line width=2.6pt,line cap=round]
      plot [smooth] coordinates {(0,0.09) (0.34,0.01) (0.68,0.21) (1.02,0.01) (1.36,0.10)};
  }%
}
% Mot-clé surligné (marqueur orange sous le texte)
\newcommand{\cle}[1]{{\popxbold\color{popdark}#1}}

% ============================================================
% En-tête / pied
% ============================================================
\pagestyle{fancy}
\fancyhf{}
\renewcommand{\headrulewidth}{0pt}
\renewcommand{\footrulewidth}{0pt}
\lhead{\raisebox{-0.15ex}{\poplogo\fontsize{13}{13}\selectfont\color{popink}OnBuch\color{poporange}+}}
\rhead{\poppill{\DOCMATIERE\ \textbullet\ \DOCNIVEAU\ \textbullet\ Leçon \DOCLECON}{white}{popink}}
\lfoot{\popxbold\fontsize{7.6}{7.6}\selectfont\color{popmuted}\MakeUppercase{Cours signé OnBuch+}\ \textbullet\ {\popsemi\DOCTITRE}}
\rfoot{\tikz[baseline=-0.6ex]\node[rounded corners=8.5pt,fill=popink,minimum height=17pt,minimum width=17pt,inner xsep=5pt,inner sep=0]{\popxbold\fontsize{8}{8}\selectfont\color{popcream}\thepage\,/\,\pageref*{LastPage}};}

% ============================================================
% Titres
% ============================================================
\newcommand{\chaptitre}[2]{%
  \begin{tcolorbox}[enhanced,colback=poporange,colframe=popink,boxrule=2.6pt,arc=11pt,
      drop shadow=popink,left=15pt,right=15pt,top=12pt,bottom=11pt,before skip=3pt,after skip=18pt]
    \poppill{#2}{popink}{popcream}\par\vspace{9pt}
    {\raggedright\hyphenpenalty=10000\exhyphenpenalty=10000\popdisplay\fontsize{22}{24}\selectfont\color{popcream}\MakeUppercase{#1}\par}
  \end{tcolorbox}
}
% Section numérotée : pastille orange avec le numéro + titre Archivo
\newcounter{popsec}
\newcommand{\coursec}[1]{%
  \stepcounter{popsec}\setcounter{popsub}{0}%
  \par\vspace{16pt}\noindent
  \begin{minipage}{\linewidth}
  \tikz[baseline=-0.7ex]\node[draw=popink,line width=1.6pt,fill=poporange,rounded corners=4pt,minimum size=22pt,inner sep=0]{\popdisplay\fontsize{12}{12}\selectfont\color{popcream}\thepopsec};%
  \hspace{8pt}{\popdisplay\fontsize{15}{17}\selectfont\color{popink}\MakeUppercase{#1}}\par
  \vspace{4pt}\noindent{\color{poporange}\rule{36pt}{3.6pt}}
  \end{minipage}\par\nopagebreak\vspace{8pt}
}
\newcounter{popsub}
\newcommand{\courssub}[1]{%
  \stepcounter{popsub}%
  \par\vspace{9pt}\noindent{\popxbold\fontsize{11.5}{13}\selectfont\color{popink}{\color{poporange}\thepopsec.\thepopsub}\hspace{6pt}#1}\par\nopagebreak\vspace{4pt}
}

% ============================================================
% Boîtes pédagogiques — même recette : fond blanc, bordure épaisse colorée,
% ombre dure encre, badge plein. Titre optionnel [#1].
% ============================================================
\tcbset{popbase/.style={breakable,enhanced,colback=white,boxrule=2.3pt,arc=9pt,
  shadow={3pt}{-3pt}{0pt}{popink},left=14pt,right=14pt,top=11pt,bottom=11pt,before skip=11pt,after skip=15pt,
  fontupper=\popsemi\fontsize{10.6}{14.5}\selectfont\color{popink},
  fontlower=\popsemi\fontsize{10.6}{14.5}\selectfont\color{popink}}}
% Coche dessinée (le glyphe ✓ n'existe pas dans les polices du document)
\newcommand{\popcheck}[1]{\tikz[baseline=-0.5ex]\draw[#1,line width=1.6pt,line cap=round,line join=round](-0.13,0.0)--(-0.04,-0.09)--(0.13,0.1);}
\providecommand{\coche}{\popcheck{popgreen}}
\providecommand{\resultat}[1]{\par\smallskip\noindent\fcolorbox{popgreen}{popgreenL}{\parbox{\dimexpr\linewidth-2\fboxsep-2\fboxrule\relax}{\textbf{Résultat.} #1}}\par\smallskip}
\newcommand{\popboxhead}[3]{\poppill{#1}{#2}{white}\ifx\relax#3\relax\else\hspace{6pt}{\popxbold\fontsize{10.6}{12}\selectfont\color{popink}#3}\fi\par\vspace{7pt}}

\newtcolorbox{aretenir}[1][]{popbase,colframe=popgreen,before upper={\popboxhead{À retenir}{popgreen}{#1}}}
\newtcolorbox{exemplebox}[1][]{popbase,colframe=poporange,before upper={\popboxhead{Exemple}{poporange}{#1}}}
\newtcolorbox{definition}[1][]{popbase,colframe=popblue,before upper={\popboxhead{Définition}{popblue}{#1}}}
\newtcolorbox{propriete}[1][]{popbase,colframe=poppurple,before upper={\popboxhead{Propriété}{poppurple}{#1}}}
\newtcolorbox{methode}[1][]{popbase,colframe=popgold,before upper={\popboxhead{Méthode}{popgold}{#1}}}
\newtcolorbox{attention}[1][]{popbase,colframe=poppink,before upper={\popboxhead{Attention}{poppink}{#1}}}
\newtcolorbox{experience}[1][]{popbase,colframe=popblue,colback=popblueL,before upper={\popboxhead{Expérience}{popblue}{#1}}}
% « Le savais-tu ? » : carte violette pleine (contexte, culture, Cameroun)
\newtcolorbox{savaistu}[1][]{popbase,colback=poppurple,colframe=popink,
  fontupper=\popsemi\fontsize{10.4}{14.2}\selectfont\color{white},
  before upper={\poppill{Le savais-tu\,?}{white}{poppurple}\ifx\relax#1\relax\else\hspace{6pt}{\popxbold\color{white}#1}\fi\par\vspace{7pt}}}
% Exercice résolu : énoncé en haut, \tcblower puis la solution sur fond crème
\newcounter{popexo}\newcounter{popexr}
\newtcolorbox{exoresolu}[1][]{popbase,colframe=poporange,colbacklower=popcream,
  segmentation style={popink,line width=1.2pt,dashed},
  before upper={\stepcounter{popexr}\popboxhead{Exercice résolu \thepopexr}{poporange}{#1}},
  before lower={\poppill{Solution}{popink}{popcream}\par\vspace{6pt}}}
% Exercice (non résolu en place) — la correction est dans la partie Corrigés
\newtcolorbox{exercice}[1][]{popbase,colframe=popink,
  before upper={\stepcounter{popexo}\popboxhead{Exercice \thepopexo}{popink}{#1}}}
% Corrigé
\newtcolorbox{corrige}[1][]{popbase,colframe=popgreen,colback=popgreenL,
  before upper={\popboxhead{Corrigé}{popgreen}{#1}}}
% Objectifs / prérequis / bilan
\newtcolorbox{objectifs}{popbase,colframe=popink,colback=poporangeL,
  before upper={\popboxhead{Objectifs de la leçon}{popink}{}}}
\newtcolorbox{prerequis}{popbase,colframe=popink,colback=popblueL,
  before upper={\popboxhead{Prérequis}{popblue}{}}}
\newtcolorbox{bilan}{popbase,colframe=popink,colback=white,boxrule=2.8pt,
  before upper={\popboxhead{Fiche bilan}{poporange}{L'essentiel à connaître pour l'examen}}}
% Figure encadrée avec légende
\newtcolorbox{popfigure}[1][]{enhanced,breakable=false,colback=white,colframe=popline,boxrule=1.4pt,arc=8pt,
  left=10pt,right=10pt,top=10pt,bottom=8pt,before skip=10pt,after skip=12pt,halign=center,
  lower separated=false,halign lower=center,
  fontlower=\popsemi\fontsize{9}{11.5}\selectfont\color{popmuted},
  #1}
\newcommand{\legende}[1]{\par\vspace{6pt}{\popsemi\fontsize{9}{11.5}\selectfont\color{popmuted}#1}}

% ============================================================
% Signature OnBuch+
% ============================================================
\newcommand{\poplogoinline}[1]{{\poplogo\fontsize{#1}{#1}\selectfont\color{popink}OnBuch\color{poporange}+}}
\newcommand{\signatureonbuch}{%
  \begin{tcolorbox}[enhanced,colback=popink,colframe=popink,boxrule=2.2pt,arc=10pt,
      drop shadow=poporange,left=14pt,right=14pt,top=10pt,bottom=10pt,before skip=0pt,after skip=0pt]
    \tikz[baseline=-0.7ex]\node[circle,fill=popgreen,draw=popcream,line width=1.3pt,minimum size=22pt,inner sep=0]{\popcheck{white}};%
    \hspace{9pt}\begin{minipage}[c]{0.62\linewidth}
      {\popxbold\fontsize{12}{13}\selectfont\color{popcream}Signé\ \textbullet\ {\poplogo OnBuch\color{poporange}+}}\par\vspace{2pt}
      {\raggedright\popsemi\fontsize{8}{10.5}\selectfont\color{popline}\MakeUppercase{Équipe pédagogique OnBuch+}\\\MakeUppercase{Conforme au programme officiel MINESEC}\par}
    \end{minipage}\hfill
    {\poplogo\fontsize{22}{22}\selectfont\color{popcream}OnBuch\color{poporange}+}
  \end{tcolorbox}%
}

% ============================================================
% Page de garde
% #1 titre · #2 matière · #3 niveau · #4 module · #5 n° leçon · #6 durée ·
% #7 description / objectifs (texte ou itemize)
% ============================================================
\newcommand{\pagegarde}[7]{%
  \thispagestyle{empty}
  \newgeometry{a4paper,margin=1.7cm,top=1.6cm,bottom=1.4cm}%
  \begin{tikzpicture}[remember picture,overlay]
    \fill[poporange!18] ([xshift=-3.2cm,yshift=-3.6cm]current page.north east) circle (4.6cm);
    \fill[poppurple!14] ([xshift=2.4cm,yshift=7.6cm]current page.south west) circle (3.8cm);
  \end{tikzpicture}%
  {\poplogo\fontsize{30}{30}\selectfont\color{popink}OnBuch\color{poporange}+}\hfill
  \raisebox{0.6ex}{\poppill{Cours signé \textbullet\ Premium}{popink}{popcream}}\par
  \vspace{1pt}\popsquiggle{poppurple}\par
  \vspace{8pt}
  \begin{tcolorbox}[enhanced,colback=poporange,colframe=popink,boxrule=2.8pt,arc=13pt,
      drop shadow=popink,left=18pt,right=18pt,top=12pt,bottom=13pt,after skip=10pt]
    \poppill{#2 \textbullet\ #3}{popink}{popcream}\hspace{5pt}\poppill{#4}{white}{popink}\par\vspace{8pt}
    {\popdisplay\fontsize{14}{14}\selectfont\color{popink}LEÇON #5}\par\vspace{4pt}
    {\raggedright\hyphenpenalty=10000\exhyphenpenalty=10000\popdisplay\fontsize{25}{27}\selectfont\color{popcream}\MakeUppercase{#1}\par}
    \vspace{8pt}
    {\popxbold\fontsize{9.5}{9.5}\selectfont\color{popink}\MakeUppercase{Durée conseillée : #6}}
  \end{tcolorbox}
  \begin{tcolorbox}[enhanced,colback=white,colframe=popink,boxrule=2.2pt,arc=10pt,
      drop shadow=popink,left=16pt,right=16pt,top=10pt,bottom=10pt,after skip=8pt]
    \poppill{Dans cette leçon}{poppurple}{white}\par\vspace{5pt}
    \popsemi\fontsize{9.3}{12.6}\selectfont\color{popink}#7
  \end{tcolorbox}
  \noindent\begin{minipage}[t]{0.487\linewidth}
    \begin{tcolorbox}[enhanced,colback=popgreen,colframe=popink,boxrule=2.2pt,arc=10pt,
        drop shadow=popink,left=11pt,right=11pt,top=10pt,bottom=10pt,height=3.1cm,valign=center]
      \tcbox[colback=white,colframe=popink,boxrule=1.3pt,arc=5pt,boxsep=0pt,nobeforeafter,
        left=3pt,right=3pt,top=3pt,bottom=3pt,valign=center]{\qrcode[height=1.9cm]{https://play.google.com/store/apps/details?id=com.luvvix.onbuch&hl=fr}}%
      \hspace{8pt}\begin{minipage}[c]{0.5\linewidth}
        {\popdisplay\fontsize{11}{12.5}\selectfont\color{white}\MakeUppercase{Télécharge OnBuch}}\par\vspace{4pt}
        {\raggedright\popsemi\fontsize{8.4}{11}\selectfont\color{white}Gratuit \textbullet\ Google Play\\Tuteur IA, annales, concours, résultats\par}
      \end{minipage}
    \end{tcolorbox}
  \end{minipage}\hfill
  \begin{minipage}[t]{0.487\linewidth}
    \begin{tcolorbox}[enhanced,colback=poppurple,colframe=popink,boxrule=2.2pt,arc=10pt,
        drop shadow=popink,left=11pt,right=11pt,top=10pt,bottom=10pt,height=3.1cm,valign=center]
      \tikz[baseline=-0.7ex]\node[circle,fill=white,draw=popink,line width=1.3pt,minimum size=1.6cm,inner sep=0]{\popdisplay\fontsize{24}{24}\selectfont\color{poppurple}+};%
      \hspace{8pt}\begin{minipage}[c]{0.6\linewidth}
        {\popdisplay\fontsize{11}{12.5}\selectfont\color{white}\MakeUppercase{Passe à OnBuch+}}\par\vspace{4pt}
        {\raggedright\popsemi\fontsize{8.4}{11}\selectfont\color{white}Tous les cours signés, Tuteur IA illimité, fiches PDF — onglet OnBuch+ de l'app\par}
      \end{minipage}
    \end{tcolorbox}
  \end{minipage}
  \vspace{8pt}
  \popcontenu
  \vfill
  \signatureonbuch
  \restoregeometry
  \newpage
}

% Rangée « Ce cours contient » (page de garde)
\newcommand{\poptile}[3]{% #1 couleur #2 gros texte #3 légende
  \begin{tcolorbox}[enhanced,colback=#1,colframe=popink,boxrule=2pt,arc=9pt,shadow={2.5pt}{-2.5pt}{0pt}{popink},
      left=6pt,right=6pt,top=9pt,bottom=9pt,halign=center,valign=center,height=1.75cm,width=\linewidth,nobeforeafter]
    {\popdisplay\fontsize{12.5}{13}\selectfont\color{white}\MakeUppercase{#2}}\par\vspace{3pt}
    {\centering\popsemi\fontsize{7.8}{9.5}\selectfont\color{white}#3\par}
  \end{tcolorbox}}
\newcommand{\popcontenu}{%
  \par\noindent{\popxboldls\fontsize{7.5}{7.5}\selectfont\color{popmuted}CE COURS SIGNÉ CONTIENT}\par\vspace{7pt}
  \noindent\begin{minipage}[t]{0.235\linewidth}\poptile{popblue}{Cours}{complet, illustré, pas à pas}\end{minipage}\hfill
  \begin{minipage}[t]{0.235\linewidth}\poptile{poporange}{Méthodes}{et exercices résolus}\end{minipage}\hfill
  \begin{minipage}[t]{0.235\linewidth}\poptile{poppink}{Exercices}{type examen + corrigés}\end{minipage}\hfill
  \begin{minipage}[t]{0.235\linewidth}\poptile{popgreen}{Fiche bilan}{l'essentiel à retenir}\end{minipage}\par}

% Sommaire
\newtcolorbox{sommaire}{enhanced,colback=white,colframe=popink,boxrule=2.2pt,arc=10pt,
  drop shadow=popink,left=18pt,right=18pt,top=13pt,bottom=13pt,before skip=6pt,after skip=20pt,
  before upper={\poppill{Sommaire}{poppurple}{white}\par\vspace{9pt}\popsemi\fontsize{10.6}{14.5}\selectfont\color{popink}}}

% Signature de fin de cours — poussée tout en bas de la dernière page
\newcommand{\findecours}{%
  \par\vfill\nopagebreak
  \begin{center}\popsquiggle{poporange}\end{center}\vspace{4pt}
  \signatureonbuch
}
\AtBeginDocument{\def\llbracket{\mathopen{[\mkern-3.5mu[}}\def\rrbracket{\mathclose{]\mkern-3.5mu]}}} % intervalles d'entiers (glyphe ⟦ absent de la police)
% Glyphes absents de Fira Math : redéfinis à partir de glyphes disponibles
\AtBeginDocument{%
  \def\setminus{\mathbin{\backslash}}\let\smallsetminus\setminus
  \def\vdots{\mathord{\vbox{\baselineskip3.2pt\lineskiplimit0pt\kern4pt\hbox{.}\hbox{.}\hbox{.}}}}%
  \def\ddots{\mathinner{\mkern1mu\raise6.5pt\hbox{.}\mkern2mu\raise3.6pt\hbox{.}\mkern2mu\raise0.7pt\hbox{.}\mkern1mu}}%
  \def\triangle{\mathord{\scalebox{0.9}{$\symup{\Delta}$}}}%
  \def\nabla{\mathord{\raisebox{\depth}{\scalebox{1}[-1]{$\symup{\Delta}$}}}}%
  \def\checkmark{\popcheck{popdark}}%
  \def\star{\mathbin{\ast}}\def\bigstar{\mathord{\ast}}%
  \def\diamond{\mathbin{\scalebox{0.6}{$\Diamond$}}}%
  \def\bigcup{\mathop{\mathchoice{\raisebox{-0.3em}{\scalebox{1.7}{$\cup$}}}{\scalebox{1.25}{$\cup$}}{\cup}{\cup}}\displaylimits}%
  \def\bigcap{\mathop{\mathchoice{\raisebox{-0.3em}{\scalebox{1.7}{$\cap$}}}{\scalebox{1.25}{$\cap$}}{\cap}{\cap}}\displaylimits}%
  \def\lesseqqgtr{\mathrel{\lessgtr}}\def\bigoplus{\mathop{\oplus}\displaylimits}%
}
\providecommand{\ssi}{si et seulement si} % abréviation fréquente
\AtBeginDocument{\providecommand{\wideparen}{\overparen}} % arc de cercle

%% FIGFIX-BEGIN v5 — correctifs globaux des figures (docs/onbuch-generation/tools/figfix)
\usepackage{adjustbox}\usepackage{etoolbox}\tcbuselibrary{raster}
% toute figure TikZ/pgfplots plus large que la ligne est réduite à la largeur disponible
\BeforeBeginEnvironment{tikzpicture}{\begin{adjustbox}{max width=\linewidth}}
\AfterEndEnvironment{tikzpicture}{\end{adjustbox}}
% légendes pgfplots lisibles par-dessus les courbes
\pgfplotsset{every axis/.append style={legend style={fill=white,fill opacity=0.92,text opacity=1,draw=black!15,inner sep=3pt}}}
% étiquettes d'arêtes (syntaxe edge["..."]) avec fond blanc
\tikzset{every edge quotes/.append style={fill=white,inner sep=1.2pt}}
%% FIGFIX-END
\begin{document}

\pagegarde{Rédiger une dissertation}{Français}{Tle Tech}{Français – classe de Terminale technique}{N20}{11 séances (fiche de progression harmonisée)}{Cette leçon outille l'élève de Terminale technique pour rédiger une dissertation complète : introduction, développement avec citations et transitions, conclusion. Elle s'appuie sur la lecture méthodique du Vieux nègre et la médaille de Ferdinand Oyono et sur les notions de cohésion, de cohérence et de texte explicatif.
\vspace{4pt}
\begin{multicols}{2}\small
\begin{itemize}
\item À la fin de cette leçon, je sais analyser un sujet de dissertation et dégager une problématique pertinente
\item À la fin de cette leçon, je sais rédiger une introduction complète (accroche, présentation du sujet, problématique, annonce du plan)
\item À la fin de cette leçon, je sais assurer la cohésion et la cohérence de mon texte grâce à la progression thématique
\item À la fin de cette leçon, je sais rédiger un paragraphe argumentatif structuré (idée, développement, exemple, bilan)
\item À la fin de cette leçon, je sais insérer et commenter correctement une citation tirée du Vieux nègre et la médaille
\item À la fin de cette leçon, je sais rédiger des transitions efficaces entre les parties et les paragraphes
\end{itemize}\end{multicols}}

\begin{sommaire}
\begin{multicols}{2}
\begin{enumerate}
\item Analyser le sujet et construire la problématique
\item La progression thématique : cohésion et cohérence
\item Rédiger l'introduction
\item Rédiger le développement : paragraphes, citations et transitions
\item Rédiger la conclusion
\item Lecture méthodique du Vieux nègre et la médaille (n° 2 et 3) et remédiation
\item Activités d'intégration : produire une dissertation complète
\item Activité d'intégration
\item Méthodes et exercices résolus
\item Exercices
\item Corrigés des exercices
\item Fiche bilan
\end{enumerate}
\end{multicols}
\end{sommaire}

\begin{objectifs}
\begin{itemize}
\item À la fin de cette leçon, je sais analyser un sujet de dissertation et dégager une problématique pertinente
\item À la fin de cette leçon, je sais rédiger une introduction complète (accroche, présentation du sujet, problématique, annonce du plan)
\item À la fin de cette leçon, je sais assurer la cohésion et la cohérence de mon texte grâce à la progression thématique
\item À la fin de cette leçon, je sais rédiger un paragraphe argumentatif structuré (idée, développement, exemple, bilan)
\item À la fin de cette leçon, je sais insérer et commenter correctement une citation tirée du Vieux nègre et la médaille
\item À la fin de cette leçon, je sais rédiger des transitions efficaces entre les parties et les paragraphes
\item À la fin de cette leçon, je sais rédiger une conclusion (bilan, réponse à la problématique, ouverture)
\item À la fin de cette leçon, je sais relire, corriger et améliorer ma production (compte rendu et remédiation)
\item À la fin de cette leçon, je sais justifier ma démarche rédactionnelle devant un correcteur
\end{itemize}
\end{objectifs}

\begin{prerequis}
\begin{itemize}
\item Les types de textes (narratif, descriptif, argumentatif) vus en classe de Première
\item La lecture méthodique n° 1 du Vieux nègre et la médaille (situation d'énonciation, thèmes)
\item Les connecteurs logiques et les marqueurs de relation
\item La structure générale d'un texte argumentatif (thèse, arguments, exemples)
\item Les règles de citation et de référence à une œuvre littéraire
\end{itemize}
\end{prerequis}

\newpage

\coursec{Analyser le sujet et construire la problématique}

Au concours d'entrée à l'ENAM de Yaoundé, de nombreux candidats camerounais maîtrisent leurs connaissances mais échouent à l'épreuve de dissertation : introductions hors sujet, paragraphes sans transitions, citations jetées sans explication. Pourtant, rédiger une dissertation n'est pas un don : c'est une technique qui s'apprend, comme construire une maison à Douala exige un plan, des fondations et des finitions. Comment passer d'un sujet d'examen à une copie claire, cohérente et convaincante ? Cette leçon donne les outils méthodiques pour réussir l'épreuve reine du Baccalauréat technique.

Tout commence par une étape décisive, souvent négligée par les candidats pressés : \cle{l'analyse du sujet}. C'est elle qui détermine la pertinence de toute la copie. Une dissertation brillamment écrite mais à côté du sujet ne vaut rien ; une copie modeste mais parfaitement centrée sur la question posée est toujours récompensée.

\courssub{Qu'est-ce qu'une dissertation ?}

\begin{definition}[La dissertation]
La \cle{dissertation} est un exercice écrit de \textbf{réflexion personnelle et argumentée} sur un sujet donné. Elle consiste à discuter une question, un jugement ou une affirmation en mobilisant ses connaissances (littéraires, historiques, sociales) et en organisant sa pensée selon un plan rigoureux : une \textbf{introduction}, un \textbf{développement} en deux ou trois parties, et une \textbf{conclusion}.
\end{definition}

Contrairement au résumé, qui restitue la pensée d'un auteur, la dissertation demande au candidat de \emph{penser par lui-même}. Le correcteur n'attend pas une vérité unique : il évalue la capacité à construire un raisonnement, à justifier ses idées par des exemples précis et à s'exprimer dans un français correct. On peut défendre une position nuancée ; ce qui compte, c'est la \textbf{cohérence} du propos.

Au Baccalauréat technique, trois grands types de sujets sont proposés :

\begin{center}
\begin{tabularx}{\linewidth}{|l|X|X|}
\hline
\rowcolor{popblueL} \textbf{Type de sujet} & \textbf{Caractéristiques} & \textbf{Exemple} \\
\hline
Sujet de \cle{réflexion} & Question ouverte invitant à exposer ses idées sur un thème de société & « Le développement technologique améliore-t-il la vie des jeunes Camerounais ? » \\
\hline
Sujet de \cle{discussion} & Citation ou affirmation d'un auteur à discuter (approuver, nuancer, réfuter) & « Un écrivain déclare : "La tradition est un frein au progrès." Discutez cette opinion. » \\
\hline
Sujet lié à l'\cle{œuvre au programme} & Question portant sur l'œuvre étudiée en classe (ex. \emph{Le Vieux nègre et la médaille} de Ferdinand Oyono) & « La médaille remise à Meka récompense-t-elle réellement ses mérites ? » \\
\hline
\end{tabularx}
\end{center}

Quel que soit le type de sujet, la démarche d'analyse reste la même. Elle comporte six étapes que nous allons détailler.

\courssub{Les six étapes de l'analyse du sujet}

\begin{methode}[Analyser un sujet de dissertation en six étapes]
\begin{enumerate}
\item \textbf{Lire le sujet plusieurs fois} (au moins trois lectures) et \textbf{souligner les mots-clés} : les noms importants, les verbes, les adverbes de nuance (« toujours », « seulement », « dans quelle mesure »).
\item \textbf{Définir chaque terme important} du sujet, au brouillon, avec ses propres mots.
\item \textbf{Reformuler le sujet} dans une phrase personnelle, pour vérifier qu'on l'a compris.
\item \textbf{Poser des questions} sur le sujet : qui ? quoi ? pourquoi ? dans quelle mesure ? quelles limites ?
\item \textbf{Dégager la problématique} : formuler la question centrale à laquelle toute la copie répondra.
\item \textbf{Mobiliser ses connaissances} (brainstorming) puis \textbf{classer les idées} en deux ou trois axes, futur plan du développement.
\end{enumerate}
\end{methode}

\begin{popfigure}
\begin{center}
\begin{tikzpicture}[
  etape/.style={draw=popblue, fill=popblueL, rounded corners=2pt, text width=3.1cm, minimum height=1cm, align=center, font=\small},
  fleche/.style={-{Stealth[length=2.5mm]}, thick, poporange}
]
\node[etape, align=center] (e1) at (0,0){\textbf{1. Lecture}\\ mots-clés soulignés};
\node[etape, align=center] (e2) at (3.6,-1.3){\textbf{2. Définitions}\\ des termes importants};
\node[etape, align=center] (e3) at (7.2,-2.6){\textbf{3. Reformulation}\\ en ses propres mots};
\node[etape, align=center] (e4) at (10.8,-3.9){\textbf{4. Questions}\\ qui ? pourquoi ? limites ?};
\node[etape, fill=poporangeL, draw=poporange, align=center] (e5) at (7.2,-5.2){\textbf{5. Problématique}\\ question centrale};
\node[etape, fill=popgreenL, draw=popgreen, align=center] (e6) at (3.6,-6.5){\textbf{6. Brainstorming}\\ idées classées en axes};
\draw[fleche] (e1.south) -- (e2.west);
\draw[fleche] (e2.south) -- (e3.west);
\draw[fleche] (e3.south) -- (e4.west);
\draw[fleche] (e4.south) -- (e5.east);
\draw[fleche] (e5.south) -- (e6.east);
\end{tikzpicture}
\end{center}
\legende{Les six étapes de l'analyse du sujet : un parcours en escalier, de la lecture attentive (étape 1) jusqu'à la collecte et au classement des idées (étape 6), en passant par la formulation de la problématique (étape 5), cœur de la démarche.}
\end{popfigure}

\courssub{Étapes 1 et 2 : lire, souligner, définir}

La première lecture donne une impression générale ; la deuxième permet de souligner les \cle{mots-clés} ; la troisième vérifie qu'aucun terme n'a été négligé. Chaque mot du libellé compte : un adverbe comme « toujours » ou « seulement » engage une discussion sur l'absolu de l'affirmation.

Prenons le sujet : « Faut-il préserver les traditions camerounaises ? » Les mots-clés sont : \textbf{faut-il} (une question de nécessité, de devoir), \textbf{préserver} (protéger, maintenir vivant), \textbf{traditions} (coutumes, rites, langues, arts transmis de génération en génération), \textbf{camerounaises} (ancrage national précis : on ne parlera pas des traditions en général).

\begin{exemplebox}[Définir les termes d'un sujet]
Pour le sujet « Le développement technologique améliore-t-il la vie des jeunes Camerounais ? », on définira au brouillon :
\begin{itemize}
\item \textbf{développement technologique} : diffusion des outils numériques (téléphone portable, internet, mobile money) et techniques dans la société ;
\item \textbf{améliorer} : rendre meilleur, faciliter, augmenter le bien-être ;
\item \textbf{vie des jeunes Camerounais} : études, emploi, loisirs, relations sociales, santé des 15-35 ans au Cameroun.
\end{itemize}
Ces définitions provisoires pourront être reprises dans l'introduction pour fixer clairement le sens du débat.
\end{exemplebox}

\begin{attention}[Ne néglige aucun mot du libellé]
Un mot oublié, c'est une partie du sujet laissée de côté, donc un hors-sujet partiel. Dans « les jeunes \textbf{Camerounais} », l'adjectif impose de puiser ses exemples dans la réalité nationale (forages, réseaux sociaux à Yaoundé, mobile money, examens en ligne), et non dans des généralités abstraites.
\end{attention}

\courssub{Étapes 3 et 4 : reformuler et interroger le sujet}

La \cle{reformulation} est un test de compréhension : si l'on sait dire le sujet autrement, c'est qu'on l'a compris. « Faut-il préserver les traditions camerounaises ? » devient par exemple : « Avons-nous le devoir de transmettre et de maintenir vivantes les coutumes de nos communautés ? »

Vient ensuite le travail d'interrogation, qui prépare la problématique. Face au sujet, le candidat se comporte comme un enquêteur :

\begin{itemize}
\item \textbf{Qui ?} Qui est concerné par la préservation des traditions ? (les jeunes, les chefs traditionnels, l'État, les écoles)
\item \textbf{Quoi ?} Que faut-il préserver exactement ? (les langues, les rites, l'artisanat, les valeurs morales)
\item \textbf{Pourquoi ?} Pourquoi la question se pose-t-elle ? (parce que la modernité, l'urbanisation et les réseaux sociaux éloignent les jeunes des coutumes)
\item \textbf{Dans quelle mesure ?} Faut-il tout préserver ? N'y a-t-il pas des traditions dépassées ou contraires aux droits humains ?
\end{itemize}

C'est de la tension entre ces questions — préserver, mais jusqu'où ? — que naît la problématique.

\courssub{Étape 5 : dégager la problématique}

\begin{definition}[La problématique]
La \cle{problématique} est la \textbf{question centrale}, issue de l'analyse du sujet, qui \textbf{oriente toute la réflexion}. Elle ne recopie pas le sujet : elle le transforme en un véritable problème intellectuel, souvent en introduisant une tension (d'une part... d'autre part...) ou une nuance (dans quelle mesure... ?). Toute la copie — chaque partie, chaque paragraphe — devra y répondre.
\end{definition}

\begin{exemplebox}[Du sujet à la problématique]
\textbf{Sujet :} « Faut-il préserver les traditions camerounaises ? »

\textbf{Problématique possible :} « Dans quelle mesure la modernité menace-t-elle les traditions, et comment concilier héritage culturel et progrès ? »

On remarque que la problématique est plus riche que le sujet : elle nomme l'adversaire (la modernité), elle admet une gradation (« dans quelle mesure ») et elle ouvre sur une recherche d'équilibre (« comment concilier »), ce qui annonce déjà un plan en deux ou trois parties.
\end{exemplebox}

\begin{attention}[Erreur fréquente : recopier le sujet comme problématique]
Beaucoup de candidats écrivent dans leur introduction : « La problématique est la suivante : faut-il préserver les traditions camerounaises ? » C'est une simple \textbf{reprise du libellé}, non une problématique. Une vraie problématique \emph{transforme} le sujet : elle en dégage l'enjeu, la tension, la difficulté. Recopier le sujet prouve au correcteur que le candidat ne l'a pas analysé.
\end{attention}

\courssub{Étape 6 : mobiliser et classer ses idées (brainstorming)}

La problématique posée, il faut rassembler la matière de la réflexion. Le \cle{brainstorming} (ou « remue-méninges ») consiste à noter au brouillon, sans trier, toutes les idées, exemples, faits, proverbes et références qui viennent à l'esprit. Puis on \textbf{classe} ces idées en deux ou trois \textbf{axes}, qui deviendront les parties du développement.

\begin{popfigure}
\begin{center}
\begin{tikzpicture}[
  sujet/.style={draw=popdark, fill=popblueL, rounded corners=3pt, text width=4.6cm, align=center, font=\small\bfseries},
  branche/.style={draw=popgreen, fill=popgreenL, rounded corners=2pt, text width=3.6cm, align=center, font=\footnotesize},
  brancheB/.style={draw=poporange, fill=poporangeL, rounded corners=2pt, text width=3.6cm, align=center, font=\footnotesize},
  lien/.style={thick, popmuted}
]
\node[sujet] (s) at (0,0) {Le développement technologique améliore-t-il la vie des jeunes Camerounais ?};
\node[branche] (p1) at (-5.6,2.6) {Études : cours en ligne, recherches sur internet, préparation du Bac};
\node[branche] (p2) at (-5.6,0.4) {Économie : mobile money, petits commerces en ligne, emplois numériques};
\node[branche] (p3) at (-5.6,-1.8) {Communication : lien avec la famille éloignée, information rapide};
\node[brancheB] (c1) at (5.6,2.6) {Dépendance : réseaux sociaux, perte de temps, baisse des résultats scolaires};
\node[brancheB] (c2) at (5.6,0.4) {Exclusion : coût des données, fracture entre villes et villages};
\node[brancheB] (c3) at (5.6,-1.8) {Mœurs : cybercriminalité, arnaques, perte des valeurs};
\draw[lien] (s.west) -- (p1.east);
\draw[lien] (s.west) -- (p2.east);
\draw[lien] (s.west) -- (p3.east);
\draw[lien] (s.east) -- (c1.west);
\draw[lien] (s.east) -- (c2.west);
\draw[lien] (s.east) -- (c3.west);
\node[font=\footnotesize\itshape, popgreen] at (-5.6,3.6) {Axe 1 : les bienfaits (pour)};
\node[font=\footnotesize\itshape, poporange] at (5.6,3.6) {Axe 2 : les limites (contre)};
\end{tikzpicture}
\end{center}
\legende{Arbre de brainstorming pour un sujet de réflexion : à gauche, les idées « pour » (axe 1) ; à droite, les idées « contre » (axe 2). Un troisième axe possible serait : « les conditions d'un usage bénéfique de la technologie ».}
\end{popfigure}

Ce classement fait apparaître naturellement un \textbf{plan dialectique} : thèse (les bienfaits), antithèse (les limites), synthèse (les conditions d'un usage maîtrisé). Le plan n'est jamais inventé de toutes pièces : il \emph{naît} du tri des idées.

\begin{exoresolu}[Analyse complète d'un sujet de discussion]
\textbf{Énoncé.} Un écrivain affirme : « La colonisation n'a apporté que du malheur à l'Afrique. » Discutez cette opinion en vous appuyant sur vos connaissances et sur \emph{Le Vieux nègre et la médaille} de Ferdinand Oyono. Effectuez l'analyse du sujet (étapes 1 à 6).
\tcblower
\textbf{Solution détaillée.}

\textbf{Étape 1 — Mots-clés :} \emph{colonisation} ; \emph{n'a apporté que} (formule absolue, à discuter) ; \emph{malheur} ; \emph{Afrique}.

\textbf{Étape 2 — Définitions :} colonisation : domination politique, économique et culturelle d'un territoire africain par une puissance européenne ; malheur : souffrances, humiliations, spoliations.

\textbf{Étape 3 — Reformulation :} « La présence coloniale en Afrique a-t-elle été uniquement destructrice, ou a-t-elle aussi eu des effets que l'on peut discuter ? »

\textbf{Étape 4 — Questions :} Quels malheurs la colonisation a-t-elle causés (expropriations, travail forcé, humiliations subies par Meka dans le roman) ? Le mot « que » est-il acceptable ? Les écoles, routes et hôpitaux coloniaux profitaient-ils réellement aux Africains ? Comment les colonisés ont-ils vécu et résisté ?

\textbf{Étape 5 — Problématique :} « Dans quelle mesure peut-on affirmer que la colonisation n'a apporté que du malheur à l'Afrique, et que révèle l'expérience de Meka des illusions et des réalités du discours colonial ? »

\textbf{Étape 6 — Brainstorming et classement :} idées recueillies : médaille dérisoire, expropriation des terres de Meka, discours paternaliste du commandant, église et école coloniales, humiliations, révolte finale de Meka, prise de conscience. Classement : \textbf{axe 1} : les malheurs réels de la colonisation (spoliation, humiliation, illusions du roman) ; \textbf{axe 2} : les limites de la formule absolue « n'a apporté que » (certains apports matériels, mais au service du colonisateur) ; \textbf{axe 3} possible : la résistance et la prise de conscience africaine.
\end{exoresolu}

\courssub{Le piège majeur : le hors-sujet}

\begin{attention}[Le hors-sujet, conséquence d'une lecture rapide]
Le \cle{hors-sujet} est la faute la plus grave de la dissertation : la copie traite un autre problème que celui posé. Il résulte presque toujours d'une \textbf{lecture rapide} du libellé : le candidat reconnaît un mot (« tradition », « colonisation ») et récite tout ce qu'il sait sur ce thème, sans répondre à la question précise. Exemple : pour « Faut-il préserver les traditions camerounaises ? », réciter une leçon sur les danses traditionnelles sans jamais discuter le « faut-il » est un hors-sujet partiel. \textbf{Parade :} vérifier à chaque paragraphe que l'on répond bien à la problématique.
\end{attention}

\begin{savaistu}[Ce que le correcteur évalue d'abord]
Au Baccalauréat technique, le correcteur accorde une part importante de la note à la \textbf{pertinence de la problématique} et au \textbf{respect du sujet}. Dès la lecture de l'introduction, il repère si le candidat a compris la question posée. Une problématique claire et fidèle au sujet oriente favorablement la correction de toute la copie ; un hors-sujet, en revanche, plafonne la note quelle que soit la qualité de l'expression. C'est pourquoi les dix premières minutes de l'épreuve doivent être consacrées à l'analyse du sujet, au brouillon, sans écrire une seule ligne au propre.
\end{savaistu}

\begin{aretenir}[L'essentiel de la section]
\begin{itemize}
\item La dissertation est un exercice de \textbf{réflexion personnelle et argumentée} ; au Baccalauréat technique, on rencontre trois types de sujets : réflexion, discussion, sujet lié à l'œuvre.
\item L'analyse du sujet suit \textbf{six étapes} : lecture et mots-clés, définitions, reformulation, questions, problématique, brainstorming et classement en axes.
\item La \textbf{problématique} est la question centrale qui transforme le sujet en problème ; elle ne recopie jamais le libellé.
\item Le \textbf{hors-sujet} naît d'une lecture rapide ; chaque paragraphe doit répondre à la problématique.
\end{itemize}
\end{aretenir}

\begin{exercice}[À toi de jouer]
Analyse le sujet suivant en appliquant les six étapes (mots-clés, définitions, reformulation, questions, problématique, deux axes d'idées) : « L'école est-elle la seule voie de la réussite pour la jeunesse camerounaise ? »
\end{exercice}

\coursec{La progression thématique : cohésion et cohérence}

Dans la section précédente, nous avons appris à analyser un sujet de dissertation et à construire une problématique. Mais une bonne problématique ne suffit pas : encore faut-il que le texte que l'on rédige « tienne debout », que les phrases s'enchaînent naturellement et que le lecteur suive le raisonnement sans effort. C'est précisément l'objet de cette section : comprendre comment un texte est \emph{tissé}, comment ses phrases sont reliées entre elles par des fils invisibles mais bien réels. Ces fils portent un nom : la \cle{cohésion} et la \cle{cohérence}, organisées par ce que les linguistes appellent la \cle{progression thématique}.

\courssub{Cohérence et cohésion : deux notions à distinguer}

\textbf{1. L'intuition de départ}\par
Lisez ces deux mini-textes :

\begin{enumerate}
\item « Meka reçoit une médaille. Cette récompense le rend célèbre dans son village. Sa célébrité attire l'attention du commandant. »
\item « Meka reçoit une médaille. Le village compte de nombreux palmiers. Le commandant porte un uniforme. »
\end{enumerate}

Les deux textes sont composés de phrases grammaticalement correctes. Pourtant, le premier « coule » et le second semble « en miettes ». Pourquoi ? Parce que dans le premier texte, les phrases sont \emph{reliées} : chacune reprend un élément de la précédente et le prolonge. Dans le second, chaque phrase repart à zéro. Cette différence, c'est celle de la cohésion et de la cohérence.

\begin{definition}[Cohérence]
La \cle{cohérence} est l'unité de sens d'un texte : toutes les idées s'organisent autour d'un même sujet et sont liées logiquement entre elles (cause, conséquence, opposition, illustration...). Un texte cohérent répond à une question unique et ne contient aucune idée étrangère au sujet.
\end{definition}

\begin{definition}[Cohésion]
La \cle{cohésion} désigne les liens \emph{formels}, visibles à la surface du texte, qui relient les phrases entre elles : les reprises nominales, les pronoms, les démonstratifs, les synonymes et les connecteurs logiques. La cohésion est la \emph{matérialisation linguistique} de la cohérence.
\end{definition}

\begin{aretenir}[La relation entre les deux notions]
La cohérence concerne le \textbf{sens} (les idées forment un tout logique) ; la cohésion concerne la \textbf{forme} (les mots et les tournures qui relient les phrases). Un texte réussi est à la fois cohérent et cohésif : la cohésion sert la cohérence, comme le ciment sert la solidité d'un mur.
\end{aretenir}

\courssub{Thème et propos : les deux pôles de la phrase}

Pour comprendre comment les phrases s'enchaînent, il faut d'abord comprendre comment chaque phrase est construite intérieurement. Toute phrase informative peut être divisée en deux parties.

\begin{definition}[Thème et propos]
Le \cle{thème} (ou \emph{topique}) est \textbf{ce dont on parle} dans la phrase : c'est l'élément connu, déjà présent dans le contexte, le point de départ du message. Le \cle{propos} (ou \emph{commentaire}, \emph{rhème}) est \textbf{ce qu'on en dit} : c'est l'information nouvelle que la phrase apporte sur le thème.
\end{definition}

\begin{exemplebox}[Repérage du thème et du propos]
Dans la phrase « \textbf{Meka} / \emph{reçoit une médaille des mains du commandant} » :
\begin{itemize}
\item le thème est « Meka » (c'est de lui qu'on parle, on le connaît déjà) ;
\item le propos est « reçoit une médaille des mains du commandant » (c'est l'information nouvelle).
\end{itemize}
Dans la phrase suivante, « \textbf{Cette récompense} / \emph{le rend célèbre dans son village} », le thème est devenu « cette récompense », qui reprend le propos de la phrase précédente.
\end{exemplebox}

\begin{attention}[Thème $\neq$ sujet grammatical]
Le thème n'est pas toujours le sujet grammatical de la phrase. Dans « \emph{Une médaille, Meka la reçoit avec fierté} », le thème est « une médaille » alors que le sujet grammatical est « Meka ». Le thème se repère par la question : « de quoi parle-t-on ici ? ».
\end{attention}

\courssub{Les trois types de progression thématique}

La \cle{progression thématique} est la manière dont les thèmes et les propos s'enchaînent de phrase en phrase dans un texte. On distingue classiquement trois types de progression.

\textbf{1. La progression à thème constant}\par
Un même thème est repris comme point de départ de chaque phrase ; seuls les propos changent. Le texte « tourne » autour d'un même personnage, d'un même objet, d'une même idée, qu'il éclaire sous différents aspects.

\begin{exemplebox}[Thème constant]
« \textbf{Meka} est un vieillard respecté de Doum. \textbf{Il} a perdu ses deux fils à la guerre. \textbf{Le vieil homme} vit modestement avec sa femme Kelara. \textbf{Il} attend avec émotion la cérémonie de remise de sa médaille. »

Le thème « Meka » (repris par « il », « le vieil homme ») demeure constant ; chaque phrase ajoute une information nouvelle à son sujet. C'est la progression typique du portrait ou de la description.
\end{exemplebox}

\textbf{2. La progression linéaire (ou en chaîne)}\par
Le propos d'une phrase devient le thème de la phrase suivante, et ainsi de suite. Le texte avance « en entonnoir » ou « en chaîne » : chaque maillon naît du précédent. C'est la progression du raisonnement et du récit, car elle fait avancer la pensée.

\begin{exemplebox}[Progression linéaire]
« \textbf{Meka} reçoit une \emph{médaille}. \textbf{Cette récompense} le rend \emph{célèbre dans son village}. \textbf{Sa célébrité} attire les \emph{regards du commandant}. \textbf{Ces regards} révèlent bientôt \emph{l'hostilité de l'administration coloniale}. »

On observe la chaîne : médaille $\rightarrow$ récompense ; célèbre $\rightarrow$ célébrité ; regards $\rightarrow$ ces regards. Chaque propos devient le thème suivant.
\end{exemplebox}

\textbf{3. La progression à thème dérivé (ou à hyperthème)}\par
Les thèmes des différentes phrases ne sont pas identiques, mais ils \emph{dérivent} tous d'un même thème général, appelé \cle{hyperthème}. Chaque phrase traite un aspect, une partie ou une conséquence de ce thème global.

\begin{exemplebox}[Thème dérivé]
Hyperthème : \textbf{la colonisation}. \\
« \textbf{L'administration coloniale} impose ses règles aux villages. \textbf{L'école} forme des auxiliaires dociles. \textbf{La religion} est utilisée pour justifier la domination. \textbf{Les impôts} accablent les paysans. »

Les thèmes « administration », « école », « religion », « impôts » sont différents, mais tous dérivent de l'hyperthème « la colonisation ». C'est la progression typique du paragraphe de dissertation, où chaque argument illustre un aspect de la thèse.
\end{exemplebox}

\begin{popfigure}
\begin{tikzpicture}[font=\small, >=stealth,
  theme/.style={rectangle, rounded corners, draw=popblue, fill=popblueL, minimum height=0.7cm, inner sep=3pt},
  propos/.style={rectangle, rounded corners, draw=poporange, fill=poporangeL, minimum height=0.7cm, inner sep=3pt},
  titre/.style={font=\bfseries\small, text=popdark},
  lien/.style={->, thick, popmuted}]
% --- Panneau 1 : thème constant ---
\node[titre] at (0,4.6) {1. Thème constant};
\node[theme] (t1) at (-1.6,3.8) {Meka};
\node[propos] (p1) at (1.6,3.8) {est respecté};
\node[theme] (t2) at (-1.6,2.9) {Il};
\node[propos] (p2) at (1.6,2.9) {a perdu ses fils};
\node[theme] (t3) at (-1.6,2.0) {Le vieillard};
\node[propos] (p3) at (1.6,2.0) {attend sa médaille};
\draw[lien] (t1) -- (p1);
\draw[lien] (t2) -- (p2);
\draw[lien] (t3) -- (p3);
\draw[->, very thick, popblue, dashed] (t1.south) -- (t2.north);
\draw[->, very thick, popblue, dashed] (t2.south) -- (t3.north);
% --- Panneau 2 : progression linéaire ---
\node[titre] at (7,4.6) {2. Progression linéaire};
\node[theme] (a1) at (5.4,3.8) {Meka};
\node[propos] (b1) at (8.6,3.8) {reçoit une médaille};
\node[theme] (a2) at (5.4,2.9) {Cette récompense};
\node[propos] (b2) at (8.6,2.9) {le rend célèbre};
\node[theme] (a3) at (5.4,2.0) {Sa célébrité};
\node[propos] (b3) at (8.6,2.0) {attire le commandant};
\draw[lien] (a1) -- (b1);
\draw[lien] (a2) -- (b2);
\draw[lien] (a3) -- (b3);
\draw[->, very thick, poporange] (b1.south west) to[bend right=20] (a2.north east);
\draw[->, very thick, poporange] (b2.south west) to[bend right=20] (a3.north east);
% --- Panneau 3 : thème dérivé ---
\node[titre] at (3.5,-0.2) {3. Thème dérivé (hyperthème)};
\node[rectangle, rounded corners, draw=poppurple, fill=poppurpleL, minimum height=0.7cm, font=\bfseries] (hyper) at (3.5,-1.0) {LA COLONISATION};
\node[theme] (d1) at (0.5,-2.2) {l'administration};
\node[theme] (d2) at (3.5,-2.2) {l'école};
\node[theme] (d3) at (6.5,-2.2) {la religion};
\node[propos] (e1) at (0.5,-3.2) {impose ses règles};
\node[propos] (e2) at (3.5,-3.2) {forme des auxiliaires};
\node[propos] (e3) at (6.5,-3.2) {justifie la domination};
\draw[->, very thick, poppurple] (hyper.south) -- (d1.north);
\draw[->, very thick, poppurple] (hyper.south) -- (d2.north);
\draw[->, very thick, poppurple] (hyper.south) -- (d3.north);
\draw[lien] (d1) -- (e1);
\draw[lien] (d2) -- (e2);
\draw[lien] (d3) -- (e3);
\end{tikzpicture}
\legende{Les trois types de progression thématique. En bleu, les thèmes (ce dont on parle) ; en orange, les propos (l'information nouvelle). (1) Le même thème est repris de phrase en phrase. (2) Le propos d'une phrase devient le thème de la suivante. (3) Plusieurs thèmes dérivent d'un même hyperthème.}
\end{popfigure}

\begin{propriete}[Combinaison des progressions]
Un texte cohérent combine généralement \textbf{plusieurs types de progression thématique}. Dans une dissertation, le paragraphe entier relève souvent du thème dérivé (chaque argument dérive de la thèse), tandis qu'à l'intérieur du paragraphe, l'explication et l'exemple suivent une progression linéaire, et le commentaire de l'exemple une progression à thème constant.
\end{propriete}

\courssub{Les outils de la cohésion}

Comment le rédacteur matérialise-t-il ces reprises de thème ? Grâce à cinq grandes familles d'outils, à mobiliser consciemment dans vos paragraphes.

\begin{center}
\begin{tabularx}{\linewidth}{|l|X|X|}
\hline
\rowcolor{popblueL} \textbf{Outil} & \textbf{Fonction} & \textbf{Exemple} \\
\hline
Reprise nominale & Répéter le nom, parfois avec un déterminant démonstratif & « une médaille » $\rightarrow$ « \textbf{cette récompense} », « \textbf{cette distinction} » \\
\hline
Pronoms personnels & Reprendre un nom déjà mentionné sans le répéter & « Meka » $\rightarrow$ « \textbf{il} », « \textbf{le} vieillard \textbf{le} reçoit » \\
\hline
Démonstratifs et possessifs & Rappeler un élément du contexte & « \textbf{ce} commandant », « \textbf{sa} célébrité », « \textbf{son} village » \\
\hline
Synonymes et périphrases & Varier la reprise pour éviter les répétitions lourdes & « le vieux nègre » $\rightarrow$ « \textbf{le vieillard} », « \textbf{le héros du roman} » \\
\hline
Connecteurs logiques & Expliciter la relation entre les idées & « \textbf{en effet} », « \textbf{cependant} », « \textbf{par conséquent} », « \textbf{d'abord... ensuite} » \\
\hline
\end{tabularx}
\end{center}

\begin{attention}[Le texte « en miettes »]
L'erreur la plus fréquente des copies est l'accumulation de phrases justes mais sans liens : « Ferdinand Oyono dénonce la colonisation. Meka est naïf. Le commandant est hypocrite. La médaille est un symbole. » Chaque phrase est correcte, mais aucune ne s'accroche à la précédente : le texte tombe en miettes. Remède : chaque phrase doit contenir au moins un élément de reprise (pronom, démonstratif, synonyme) ou un connecteur logique qui la rattache à ce qui précède.
\end{attention}

\begin{methode}[Vérifier la cohésion de son paragraphe]
\begin{enumerate}
\item Relisez votre paragraphe phrase par phrase et \textbf{soulignez} toutes les reprises (pronoms, démonstratifs, synonymes, reprises nominales).
\item \textbf{Encadrez} les connecteurs logiques (d'abord, en effet, mais, ainsi...).
\item Pour chaque phrase, posez la question : « quel mot ou quelle idée la rattache à la phrase précédente ? » Si vous ne trouvez rien, la phrase est un maillon cassé : ajoutez une reprise ou un connecteur.
\item Identifiez le type de progression dominante (constante, linéaire, dérivée) et vérifiez qu'elle sert votre objectif : linéaire pour raisonner, dérivée pour énumérer des arguments, constante pour analyser un exemple.
\item Vérifiez enfin la cohérence globale : toutes les phrases parlent-elles bien de la même idée directrice, annoncée en tête de paragraphe ?
\end{enumerate}
\end{methode}

\courssub{Application : repérage dans \emph{Le Vieux nègre et la médaille}}

Appliquons ces notions à un passage inspiré du roman de Ferdinand Oyono, au moment où Meka apprend qu'il va recevoir une médaille.

\begin{exemplebox}[Extrait annoté]
« \textbf{Meka} venait d'apprendre qu'\textbf{il} allait recevoir une \emph{médaille}. \textbf{Cette distinction}, \emph{en effet}, récompensait ses terres données à la mission et ses fils morts pour la France. \textbf{Le vieillard}, \emph{d'abord} incrédule, finit par éprouver une grande fierté. \emph{Cependant}, \textbf{cette fierté} allait bientôt se changer en amertume. »

Relevé des outils de cohésion :
\begin{itemize}
\item reprises pronominales : « Meka » $\rightarrow$ « il » $\rightarrow$ « le vieillard » ;
\item reprise nominale par synonyme : « médaille » $\rightarrow$ « cette distinction » ;
\item chaîne linéaire : « fierté » (propos) $\rightarrow$ « cette fierté » (thème) ;
\item connecteurs : « en effet » (explication), « d'abord » (ordre), « cependant » (opposition, qui annonce le retournement).
\end{itemize}
La progression est d'abord \textbf{linéaire} (médaille $\rightarrow$ distinction $\rightarrow$ fierté $\rightarrow$ cette fierté), ce qui fait avancer le récit, avec un thème constant sous-jacent (Meka) qui unit le passage.
\end{exemplebox}

\begin{exoresolu}[Identifier et améliorer la progression thématique]
Énoncé. Voici un paragraphe d'élève : « La médaille est remise à Meka lors d'une grande cérémonie. Les Blancs boivent et dansent. Meka est oublié sous la pluie. La médaille ne protège de rien. » 1) Relevez les défauts de cohésion. 2) Réécrivez le paragraphe en assurant une progression linéaire et en ajoutant des connecteurs.
\tcblower
Solution détaillée. 1) Le paragraphe contient quatre phrases justes mais juxtaposées : aucune reprise ne relie « les Blancs » à la cérémonie, ni « Meka » (phrase 3) à la phrase 2 ; aucun connecteur n'explicite les relations logiques (opposition entre la fête et l'abandon, conséquence entre l'abandon et la leçon finale). Le texte est « en miettes ».

2) Réécriture possible : « \textbf{Meka} reçoit sa \emph{médaille} lors d'une grande cérémonie. \textbf{Cette fête}, \emph{cependant}, tourne vite à l'humiliation : \emph{tandis que} les Blancs boivent et dansent, \textbf{le vieillard} est oublié sous la pluie. \textbf{Cet abandon} révèle la vérité : \textbf{la médaille}, symbole de reconnaissance, \emph{en réalité} ne protège de rien. » On observe désormais une progression linéaire (cérémonie $\rightarrow$ cette fête ; oublié $\rightarrow$ cet abandon), des reprises (« Meka » $\rightarrow$ « le vieillard ») et des connecteurs (« cependant », « tandis que », « en réalité ») qui rendent le raisonnement lisible et cohérent.
\end{exoresolu}

\begin{savaistu}[Pourquoi cette notion est décisive au Baccalauréat technique]
Au correcteur, la cohésion saute aux yeux dès les premières lignes. Les barèmes de la dissertation réservent toujours une part importante à la « cohérence du texte » et à la « qualité de l'enchaînement des idées ». Un candidat qui maîtrise la progression thématique gagne des points sans connaissances supplémentaires : il transforme les mêmes idées en un texte qui se lit avec plaisir. C'est une compétence directement transférable aux rapports de stage, aux comptes rendus techniques et à toute la vie professionnelle.
\end{savaistu}

\begin{exercice}[À vous de jouer]
1) Dans le passage suivant, soulignez les thèmes, entourez les propos et indiquez le type de progression : « Le commandant de Doum prépare la cérémonie avec soin. Cette cérémonie doit impressionner le grand blanc de Yaoundé. Sa venue représente un honneur rare pour le poste. »

2) Classez les connecteurs suivants selon leur fonction (addition, opposition, cause, conséquence, illustration) : \emph{cependant, en effet, par conséquent, de plus, par exemple, car, néanmoins, ainsi}.

3) Rédigez un paragraphe de cinq phrases sur le thème « la naïveté de Meka », en imposant une progression à thème constant, puis transformez-le en progression linéaire. Comparez les deux effets obtenus.
\end{exercice}

\begin{corrige}[Exercice 1]
1) Thèmes : « le commandant », « cette cérémonie », « sa venue » ; propos : « prépare la cérémonie », « doit impressionner le grand blanc », « représente un honneur rare ». Progression \textbf{linéaire} : cérémonie $\rightarrow$ cette cérémonie ; impressionner le grand blanc $\rightarrow$ sa venue.

2) Addition : de plus ; opposition : cependant, néanmoins ; cause : car ; conséquence : par conséquent, ainsi ; illustration : en effet, par exemple.

3) Exemple à thème constant : « Meka est un homme simple. Il croit sincèrement à la gratitude des Blancs. Le vieillard imagine la médaille comme une consécration. Il ne voit pas le piège qu'elle recouvre. » Version linéaire : « Meka croit à la gratitude des Blancs. Cette croyance nourrit son attente de la médaille. Cette attente l'aveugle sur la réalité coloniale. Cet aveuglement causera sa chute. » La première version dresse un portrait (on accumule les traits) ; la seconde construit un raisonnement (on enchaîne les causes et les conséquences).
\end{corrige}

En résumé, la cohérence assure l'unité de sens, la cohésion la rend visible, et la progression thématique organise l'enchaînement des thèmes et des propos selon trois modèles complémentaires. Armés de ces outils, nous pouvons maintenant passer à la pratique : la rédaction de l'introduction, première étape concrète de la dissertation, fera l'objet de la section suivante.

\coursec{Rédiger l'introduction}

Dans la section précédente, nous avons vu comment la \cle{progression thématique} assure la cohésion et la cohérence d'un texte : chaque phrase reprend un élément connu pour y ajouter une information nouvelle, de sorte que le lecteur avance sans jamais se perdre. Cette règle vaut pour l'ensemble de la dissertation, mais elle trouve sa première application dans l'\cle{introduction}, car c'est là que le correcteur vérifie, dès les premières lignes, si votre pensée est organisée. Une introduction réussie est une introduction qui progresse : du général vers le particulier, du constat vers la question, de la question vers le plan. C'est cette mécanique que nous allons maintenant étudier et pratiquer.

\courssub{Fonction de l'introduction : mettre le lecteur en situation et annoncer la démarche}

Imaginez un visiteur qui entre dans une concession de véhicules à Douala. Le vendeur ne lui tend pas immédiatement un contrat à signer : il l'accueille, lui parle de la marque, lui montre le modèle, puis seulement lui propose un essai. L'introduction d'une dissertation joue exactement ce rôle d'accueil. Elle a une double fonction :

\begin{itemize}
\item \textbf{Mettre le lecteur en situation} : le placer dans le domaine du sujet, lui rappeler de quoi l'on parle, éveiller son intérêt. Le lecteur ne doit jamais être « jeté » brutalement sur la question.
\item \textbf{Annoncer la démarche} : formuler la question qui guidera la réflexion (la problématique) et indiquer les étapes par lesquelles on y répondra (l'annonce du plan).
\end{itemize}

\begin{aretenir}[La règle d'or de l'introduction]
L'introduction est un \cle{entonnoir} : elle part du large (le monde, la littérature, la vie) et se resserre progressivement jusqu'au point précis que sera votre devoir (le plan). Toute phrase qui rompt ce mouvement de resserrement — un détail donné trop tôt, une citation de l'œuvre placée avant l'amorce — est une faute de construction.
\end{aretenir}

\begin{popfigure}
\begin{center}
\begin{tikzpicture}[scale=1]
% Entonnoir : 4 trapèzes de largeur décroissante
\fill[popblueL] (-4.5,4.5) -- (4.5,4.5) -- (3.4,3.2) -- (-3.4,3.2) -- cycle;
\fill[popgreenL] (-3.4,3.2) -- (3.4,3.2) -- (2.3,1.9) -- (-2.3,1.9) -- cycle;
\fill[poporangeL] (-2.3,1.9) -- (2.3,1.9) -- (1.3,0.6) -- (-1.3,0.6) -- cycle;
\fill[poppinkL] (-1.3,0.6) -- (1.3,0.6) -- (0.5,-0.7) -- (-0.5,-0.7) -- cycle;
% Contours pour lisibilité
\draw[popline, thick] (-4.5,4.5) -- (4.5,4.5) -- (3.4,3.2) -- (-3.4,3.2) -- cycle;
\draw[popline, thick] (-3.4,3.2) -- (3.4,3.2) -- (2.3,1.9) -- (-2.3,1.9) -- cycle;
\draw[popline, thick] (-2.3,1.9) -- (2.3,1.9) -- (1.3,0.6) -- (-1.3,0.6) -- cycle;
\draw[popline, thick] (-1.3,0.6) -- (1.3,0.6) -- (0.5,-0.7) -- (-0.5,-0.7) -- cycle;
% Étiquettes
\node[font=\small\bfseries, popdark] at (0,3.85) {1. L'amorce (le plus large)};
\node[font=\small\bfseries, popdark] at (0,2.55) {2. Présentation du sujet};
\node[font=\small\bfseries, popdark] at (0,1.25) {3. Problématique};
\node[font=\small\bfseries, popdark] at (0,0.0) {4. Plan};
% Flèche de resserrement
\draw[-{Stealth[length=3mm]}, very thick, poppurple] (5.2,4.5) -- (5.2,-0.7);
\node[rotate=90, font=\small\itshape, poppurple] at (5.7,1.9) {resserrement progressif};
\end{tikzpicture}
\end{center}
\legende{Schéma en entonnoir de l'introduction : chaque composante est plus étroite que la précédente. L'amorce ouvre sur le monde (1), la présentation cerne le sujet (2), la problématique pose la question unique (3), l'annonce du plan désigne le chemin précis de la réponse (4).}
\end{popfigure}

\courssub{Composante 1 : l'amorce (ou accroche)}

\begin{definition}[L'amorce]
L'\cle{amorce} (ou \cle{accroche}) est la \textbf{première phrase de l'introduction}, destinée à \textbf{capter l'attention du lecteur}. Elle situe le sujet dans un cadre général : un fait de société, un constat sur la nature humaine, une citation d'auteur, un événement historique ou culturel lié au thème du sujet.
\end{definition}

L'intuition est simple : on n'entre jamais dans une conversation en criant la conclusion. On commence par un terrain commun. L'amorce crée ce terrain commun entre vous et le correcteur. Quatre types d'amorce sont recevables au Baccalauréat technique :

\begin{center}
\begin{tabularx}{\linewidth}{|l|X|X|}
\hline
\rowcolor{popblueL} \textbf{Type d'amorce} & \textbf{Principe} & \textbf{Exemple (sujet : la médaille de Meka)} \\
\hline
Fait général & Un constat large sur la société ou l'homme & « Dans toutes les sociétés, l'homme recherche la reconnaissance de ses pairs. » \\
\hline
Constat historique & Un repère lié au contexte de l'œuvre & « La colonisation a profondément bouleversé les rapports entre les peuples d'Afrique et d'Europe. » \\
\hline
Citation & Une parole d'auteur (exacte et attribuée) & « "Le colonialisme n'est pas une machine à penser", écrit Aimé Césaire. » \\
\hline
Événement lié au sujet & Un fait précis, daté, en rapport avec le thème & « En 1956, Ferdinand Oyono publie un roman qui dénonce les illusions de l'assimilation coloniale. » \\
\hline
\end{tabularx}
\end{center}

\begin{attention}[L'amorce creuse : « Depuis la nuit des temps... »]
L'erreur la plus fréquente — et la plus sanctionnée — consiste à ouvrir par une formule vide : « Depuis la nuit des temps, l'homme... », « De tout temps... », « L'homme a toujours... ». Ces amorces ne disent rien, s'appliquent à tous les sujets et signalent au correcteur un candidat sans culture mobilisable. Préférez toujours un constat \textbf{précis} et \textbf{en rapport direct} avec le sujet. Autre piège : l'amorce qui cite l'œuvre ou l'auteur dès la première phrase — c'est trop étroit, cela viole la logique de l'entonnoir.
\end{attention}

\courssub{Composante 2 : la présentation du sujet}

Après l'amorce, il faut \textbf{resserrer} : ramener le lecteur du constat général vers le sujet précis posé par l'examen. Cette présentation accomplit trois tâches :

\begin{enumerate}
\item \textbf{Définir les termes clés du sujet.} Si le sujet est « La médaille de Meka est-elle une récompense ou une humiliation ? », il faut dire brièvement ce qu'on entend par \emph{récompense} (une distinction qui honore) et par \emph{humiliation} (un acte qui rabaisse, qui nie la dignité).
\item \textbf{Situer le contexte.} Pour un sujet portant sur \emph{Le Vieux Nègre et la Médaille}, on rappelle en une ou deux phrases le cadre : le roman de Ferdinand Oyono, le village de Doum, le vieux Meka qui reçoit une médaille de l'administration coloniale pour avoir donné ses terres à la mission et ses deux fils à la guerre.
\item \textbf{Dégager l'enjeu.} Pourquoi cette question mérite-t-elle d'être posée ? Parce que derrière la médaille se joue le rapport de domination entre colonisateur et colonisé : l'enjeu est donc à la fois littéraire, historique et humain.
\end{enumerate}

\begin{exemplebox}[Présentation du sujet bien menée]
« Dans le roman de Ferdinand Oyono, \emph{Le Vieux Nègre et la Médaille} (1956), le vieux Meka, paysan de Doum, reçoit des mains du commandant une médaille censée couronner ses sacrifices : ses terres offertes à la mission catholique et ses deux fils morts à la guerre pour la France. Mais cette distinction, accordée par l'administration coloniale, soulève une question de fond : honore-t-elle réellement celui qui la reçoit ? »
\end{exemplebox}

\courssub{Composante 3 : la problématique}

\begin{definition}[La problématique]
La \cle{problématique} est la \textbf{question directrice} de la dissertation : une question unique, claire, qui naît naturellement de la présentation du sujet et à laquelle tout le développement devra répondre. Elle est généralement formulée par une interrogation directe ou indirecte.
\end{definition}

La problématique n'est pas une question inventée : elle est \textbf{déjà contenue dans le sujet}. Le travail consiste à la dégager, souvent en reprenant les termes du sujet sous forme interrogative. Trois formulations sont possibles :

\begin{itemize}
\item \textbf{Interrogation directe} : « La médaille de Meka est-elle une véritable récompense ou une humiliation déguisée ? »
\item \textbf{Interrogation indirecte} : « On peut alors se demander si la médaille de Meka constitue une récompense ou une humiliation. »
\item \textbf{Formulation problématisée} : « Comment une distinction honorifique peut-elle devenir l'instrument d'une domination ? »
\end{itemize}

\begin{attention}[Pièges de la problématique]
Évitez : (1) la \textbf{multiplication des questions} — une seule question directrice, les autres sont des sous-questions ; (2) la question \textbf{hors sujet} — elle doit reprendre les termes du sujet, pas un thème voisin ; (3) la question \textbf{trop factuelle} (« Qui est Meka ? ») qui n'appelle aucune réflexion. La problématique doit ouvrir un \emph{débat}, pas demander une définition.
\end{attention}

\courssub{Composante 4 : l'annonce du plan}

Dernière étape de l'entonnoir : indiquer au lecteur le chemin que suivra la réponse. Au Baccalauréat technique, le plan est le plus souvent \textbf{dialectique en deux parties} (thèse / antithèse, éventuellement dépassées) ou \textbf{thématique en deux ou trois parties}.

\begin{methode}[Rédiger une introduction en 4 temps]
\begin{enumerate}
\item \textbf{Amorce} : une phrase générale (fait, constat, citation) en rapport avec le thème du sujet — sans citer l'œuvre.
\item \textbf{Présentation du sujet} : deux à quatre phrases qui définissent les termes, situent l'œuvre et le contexte, dégagent l'enjeu.
\item \textbf{Problématique} : une question unique, claire, formulée à partir des termes du sujet.
\item \textbf{Annonce du plan} : une ou deux phrases qui présentent les grandes parties : « Nous verrons d'abord que..., puis nous montrerons que... »
\end{enumerate}
Vérifiez ensuite la longueur : \textbf{8 à 12 lignes} au Baccalauréat technique. En dessous, une composante manque ; au-delà, vous empiétez sur le développement.
\end{methode}

\begin{savaistu}[Pourquoi le correcteur aime les annonces de plan]
Un correcteur de Baccalauréat lit des centaines de copies en quelques jours. Une annonce de plan claire — « Nous verrons d'abord... puis... » — lui permet de connaître l'organisation de votre devoir avant même de le lire, et le rassure : il sait où vous allez et peut suivre votre progression. À l'inverse, une copie sans annonce de plan oblige le correcteur à deviner votre démarche, ce qui fatigue et irrite. L'annonce du plan est donc un contrat de lecture que vous signez avec lui : tenez-le !
\end{savaistu}

\courssub{Exemple complet et amélioration par la relecture ciblée}

\begin{exemplebox}[Introduction rédigée — Sujet : « La médaille de Meka est-elle une récompense ou une humiliation ? »]
« Dans toutes les sociétés, les distinctions honorifiques sont censées traduire la gratitude d'un groupe envers celui qui s'est dévoué pour lui. \textbf{(amorce)} La littérature africaine francophone, née dans le contexte colonial, s'est souvent interrogée sur la sincérité de ces marques de reconnaissance. Dans \emph{Le Vieux Nègre et la Médaille} (1956), Ferdinand Oyono raconte l'histoire de Meka, vieux paysan de Doum, qui reçoit une médaille de l'administration coloniale après avoir donné ses terres à la mission et perdu ses deux fils à la guerre. Une récompense honore celui qui la reçoit ; une humiliation, au contraire, rabaisse et nie sa dignité. Or la médaille de Meka, loin de l'élever, semble précipiter sa chute. \textbf{(présentation)} Cette médaille est-elle alors une véritable récompense ou une humiliation déguisée ? \textbf{(problématique)} Nous verrons d'abord que la médaille apparaît comme une récompense aux yeux de Meka et de son village, puis nous montrerons qu'elle constitue en réalité une humiliation qui révèle le mépris de l'administration coloniale. \textbf{(annonce du plan)} »
\end{exemplebox}

\begin{methode}[La relecture ciblée de l'introduction]
Après avoir rédigé, relisez votre introduction en ne cherchant qu'\textbf{une chose à la fois}, dans l'ordre :
\begin{enumerate}
\item L'amorce est-elle présente, générale, sans cliché ?
\item Les termes du sujet sont-ils définis, le contexte situé ?
\item La problématique est-elle une question unique reprenant les termes du sujet ?
\item Le plan est-il annoncé avec des connecteurs d'ordre (d'abord, ensuite, enfin) ?
\item La longueur est-elle comprise entre 8 et 12 lignes ?
\end{enumerate}
Cochez mentalement chaque point ; corrigez uniquement ce qui manque. Cette technique de relecture ciblée est plus efficace qu'une relecture globale, où l'œil glisse sans rien voir.
\end{methode}

\begin{exoresolu}[Introduction ratée vs introduction améliorée]
Sujet : « La médaille de Meka est-elle une récompense ou une humiliation ? » Voici l'introduction d'un candidat : « Depuis la nuit des temps, l'homme aime les médailles. Meka est un vieux nègre qui reçoit une médaille. Il est très content. Mais après il est arrêté. Est-ce que la médaille est bonne ou mauvaise ? Je vais parler de la médaille. » Identifiez les fautes, puis réécrivez cette introduction.
\tcblower
\textbf{Diagnostic des fautes :} (1) amorce creuse et banale (« Depuis la nuit des temps ») ; (2) présentation réduite à un résumé enfantin, sans définition des termes ni contexte (ni l'auteur, ni l'œuvre, ni la colonisation ne sont nommés) ; (3) problématique maladroite et familière (« bonne ou mauvaise ») ; (4) annonce de plan inexistante (« Je vais parler de... » n'annonce aucune partie) ; (5) registre oral (« très content »).
\textbf{Version améliorée :} « Les distinctions honorifiques traduisent, en principe, la reconnaissance d'une communauté envers celui qui s'est sacrifié pour elle. Mais dans l'Afrique coloniale, les gestes du colonisateur envers le colonisé étaient rarement désintéressés. C'est ce que montre Ferdinand Oyono dans \emph{Le Vieux Nègre et la Médaille} (1956) : le vieux Meka, qui a offert ses terres à la mission et ses deux fils à la guerre, reçoit solennellement une médaille de l'administration. Si une récompense honore, une humiliation rabaisse : or cette médaille semble faire basculer Meka de la fierté vers la honte. La médaille de Meka est-elle donc une récompense ou une humiliation ? Nous verrons d'abord qu'elle revêt l'apparence d'une récompense, puis qu'elle constitue en vérité une humiliation. »
\textbf{Commentaire :} chaque composante est désormais présente, dans l'ordre de l'entonnoir ; le registre est soutenu ; la longueur est d'environ 10 lignes.
\end{exoresolu}

\courssub{Production guidée : rédiger votre introduction}

\begin{exercice}[À vous de rédiger]
Sujet : « Dans \emph{Le Vieux Nègre et la Médaille}, le village de Doum admire Meka après la remise de la médaille. Cette admiration est-elle justifiée ? » Rédigez une introduction complète de 8 à 12 lignes en suivant la méthode en 4 temps. Consignes : (1) choisissez un type d'amorce dans le tableau du cours ; (2) définissez le mot \emph{admiration} ; (3) situez l'œuvre et l'auteur ; (4) formulez une problématique unique ; (5) annoncez un plan en deux parties. Relisez ensuite avec la grille de relecture ciblée et faites corriger par votre voisin.
\end{exercice}

\begin{corrige}[Exercice : production guidée]
Éléments de correction (une formulation possible parmi d'autres) : \textbf{Amorce} (constat) : « L'approbation de la communauté est l'une des formes les plus anciennes de la consécration sociale. » \textbf{Présentation} : admirer, c'est porter sur quelqu'un un regard de respect mêlé d'éloge ; dans \emph{Le Vieux Nègre et la Médaille} (1956) de Ferdinand Oyono, le village de Doum célèbre Meka, décoré par l'administration coloniale, et l'envie pour cette distinction. \textbf{Problématique} : « Cette admiration du village est-elle justifiée, ou repose-t-elle sur une illusion ? » \textbf{Annonce du plan} : « Nous verrons d'abord les raisons qui rendent cette admiration compréhensible, puis nous montrerons qu'elle repose sur un malentendu que la chute de Meka révèle cruellement. » Critères d'évaluation : présence des 4 composantes (4 pts), ordre de l'entonnoir respecté (2 pts), problématique unique et fidèle au sujet (2 pts), correction de la langue (2 pts).
\end{corrige}

\begin{aretenir}[L'essentiel de la section]
Une introduction de dissertation comporte \textbf{quatre composantes ordonnées en entonnoir} : l'\cle{amorce} (générale, sans cliché), la \cle{présentation du sujet} (définitions, contexte, enjeu), la \cle{problématique} (une question unique) et l'\cle{annonce du plan} (deux ou trois parties). Longueur : 8 à 12 lignes. Améliorez toujours votre introduction par une relecture ciblée, composante par composante. Une introduction réussie est un contrat signé avec le correcteur : le développement, que nous étudierons dans la prochaine section, devra en honorer chaque clause.
\end{aretenir}

\coursec{Rédiger le développement : paragraphes, citations et transitions}

Vous avez rédigé une introduction solide : accroche, présentation de l'œuvre ou du thème, problématique claire, annonce du plan. Mais une belle porte d'entrée ne suffit pas : c'est dans le développement que se joue l'essentiel de la note. Le développement est le \emph{corps} de la dissertation, le lieu où vous démontrez réellement votre thèse. Cette section vous apprend à bâtir ce corps, étage par étage : les paragraphes, les citations qui les nourrissent, les transitions qui les relient, et la progression thématique qui assure la cohésion de l'ensemble.

\courssub{La structure générale du développement}

Une dissertation au Baccalauréat technique se développe en \cle{deux ou trois parties}, chacune correspondant à une étape de votre plan annoncé en introduction. Chaque partie est elle-même composée de \cle{deux ou trois paragraphes}. On obtient donc un développement de quatre à neuf paragraphes au total.

\begin{exemplebox}[Le poids du développement dans la copie]
Au Baccalauréat technique, le développement représente environ \textbf{les deux tiers de la copie}, soit \textbf{25 à 35 lignes} sur une copie d'une quarantaine de lignes. L'introduction occupe 6 à 8 lignes, la conclusion 4 à 6 lignes. C'est donc dans le développement que le correcteur juge votre capacité à argumenter : négliger cette partie, c'est renoncer à la majorité des points.
\end{exemplebox}

\begin{attention}[Piège : le développement sans plan visible]
Chaque partie doit être \emph{visible} : on saute une ligne entre deux parties, et on va à la ligne en alinéa entre deux paragraphes. Un développement rédigé en un seul bloc donne l'impression d'une pensée confuse, même si les idées sont bonnes. Le correcteur doit \emph{voir} votre plan sans avoir à le deviner.
\end{attention}

\courssub{La progression thématique : cohésion et cohérence}

\emph{(Point de programme : Morphosyntaxe — Progression thématique)}

Au-delà de la structure visible (parties, paragraphes, alinéas), la fluidité du développement repose sur la \cle{progression thématique} : la manière dont l'information neuve (le rhème) d'une phrase devient l'information connue (le thème) de la phrase suivante. Maîtriser cette progression, c'est garantir la \textbf{cohérence} (liens logiques entre les idées) et la \textbf{cohésion} (liens linguistiques explicites) de votre texte.

\begin{definition}[Les trois types de progression thématique]
\begin{itemize}
\item \textbf{Progression linéaire (en chaînage) :} Le rhème de la phrase $n$ devient le thème de la phrase $n+1$. \emph{Exemple : « Meka reçoit une médaille. Cette médaille est remise par le commandant. Le commandant prononce un discours. »} Idéal pour narrer ou expliquer un mécanisme pas à pas.
\item \textbf{Progression à thème constant (parallèle) :} Le même thème est repris au début de chaque phrase. \emph{Exemple : « Meka est fier. Meka tremble d'émotion. Meka espère un avenir meilleur. »} Idéal pour décrire un personnage ou accumuler des arguments sur une même thèse.
\item \textbf{Progression à thèmes dérivés (rayonnante) :} À partir d'un thème principal (hyperthème), on développe des sous-thèmes. \emph{Exemple : « La cérémonie est solennelle. Le décor impressionne la foule. Les discours flattent Meka. Les honneurs masquent la réalité. »} Idéal pour structurer un paragraphe entier autour de l'idée directrice.
\end{itemize}
\end{definition}

\begin{methode}[Assurer la progression thématique dans un paragraphe]
\begin{enumerate}
\item Identifiez le \textbf{thème principal} du paragraphe (souvent le sujet de l'idée directrice).
\item Choisissez le type de progression adapté : \emph{linéaire} pour l'explication pas à pas, \emph{constant} pour l'accumulation d'arguments, \emph{dérivé} pour l'analyse d'un concept.
\item Utilisez des \textbf{reprises} (pronoms, synonymes, anaphores) et des \textbf{connecteurs logiques} pour rendre les liens explicites.
\end{enumerate}
\end{methode}

\begin{exemplebox}[Progression thématique dans le paragraphe modèle]
\emph{« Dans un premier temps, la médaille apparaît comme un honneur exceptionnel pour le vieux Meka [Thème constant : la médaille/Meka], \textbf{parce qu'}elle couronne publiquement toute une vie de soumission [Rhème $\rightarrow$ Thème suivant : cette soumission]. \textbf{En effet}, Meka a toujours été un sujet docile [Thème constant : Meka] : il a donné ses terres, ses fils à la guerre, sa foi au christianisme [Progression linéaire : terres $\rightarrow$ fils $\rightarrow$ foi]. »}
\end{exemplebox}

\courssub{La structure d'un paragraphe argumentatif}

Imaginez un paragraphe comme une petite dissertation en miniature : il annonce une idée, il la développe, il la prouve, il conclut. Un paragraphe argumentatif bien construit comporte \textbf{quatre étages} :

\begin{enumerate}
\item \textbf{L'idée directrice} (ou argument) : une phrase qui affirme clairement ce que le paragraphe va démontrer. C'est la première phrase du paragraphe.
\item \textbf{L'explication} : deux ou trois phrases qui précisent l'idée, la rendent intelligible, en dégagent le sens et la portée. C'est ici que s'applique la progression thématique et le texte explicatif.
\item \textbf{L'exemple ou la citation} : une preuve concrète tirée de l'œuvre étudiée ou de la culture générale, qui vient étayer l'argument.
\item \textbf{La mini-conclusion} (ou bilan) : une phrase qui tire la leçon de l'exemple et ramène à la problématique.
\end{enumerate}

\begin{popfigure}
\begin{center}
\begin{tikzpicture}[
  bloc/.style={rectangle, rounded corners=3pt, draw=popdark, thick, align=center, text width=0.62\linewidth, minimum height=1.05cm, inner sep=5pt},
  fleche/.style={-{Stealth[length=3mm]}, very thick, popmuted},
  num/.style={circle, draw=popdark, fill=popgoldL, thick, minimum size=0.7cm, inner sep=1pt, font=\bfseries}
]
\node[bloc, fill=popblueL, align=center] (idee) at (0,0){\textbf{1. Idée directrice}\\ \small L'argument du paragraphe, affirmé en une phrase};
\node[bloc, fill=popgreenL, align=center] (expl) at (0,-1.8){\textbf{2. Explication}\\ \small On précise, on rend l'idée intelligible (progression thématique + marqueurs explicatifs)};
\node[bloc, fill=poporangeL, align=center] (cit) at (0,-3.6){\textbf{3. Exemple / citation}\\ \small La preuve tirée de l'œuvre, entre guillemets};
\node[bloc, fill=poppinkL, align=center] (bil) at (0,-5.4){\textbf{4. Mini-conclusion}\\ \small Ce que la preuve démontre, retour à la problématique};
\draw[fleche] (idee) -- (expl);
\draw[fleche] (expl) -- (cit);
\draw[fleche] (cit) -- (bil);
\node[num] at (-5.2,0) {1};
\node[num] at (-5.2,-1.8) {2};
\node[num] at (-5.2,-3.6) {3};
\node[num] at (-5.2,-5.4) {4};
\end{tikzpicture}
\end{center}
\legende{Architecture d'un paragraphe argumentatif : quatre étages solidaires. Supprimer l'un des étages fragilise tout l'édifice.}
\end{popfigure}

\begin{aretenir}[La règle d'or du paragraphe]
Un paragraphe = \cle{une seule idée}, expliquée, prouvée, conclue. Si vous changez d'idée, changez de paragraphe. Un paragraphe trop court (une ou deux lignes) est une idée non développée ; un paragraphe d'une page entière mélange probablement plusieurs idées.
\end{aretenir}

\courssub{La citation : choisir, insérer, commenter}

\courssub{Choisir une citation pertinente}

La citation est la \cle{preuve} de votre argument. Pour être efficace, elle doit respecter trois critères :

\begin{itemize}
\item \textbf{Courte} : une phrase, deux au maximum. Une citation d'un demi-paragraphe montre que vous n'avez pas su sélectionner l'essentiel.
\item \textbf{Pertinente} : elle doit illustrer \emph{exactement} l'argument du paragraphe, pas un argument voisin. Avant de citer, demandez-vous : « Que prouve cette phrase ? »
\item \textbf{Exacte} : on cite mot pour mot, sans transformer le texte de l'auteur. Si vous n'êtes pas sûr des termes exacts, mieux vaut rapporter l'épisode avec vos propres mots (référence) plutôt que de fausser une citation.
\end{itemize}

\begin{methode}[Insérer une citation en trois temps]
\begin{enumerate}
\item \textbf{Annoncer} : préparez la citation par une phrase qui la situe (qui parle ? à quel moment de l'œuvre ? dans quelle situation ?).
\item \textbf{Citer entre guillemets} : reproduisez le texte exact, précédé des guillemets « ... », éventuellement introduit par un deux-points.
\item \textbf{Commenter} : expliquez ce que la citation prouve, en quoi elle confirme votre argument. C'est l'étape la plus importante.
\end{enumerate}
\end{methode}

\courssub{Les techniques d'insertion}

Il existe trois manières principales d'intégrer une citation à votre phrase :

\begin{center}
\begin{tabularx}{\linewidth}{|l|X|X|}
\hline
\rowcolor{popblueL} \textbf{Technique} & \textbf{Principe} & \textbf{Exemple} \\
\hline
Citation introduite par un deux-points & La phrase d'annonce est complète, le deux-points appelle la citation & Ferdinand Oyono souligne l'ironie de la scène : « Meka était tout ému, tout tremblant de joie. » \\
\hline
Citation intégrée à la phrase & La citation s'insère dans la syntaxe de votre propre phrase & Meka découvre que sa médaille n'est que « poudre aux yeux », ce qui révèle la vanité de la récompense. \\
\hline
Citation avec référence explicite & On précise l'auteur, l'œuvre, voire le chapitre & Dans \emph{Le Vieux nègre et la médaille}, Oyono écrit : « La médaille ne servait à rien. » \\
\hline
\end{tabularx}
\end{center}

\courssub{Commenter la citation : ne jamais la laisser « parlante seule »}

Une citation ne prouve rien toute seule : c'est \emph{vous} qui devez montrer ce qu'elle prouve. Le commentaire explicite le sens de la citation, souligne les mots forts, et la relie à votre argument. Il commence souvent par des formules comme : \emph{Cette réplique montre que...}, \emph{Par ces mots, l'auteur dénonce...}, \emph{On comprend ici que...}.

\begin{attention}[Erreur fréquente : citer sans commenter, ou paraphraser]
Deux fautes coûtent cher à l'examen :
\begin{itemize}
\item \textbf{Laisser la citation « parlante seule »} : juxtaposer argument et citation sans commentaire revient à dire au correcteur « débrouillez-vous avec cette preuve ». Or c'est à vous de l'exploiter.
\item \textbf{Paraphraser la citation} : répéter la citation avec d'autres mots n'est pas un commentaire. Commenter, c'est \emph{expliquer ce que la citation apporte de plus} : son ironie, sa portée critique, ce qu'elle révèle d'un personnage ou d'une société.
\end{itemize}
\end{attention}

\courssub{La transition : le pont entre les parties}

\courssub{Définition et fonction}

La \cle{transition} est une phrase ou un court paragraphe (deux ou trois lignes maximum) placé entre deux parties du développement. Elle assure la continuité logique du raisonnement : sans elle, le lecteur a l'impression de sauter d'un bloc à un autre sans lien.

\begin{propriete}[La double fonction de la transition]
Une bonne transition fait \textbf{double emploi} : elle \textbf{ferme} ce qui précède (bilan de la partie qui s'achève) et elle \textbf{ouvre} ce qui suit (annonce de la partie suivante). Une transition qui ne fait que conclure laisse le lecteur en suspens ; une transition qui ne fait qu'annoncer donne l'impression que la partie précédente est inachevée.
\end{propriete}

\begin{popfigure}
\begin{center}
\begin{tikzpicture}[
  rive/.style={rectangle, draw=popdark, thick, fill=popblueL, align=center, minimum width=3.4cm, minimum height=1.3cm, rounded corners=3pt},
  pont/.style={rectangle, draw=poporange, very thick, fill=poporangeL, align=center, minimum width=4.6cm, minimum height=1.5cm, rounded corners=6pt},
  fleche/.style={{Stealth[length=3mm]}-{Stealth[length=3mm]}, thick, popdark}
]
\node[rive, align=center] (p1) at (-4.6,0){\textbf{Partie I}\\ \small Thèse / premier aspect};
\node[rive, align=center] (p2) at (4.6,0){\textbf{Partie II}\\ \small Antithèse / second aspect};
\node[pont, align=center] (pont) at (0,0){\textbf{TRANSITION}\\ \small Bilan de I + annonce de II};
\draw[fleche] (p1) -- (pont);
\draw[fleche] (pont) -- (p2);
\node[align=center, font=\small, text=popmuted] at (0,-1.5) {Le lecteur traverse le raisonnement sans rupture};
\end{tikzpicture}
\end{center}
\legende{La transition comme un pont : elle s'appuie sur la partie qui se termine (bilan) et débouche sur la partie qui commence (annonce).}
\end{popfigure}

\courssub{Les trois types de transitions}

\begin{itemize}
\item \textbf{La transition « bilan + annonce »} : on résume l'acquis de la partie, puis on annonce la suite. \emph{Exemple : « Ainsi, la médaille apparaît d'abord comme un honneur enviable. Mais cet honneur est-il réel ? C'est ce que nous allons examiner. »}
\item \textbf{La question relais} : on conclut par une question dont la partie suivante apportera la réponse. \emph{Exemple : « La récompense flatte donc l'orgueil du vieux Meka. Pourtant, que cache cette générosité soudaine du commandant ? »}
\item \textbf{La concession (« Certes... mais... »)} : on accorde provisoirement un point à la thèse adverse avant de la dépasser. \emph{Exemple : « Certes, Meka reçoit les honneurs de l'administration coloniale ; mais cette reconnaissance n'est qu'un leurre, comme le montrera la suite du récit. »}
\end{itemize}

\begin{attention}[Pièges de la transition]
\begin{itemize}
\item Ne transformez pas la transition en nouveau paragraphe argumentatif : elle ne contient \textbf{ni argument nouveau, ni citation}.
\item Évitez les transitions mécaniques et vides (« Nous venons de voir la première partie, nous allons voir la deuxième ») : elles ne disent rien du contenu.
\item La transition se place \emph{entre} les parties, pas entre chaque paragraphe : entre deux paragraphes d'une même partie, un simple connecteur logique (\emph{de plus, en outre, par ailleurs}) suffit.
\end{itemize}
\end{attention}

\courssub{Le texte explicatif : rendre le raisonnement intelligible}

Dans le développement, vous ne vous contentez pas d'affirmer : vous \emph{expliquez}. C'est la fonction du texte explicatif, qui structure l'étage 2 (l'explication) de votre paragraphe.

\begin{definition}[Le texte explicatif]
Le \cle{texte explicatif} est un texte qui rend compte de \textbf{causes}, de \textbf{mécanismes} ou de \textbf{raisons} afin de rendre un phénomène intelligible au lecteur. Il s'appuie sur des relations logiques — cause, but, conséquence — signalées par des \textbf{marqueurs explicatifs} : \emph{parce que, car, en effet, puisque, c'est pourquoi, ainsi, donc, de sorte que, afin de, c'est la raison pour laquelle}.
\end{definition}

Expliquer, c'est répondre à la question « pourquoi ? » ou « comment ? ». Dans un paragraphe de dissertation, l'explication (deuxième étage) est précisément un mini-texte explicatif : elle décompose le mécanisme de votre argument en assurant la progression thématique.

\begin{exemplebox}[Marqueurs explicatifs et progression thématique à l'œuvre]
« Meka se réjouit de recevoir la médaille \textbf{parce qu'}il y voit la reconnaissance de toute une vie de soumission [Cause]. \textbf{En effet}, il a toujours obéi à l'administration coloniale [Justification / Preuve]. \textbf{C'est pourquoi} la remise de la médaille lui semble être le couronnement de son existence [Conséquence]. » — Chaque marqueur tisse un lien logique visible ; la progression thématique (médaille $\rightarrow$ soumission $\rightarrow$ obéissance $\rightarrow$ couronnement) guide le lecteur pas à pas.
\end{exemplebox}

\begin{savaistu}[La scène la plus citée des examens camerounais]
Dans \emph{Le Vieux nègre et la médaille} de Ferdinand Oyono, la scène de la \textbf{remise de la médaille} (le 14 juillet, en présence du commandant et des notables) est la scène la plus citée aux examens camerounais. Il faut la connaître précisément : le discours pompeux du commandant, l'attitude fière et naïve de Meka, la foule des invités, et l'ironie de l'auteur qui transparaît dans chaque détail. Apprenez deux ou trois citations courtes de cette scène : elles serviront dans presque tous les sujets sur l'œuvre.
\end{savaistu}

\courssub{Application guidée : un paragraphe complet avec citation et transition}

\begin{exoresolu}[Rédiger un paragraphe argumentatif sur \emph{Le Vieux nègre et la médaille}]
\textbf{Énoncé.} Sujet : \emph{La médaille donnée à Meka est-elle une véritable récompense ?} Rédigez un paragraphe argumentatif (première partie : la médaille comme honneur apparent) comportant une citation insérée et commentée, puis la transition vers la deuxième partie (la médaille comme leurre).
\tcblower
\textbf{Solution détaillée.}

\emph{Paragraphe :}

Dans un premier temps, la médaille apparaît comme un honneur exceptionnel pour le vieux Meka, \textbf{parce qu'}elle couronne publiquement toute une vie de soumission à l'administration coloniale. \textbf{En effet}, Meka a toujours été un sujet docile : il a donné ses terres, ses fils à la guerre, sa foi au christianisme. Le jour de la cérémonie, sa fierté éclate lorsque Oyono écrit : « Meka était tout ému, tout tremblant de joie ». Cette réplique montre que le vieil homme vit cet instant comme l'aboutissement de son existence : l'émotion et le tremblement traduisent un bonheur sincère, celui d'être enfin reconnu par ceux qu'il a servis. La médaille semble donc bien, à ce stade du récit, une récompense authentique et méritée.

\emph{Transition (bilan + question relais) :}

Ainsi, la médaille comble d'abord l'orgueil du vieux nègre et flatte sa vanité. Mais cette gloire d'un jour résiste-t-elle à l'épreuve de la réalité ? C'est ce que révèle la suite du roman, où la récompense se transforme en instrument d'humiliation.

\emph{Commentaire de la copie :} le paragraphe respecte les quatre étages (idée directrice, explication avec marqueurs explicatifs et progression thématique, citation annoncée, citée entre guillemets puis commentée, mini-conclusion). La transition fait le bilan de la première partie et annonce la seconde par une question relais.
\end{exoresolu}

\begin{exercice}[À vous de jouer]
Sur le même sujet, rédigez le premier paragraphe de la deuxième partie (la médaille comme leurre) : idée directrice, explication avec deux marqueurs explicatifs, citation courte de l'œuvre insérée et commentée, mini-conclusion. Rédigez ensuite la transition qui mènerait à la conclusion de la dissertation.
\end{exercice}

\begin{corrige}[Exercice : éléments de réponse]
\emph{Paragraphe attendu (modèle) :} Cependant, la médaille se révèle être un leurre, \textbf{car} elle n'apporte à Meka aucun bien réel. \textbf{En effet}, dès la soirée qui suit la cérémonie, le vieil homme est arrêté et humilié par les mêmes forces coloniales qui venaient de l'honorer. Oyono note ironiquement : « La médaille ne servait à rien. » Cette constatation dénonce l'absurdité de la récompense : purement symbolique, elle masque la domination réelle exercée sur le vieux nègre. La médaille n'est donc qu'un instrument de la propagande coloniale.

\emph{Transition attendue :} Ainsi, derrière l'honneur apparent se cache une manipulation : la médaille glorifie celui qui la donne bien plus que celui qui la reçoit. Il reste à se demander ce que la chute finale du roman révèle de la prise de conscience de Meka.

\emph{Barème indicatif :} idée directrice claire (1 pt), explication avec marqueurs explicatifs et progression thématique (1 pt), citation exacte, insérée et commentée (2 pts), mini-conclusion (1 pt), transition à double fonction (1 pt).
\end{corrige}

\begin{aretenir}[L'essentiel de la section]
\begin{itemize}
\item Le développement compte deux ou trois parties de deux ou trois paragraphes ; il occupe environ les deux tiers de la copie.
\item Chaque paragraphe suit quatre étages : idée directrice, explication (progression thématique + marqueurs explicatifs), citation, mini-conclusion.
\item La citation s'insère en trois temps : annoncer, citer entre guillemets, commenter. Jamais de citation laissée seule.
\item La transition fait double emploi : elle ferme une partie (bilan) et ouvre la suivante (annonce), sous forme de bilan + annonce, de question relais ou de concession.
\item Le texte explicatif rend le raisonnement intelligible grâce aux marqueurs de cause, de but et de conséquence (\emph{parce que, en effet, c'est pourquoi, ainsi}) et à une progression thématique maîtrisée (linéaire, constante, dérivée).
\end{itemize}
\end{aretenir}

\coursec{Rédiger la conclusion}

Après avoir construit le développement — paragraphes argumentés, citations insérées avec rigueur et transitions fluides assurant la progression thématique — il reste à poser la pierre angulaire de l'édifice : la conclusion. C'est le dernier contact avec le correcteur, l'ultime occasion de prouver que la démonstration tient debout et que la réflexion a abouti. Une conclusion bâclée ruine souvent une copie par ailleurs solide ; une conclusion maîtrisée, au contraire, scelle la réussite de l'épreuve.

\courssub{La fonction et la structure de la conclusion}

La conclusion n'est pas un simple résumé, ni une politesse de fin de copie. Elle remplit une fonction \textbf{rhétorique} et \textbf{cognitive} précise : elle doit clore la réflexion en donnant au lecteur le sentiment d'une pensée achevée, tout en manifestant l'intelligence de la problématique. Elle opère un mouvement \textbf{centrifuge} : du particulier (les arguments du développement) vers le général (la réponse, puis l'horizon élargi). C'est le miroir exact de l'introduction, qui allait du général (amorce, sujet) vers le particulier (problématique, annonce du plan).

\begin{definition}[La conclusion de dissertation]
La conclusion est la partie terminale de la dissertation qui synthétise la démonstration, apporte une réponse définitive et nuancée à la problématique posée en introduction, et ouvre la réflexion sur une perspective nouvelle sans la traiter. Elle se compose de trois temps obligatoires : le bilan, la réponse à la problématique, l'ouverture.
\end{definition}

\begin{definition}[L'ouverture]
Dernière phrase (ou dernier groupe de phrases) de la conclusion, l'ouverture prolonge la réflexion vers un sujet connexe, une question voisine ou un enjeu actuel, sans jamais entamer un nouveau développement. Elle montre que la problématique traitée s'inscrit dans un champ plus vaste (ex. : enjeux sociétaux, littéraires, philosophiques, perspectives d'avenir).
\end{definition}

\courssub{Les trois composantes de la conclusion}

Chaque composante répond à une exigence logique distincte. Leur enchaînement doit être fluide, sans rupture brutale.

\textbf{1. Le bilan (ou synthèse) :} C'est le rappel \emph{synthétique} des grandes étapes du raisonnement. On ne répète pas les arguments un par un (ce serait une redite fastidieuse), on en restitue l'architecture : « Nous avons vu d'abord que... ; ensuite que... ; enfin que... ». Le bilan valide la cohérence du plan annoncé.

\textbf{2. La réponse à la problématique :} C'est le cœur de la conclusion. Elle doit être \textbf{claire}, \textbf{nuancée} et \textbf{directement issue} du développement. On évite le « Oui » ou le « Non » binaire. On formule une thèse aboutie : « Au terme de cette analyse, il apparaît que... » ou « Force est de constater que... ». Cette réponse doit reprendre les termes mêmes de la problématique pour prouver qu'on n'a pas dévié.

\textbf{3. L'ouverture :} Elle élargit l'horizon. Elle peut prendre plusieurs formes : une question nouvelle (sans y répondre), une perspective historique ou prospective (ex. : « La place des traditions dans le Cameroun de 2035 »), une référence à une autre œuvre, un autre auteur, un autre contexte. Elle prouve que le candidat a une culture et une vision au-delà du sujet imposé.

\begin{methode}[Rédiger une conclusion en trois temps]
\begin{enumerate}
    \item \textbf{Relire l'introduction et le développement} : identifier la problématique exacte et les 2 ou 3 grandes parties du plan.
    \item \textbf{Rédiger le bilan} (1 à 2 phrases) : utiliser des connecteurs de synthèse (\emph{au total, en définitive, pour résumer, en somme}) pour articuler les grandes étapes sans citer les exemples.
    \item \textbf{Formuler la réponse} (1 à 2 phrases) : reprendre les mots-clés de la problématique ; affirmer une position nuancée (\emph{partiellement, dans une large mesure, toutefois, néanmoins}).
    \item \textbf{Construire l'ouverture} (1 phrase) : choisir un prolongement pertinent (autre œuvre, actualité, enjeu futur, question philosophique) ; ne pas développer.
    \item \textbf{Relier le tout} : assurer la fluidité par des connecteurs logiques (\emph{ainsi, par conséquent, dès lors, c'est pourquoi}).
    \item \textbf{Vérifier la longueur} : 5 à 8 lignes manuscrites maximum (environ 60-100 mots).
\end{enumerate}
\end{methode}

\courssub{Correspondance introduction – conclusion : le miroir}

La conclusion est le reflet fidèle de l'introduction. Chaque élément de l'introduction trouve son écho symétrique dans la conclusion. Le tableau ci-dessous formalise cette correspondance structurelle, garantie de la cohérence globale de la copie.

\begin{popfigure}
\begin{tikzpicture}[
    cell/.style={rectangle, draw=popblue, thick, minimum height=1.2cm, minimum width=3.5cm, align=center, font=\small, inner sep=4pt},
    header/.style={fill=popblue, text=white, font=\bfseries\small, minimum height=1cm},
    arrow/.style={->, >=Stealth, thick, poporange},
    title/.style={font=\bfseries\normalsize, text=popdark}
]
% En-têtes
\node[header] (h1) at (0,0) {Introduction (Entonnoir)};
\node[header] (h2) at (7.5,0) {Conclusion (Entonnoir inversé)};
% Lignes Introduction
\node[cell, fill=popblueL, align=center] (i1) at (0,-1.5){Amorce / Sujet posé\\ (Général)};
\node[cell, fill=popblueL, align=center] (i2) at (0,-3){Problématique\\ (Question précise)};
\node[cell, fill=popblueL, align=center] (i3) at (0,-4.5){Annonce du plan\\ (2-3 grandes parties)};
% Lignes Conclusion
\node[cell, fill=poporangeL, align=center] (c1) at (7.5,-1.5){Ouverture\\ (Horizon élargi / Général)};
\node[cell, fill=poporangeL, align=center] (c2) at (7.5,-3){Réponse à la problématique\\ (Thèse aboutie / Précise)};
\node[cell, fill=poporangeL, align=center] (c3) at (7.5,-4.5){Bilan / Synthèse du plan\\ (Rappel des grandes étapes)};
% Flèches de correspondance
\draw[arrow] (i1.east) -- node[above, font=\tiny, sloped] {Miroir} (c1.west);
\draw[arrow] (i2.east) -- node[above, font=\tiny, sloped] {Répondre} (c2.west);
\draw[arrow] (i3.east) -- node[above, font=\tiny, sloped] {Synthétiser} (c3.west);
% Flèches internes Intro (descendantes)
\draw[popmuted, thick, ->] (i1.south) -- (i2.north);
\draw[popmuted, thick, ->] (i2.south) -- (i3.north);
% Flèches internes Conclusion (montantes)
\draw[popmuted, thick, ->] (c3.north) -- (c2.south);
\draw[popmuted, thick, ->] (c2.north) -- (c1.south);
% Légendes de sens
\node[title, rotate=90] at (-3.5, -3) {Du général au particulier};
\node[title, rotate=-90] at (11, -3) {Du particulier au général};
\end{tikzpicture}
\legende{Tableau de correspondance symétrique entre l'introduction et la conclusion. La flèche centrale indique le sens de la pensée : l'introduction resserre (entonnoir), la conclusion élargit (entonnoir inversé).}
\end{popfigure}

\courssub{Schémas visuels : l'entonnoir inversé}

Pour ancrer visuellement la structure de la conclusion, le schéma en entonnoir inversé est le modèle mental le plus efficace. Il montre la progression du \textbf{étroit} (le bilan, centré sur le sujet précis) vers le \textbf{large} (l'ouverture, vers l'universel).

\begin{popfigure}
\begin{tikzpicture}[
    funnel/.style={rectangle, rounded corners, draw=popdark, thick, minimum height=1.5cm, text width=6cm, align=center, font=\small},
    funnel top/.style={funnel, fill=popblueL, text width=8cm, minimum width=8cm},
    funnel mid/.style={funnel, fill=poporangeL, text width=5cm, minimum width=5cm},
    funnel bot/.style={funnel, fill=popgreenL, text width=2.5cm, minimum width=2.5cm},
    lab/.style={font=\bfseries\small, text=popdark},
    desc/.style={font=\tiny, text=popmuted, align=center, text width=5cm}
]
% Nœuds de l'entonnoir inversé (du bas vers le haut visuellement, mais logique du haut vers le bas)
\node[funnel bot, label=below:{BILAN}, align=center] (bilan) at (0,0){\textbf{Synthèse des grandes étapes}\\\emph{(« Nous avons montré que... »)}};
\node[funnel mid, above=0.2cm of bilan, label=above:{RÉPONSE}, align=center] (reponse){\textbf{Réponse à la problématique}\\\emph{(« Il apparaît que... »)}};
\node[funnel top, above=0.2cm of reponse, label=above:{OUVERTURE}, align=center] (ouverture){\textbf{Élargissement / Perspective}\\\emph{(« Cela invite à s'interroger sur... »)}};
% Flèches montantes (logique de construction)
\draw[poporange, thick, ->, >=Stealth] (bilan.north) -- (reponse.south) node[midway, right=0.2cm, desc] {Conduit à};
\draw[poporange, thick, ->, >=Stealth] (reponse.north) -- (ouverture.south) node[midway, right=0.2cm, desc] {Appelle};
% Annotations latérales
\node[lab, left=1.5cm of bilan] (l1) {Étroit};
\node[lab, left=1.5cm of reponse] (l2) {Intermédiaire};
\node[lab, left=1.5cm of ouverture] (l3) {Large};
\draw[popmuted, thick] (l1.west) -- (l3.west) node[midway, left=0.5cm, font=\small\bfseries, rotate=90] {Amplitude de la réflexion};
% Titre
\node[font=\bfseries\normalsize, text=popdark, above=1cm of ouverture] {Structure de la conclusion : l'entonnoir inversé};
\end{tikzpicture}
\legende{Schéma de la conclusion en entonnoir inversé : on part du bilan (étroit, centré sur le sujet) pour aboutir à l'ouverture (large, horizon élargi). C'est le mouvement inverse de l'introduction.}
\end{popfigure}

\courssub{Longueur, amélioration et pièges à éviter}

La conclusion doit tenir en \textbf{5 à 8 lignes} manuscrites (environ 60 à 100 mots). Au-delà, elle devient un développement déguisé ; en deçà, elle paraît expédiée. L'amélioration systématique passe par une vérification croisée avec l'introduction.

\begin{attention}[Erreur fréquente : introduire un nouvel argument]
C'est la faute capitale. Écrire dans la conclusion : « \emph{Par ailleurs, on pourrait ajouter que...} » ou « \emph{En outre, l'exemple de X prouve que...} » déséquilibre la copie. Le développement est clos. La conclusion ne fait que \textbf{regarder en arrière} (bilan, réponse) et \textbf{jeter un regard latéral} (ouverture), jamais avancer de nouveau. Si une idée forte vous vient à la fin, c'est qu'elle manquait au développement : tant pis, ne l'insérez pas en conclusion.
\end{attention}

\begin{attention}[Autres pièges classiques]
\begin{itemize}
    \item \textbf{La redite servile :} recopier l'annonce du plan mot pour mot au lieu de synthétiser.
    \item \textbf{La réponse absente ou floue :} « \emph{Ce sujet est complexe} » n'est pas une réponse. Il faut trancher.
    \item \textbf{L'ouverture hors-sujet :} ouvrir sur un thème sans lien logique avec la problématique (ex. : ouvrir sur l'écologie après une dissertation sur le théâtre classique).
    \item \textbf{La phrase « accroche » finale :} terminer par une citation non analysée ou une formule creuse (« \emph{L'avenir nous le dira} »).
\end{itemize}
\end{attention}

\courssub{Production guidée : conclusion sur le sujet de la médaille de Meka}

Reprenons le sujet travaillé en Section 3 (Introduction) et Section 4 (Développement) : \emph{« La médaille remise au Vieux Nègre (Meka) dans \og Le Vieux Nègre et la Médaille \fg de Ferdinand Oyono : quelle valeur symbolique et quelle ironie ? »}

\emph{Rappel de la problématique (Section 3) :} « En quoi la remise de la médaille au Vieux Nègre Meka révèle-t-elle l'absurdité du système colonial et la tragédie de l'aliénation, tout en conférant au personnage une dignité involontaire ? »

\emph{Rappel du plan (Section 4) :}
\begin{enumerate}
    \item Une distinction vide de sens : l'ironie d'une récompense pour services imaginaires.
    \item L'aliénation consentie : Meka, acteur de sa propre duperie par soif de reconnaissance.
    \item La subversion finale : la médaille, symbole retourné, devient le témoin muet de la folie coloniale.
\end{enumerate}

\begin{exemplebox}[Conclusion rédigée : « La médaille de Meka »]
\textbf{Bilan :} Au terme de cette étude, il apparaît que la médaille remise à Meka fonctionne comme un miroir grossissant de l'entreprise coloniale : une distinction vide, octroyée pour des « services » inexistants, qui dévoile le mépris sous-jacent de l'administration. \\
\textbf{Réponse :} Loin d'honorer le vieil homme, cette décoration scelle son aliénation tragique ; Meka, en l'acceptant avec fierté, devient le complice inconscient de sa propre duperie, illustrant la violence symbolique d'un système qui nie l'Autre en le décorant. \\
\textbf{Ouverture :} Cette ironie cruelle invite à s'interroger, aujourd'hui encore, sur la portée réelle des distinctions honorifiques au Cameroun : sont-elles le reflet d'un mérite authentique ou le prolongement d'une logique de clientélisme et de reconnaissance symbolique détournée ?
\end{exemplebox}

\begin{exoresolu}[Analyse de la conclusion produite]
\begin{itemize}
    \item \textbf{Bilan (2 lignes) :} Synthétise les 3 parties sans les lister mécaniquement. Mots-clés : « miroir grossissant », « distinction vide », « mépris », « aliénation ». \checkmark
    \item \textbf{Réponse (2 lignes) :} Reprend « aliénation », « duperie », « violence symbolique » de la problématique. Position nuancée : « complice inconscient », « tragique ». \checkmark
    \item \textbf{Ouverture (2 lignes) :} Prolongement vers l'actualité camerounaise (distinctions honorifiques, clientélisme). Question ouverte, sans développement. Lien logique fort : la médaille d'hier $\rightarrow$ les décorations d'aujourd'hui. \checkmark
    \item \textbf{Longueur :} 6 lignes manuscrites environ. Respect du format. \checkmark
    \item \textbf{Cohérence intro/conclusion :} Problématique (absurdité, aliénation, dignité) $\leftrightarrow$ Réponse (mépris, aliénation tragique, violence symbolique). Parfait miroir. \checkmark
\end{itemize}
\end{exoresolu}

\courssub{Exercice d'application : rédaction d'une conclusion}

\begin{exercice}[Sujet : « Le rôle du griot dans \og Soundiata \fg ou \og Sundiata \fg (selon version étudiée) : gardien de la mémoire ou acteur du pouvoir ? »]
À partir des éléments ci-dessous, rédigez une conclusion complète (bilan, réponse, ouverture) en respectant la méthode en 3 temps et la limite de 8 lignes.

\textbf{Problématique :} « Dans quelle mesure le griot, par sa parole, assure-t-il la légitimité du pouvoir mandingue tout en préservant l'identité collective ? »

\textbf{Plan du développement :}
\begin{enumerate}
    \item Le griot, dépositaire de la parole ancestrale : garant de la mémoire et de la généalogie.
    \item Le griot, stratège politique : son verbe fonde la légitimité de Soundiata (ex. : Balla Fasséké).
    \item La parole griotique comme ciment social : au-delà du roi, elle unit la communauté par les valeurs (honneur, courage, alliance).
\end{enumerate}

\textbf{Contraintes :} Utiliser au moins deux connecteurs de synthèse (\emph{au total, en définitive, pour résumer}) et un connecteur de conséquence (\emph{ainsi, par conséquent, c'est pourquoi}) pour la réponse. L'ouverture doit porter sur la place de l'oralité dans l'Afrique contemporaine (médias, éducation, politique).
\end{exercice}

\begin{corrige}[Exercice n°1 : Conclusion sur le griot]
\textbf{Bilan :} \emph{Au total}, l'analyse du rôle du griot dans l'épopée mandingue révèle une figure aux fonctions indissociables : gardien sacré de la mémoire collective par la transmission des généalogies et des hauts faits, il est aussi l'artisan discret mais indispensable de la légitimité politique du roi. \\
\textbf{Réponse :} \emph{C'est pourquoi} le griot ne saurait être réduit à un simple historien ; il est l'opérateur symbolique par qui le pouvoir devient droit et la communauté devient peuple, conjuguant ainsi l'autorité temporelle et l'éternité des ancêtres. \\
\textbf{Ouverture :} Cette centralité de la parole maîtrisée interroge la place de l'oralité dans nos sociétés africaines contemporaines, tiraillées entre la tentation de l'écrit unique et la vitalité des nouveaux griots que sont les journalistes, rappeurs ou animateurs de la parole publique.
\end{corrige}

\coursec{Lecture méthodique du Vieux nègre et la médaille (n° 2 et 3) et remédiation}

Après avoir maîtrisé la rédaction de la conclusion, nous devons maintenant ancrer nos acquis dans l'étude d'une œuvre au programme : \emph{Le Vieux nègre et la médaille} de Ferdinand Oyono. Cette œuvre n'est pas seulement un roman à lire ; c'est un réservoir d'arguments, de citations et de procédés stylistiques que vous devez savoir mobiliser dans une dissertation. Les lectures méthodiques n° 2 et 3 portent sur deux moments charnières : la remise de la médaille (apogée de l'illusion) et ses conséquences (chute dans la réalité). La remédiation finale vous permettra de corriger vos défauts récurrents avant l'examen.

\courssub{Transition : de la technique dissertation à l'exploitation de l'œuvre}

\emph{Le Vieux nègre et la médaille} est l'œuvre de référence pour illustrer les thèmes de la colonisation, de l'aliénation, de la dignité et de la révolte. Dans une dissertation, vous ne \textbf{racontez} pas l'histoire ; vous \textbf{exploitez} des extraits précis pour prouver vos arguments. Les deux lectures méthodiques qui suivent vous apprennent à transformer une page de roman en matériel argumentatif de haut niveau.

\courssub{Lecture méthodique n° 2 : La remise de la médaille (Chapitre 6) — L'apogée de l'ironie}

\textbf{Contexte et situation d'énonciation.}\par
Nous sommes au chapitre 6. Meka, après des mois d'attente et de démarches humiliantes, reçoit enfin la médaille de la part de l'Administrateur. La scène se déroule en public, devant la chefferie, les notables et les Blancs.
\begin{itemize}
    \item \textbf{Locuteur :} Narrateur omniscient, focalisation interne sur Meka (on accède à ses pensées, sa fierté naïve).
    \item \textbf{Allocutaire :} Le lecteur, mis en position de juge lucide.
    \item \textbf{Ton :} Solennel en surface (cérémonie officielle), profondément \cle{ironique} en profondeur.
\end{itemize}

\textbf{Analyse des champs lexicaux : le choc des mondes.}\par
Oyono construit la scène sur l'opposition de deux univers lexicaux :
\begin{center}
\begin{tabularx}{\linewidth}{|l|X|X|}
\hline
\rowcolor{popblueL}
\textbf{Champ lexical} & \textbf{Mots-clés (extrait)} & \textbf{Effet produit} \\
\hline
La grandeur officielle & « cérémonie », « décoré », « distinction », « mérite », « autorité », « République » & Légitime le pouvoir colonial, langage vide de sens réel. \\
\hline
La misère réelle & « loques », « poussière », « pieds nus », « ventre vide », « cases en terre » & Rappelle la condition véritable du « vieux nègre », décalage cruel. \\
\hline
Le corps de Meka & « tremblait », « gonflé », « yeux brillants », « médaille au cou » & Corps théâtralisé, objet de la mise en scène coloniale. \\
\hline
\end{tabularx}
\end{center}

\begin{definition}[Ironie (arme majeure d'Oyono)]
L'\cle{ironie} est un procédé stylistique qui consiste à \textbf{dire le contraire de ce que l'on pense} (ou à présenter une situation comme positive alors qu'elle est humiliante) pour \textbf{critiquer} avec plus de force. Le lecteur comprend le second degré ; le personnage (Meka) y reste aveugle. C'est une \cle{ironie dramatique} : le spectateur sait, l'acteur ignore.
\end{definition}

\textbf{Exemple d'exploitation en dissertation.}\par
\begin{exemplebox}[Citation exploitable : thème de l'illusion]
\textit{« Meka, médaille au cou, se croyait devenu l'égal des Blancs. »} (Chapitre 6)
\begin{itemize}
    \item \textbf{Analyse :} Le verbe « se croyait » marque l'erreur, l'illusion. La médaille (objet matériel) crée une fausse égalité symbolique.
    \item \textbf{Insertion dans un paragraphe :} « La colonisation repose sur une illusion entretenue : le colonisé pense accéder à la dignité par les symboles du colonisateur. Ainsi, Meka, "médaille au cou, se croyait devenu l'égal des Blancs", alors qu'il n'est qu'un instrument de plus de l'administration. »
\end{itemize}
\end{exemplebox}

\begin{methode}[Comment exploiter cet extrait en 3 étapes]
1. \textbf{Contextualiser en une phrase :} « Lors de la cérémonie officielle de remise de la médaille (ch. 6), Meka croit toucher au but. »
2. \textbf{Citer court et précis :} « "Meka, médaille au cou, se croyait devenu l'égal des Blancs." »
3. \textbf{Commenter le procédé (l'ironie) :} « L'ironie d'Oyono dévoile le piège : la médaille ne change pas la nature du rapport de force ; elle l'entérine. Le "vieux nègre" reste un subalterne décoré. »
\end{methode}

\begin{exoresolu}[Analyse d'un court extrait : l'arrivée de l'Administrateur]
\textbf{Énoncé :} Étudier l'ironie dans ce passage : « L'Administrateur arriva. Il était gros, rouge, suant. Il sourit. Il serra la main de Meka. Il lui épingla la médaille sur la poitrine. Meka ne respirait plus. Il était heureux. »
\textbf{Solution détaillée :}
\begin{enumerate}
    \item \textbf{Description physique dégradante :} « gros, rouge, suant » — le colonisateur n'a rien de majestueux, il est humain, lourd, transpirant. L'ironie commence ici : le représentant de la « grandeur » française est physiquement ridicule.
    \item \textbf{Gestes mécaniques :} « sourit », « serra la main », « épingla » — verbes d'action rapides, sans affect réel. C'est une routine administrative, pas une reconnaissance humaine.
    \item \textbf{Contraste interne :} « Meka ne respirait plus. Il était heureux. » — L'apnée marque la tension sacrée ; le bonheur est naïf. Le lecteur perçoit le vide de l'instant : Meka est heureux \emph{parce qu'il ne comprend pas}.
    \item \textbf{Conclusion argumentative :} Cet extrait prouve que la reconnaissance coloniale est une \emph{mascarade} (thème : illusion/aliénation).
\end{enumerate}
\end{exoresolu}

\courssub{Lecture méthodique n° 3 : Les conséquences de la médaille (Chapitres 7 à 9) — Fierté, humiliation, révolte}

\textbf{La bascule : de l'euphorie à la désillusion.}\par
Après la cérémonie, Meka rentre au village. La médaille devient un aimant à ennuis.
\begin{itemize}
    \item \textbf{Fierté éphémère (Ch. 7) :} Meka exige le respect, fait appeler « Monsieur le Décoré ». Il inverse la hiérarchie villageoise (les jeunes doivent le saluer). \emph{Registre : satirique / pathétique.}
    \item \textbf{Humiliation administrative (Ch. 8) :} Pour toucher la prime (1000 francs), il doit subir l'interrogatoire du Chef de Poste, payer des pots-de-vin, subir le mépris du commis blanc. La médaille ne lui donne \textbf{aucun droit réel}.
    \item \textbf{Révolte finale (Ch. 9) :} Le vol de la médaille par Kelara (sa femme) et la découverte de la supercherie (la médaille est en toc, la prime minable) provoquent l'effondrement. Meka arrache la médaille : « Il la jeta dans la latrine. »
\end{itemize}

\textbf{Procédés stylistiques exploitables :}\par
\begin{propriete}[Registre satirique et antithèse]
Oyono utilise la \cle{satire} pour ridiculiser la société coloniale ET la naïveté des colonisés.
\begin{itemize}
    \item \textbf{Antithèse structurelle :} Cérémonie fastueuse (Ch. 6) / Réalité sordide (Ch. 8-9).
    \item \textbf{Gradation dans l'humiliation :} Attente $\rightarrow$ Interrogatoire $\rightarrow$ Pot-de-vin $\rightarrow$ Prime dérisoire.
    \item \textbf{Chute finale (catharsis) :} Le jet de la médaille dans la latrine est un acte \textbf{symbolique fort} : rejet de l'aliénation, reprise de dignité par le mépris de l'objet trompeur.
\end{itemize}
\end{propriete}

\begin{exemplebox}[Citation exploitable : thème de la révolte / dignité retrouvée]
\textit{« Il la jeta dans la latrine. »} (Chapitre 9, dernière phrase du roman)
\begin{itemize}
    \item \textbf{Analyse :} Verbe d'action brut (« jeter »), lieu méprisable (« latrine »). Geste libérateur. La médaille, symbole de l'asservissement consenti, finit dans les excréments.
    \item \textbf{Thèmes :} Rejet de l'illusion, reconquête de la lucidité, révolte silencieuse mais radicale.
\end{itemize}
\end{exemplebox}

\begin{attention}[Erreur fréquente : Raconter au lieu d'argumenter]
\emph{Interdit :} « Dans le chapitre 7, Meka rentre au village et tout le monde le salue. Il est content. Au chapitre 8, il va voir le chef de poste pour l'argent mais on lui demande de l'argent. Au chapitre 9, sa femme vole la médaille et il la jette aux toilettes. » $\rightarrow$ \textbf{Note : 0/20 en argumentation.}
\emph{Exigé :} « La médaille, loin d'apporter la puissance, expose Meka à la prédation (vol de Kelara) et à l'extorsion (chef de poste). Le geste final — "Il la jeta dans la latrine" — signe le refus définitif de l'ordre colonial. »
\end{attention}

\courssub{Constitution d'une banque personnelle de citations classées par thème}

Pour le Probatoire, vous devez avoir \textbf{3 à 4 citations courtes par thème}, mémorisées avec leur numéro de chapitre. Voici un modèle à compléter et à apprendre par cœur.

\begin{popfigure}
\legende{Banque de citations type — Le Vieux nègre et la médaille (à mémoriser pour la dissertation)}
\begin{center}
\begin{tabularx}{\linewidth}{|>{\bfseries}p{2.5cm}|p{1.5cm}|>{\itshape}X|X|}
\hline
\rowcolor{popblue}
\textcolor{white}{\textbf{Thème}} & \textcolor{white}{\textbf{Chapitre}} & \textcolor{white}{\textbf{Citation (extrait)}} & \textcolor{white}{\textbf{Exploitation en dissertation}} \\
\hline
\rowcolor{popblueL}
Colonisation / Aliénation & Ch. 1, 6 & « On ne peut pas servir deux maîtres à la fois. » (Père de Meka, Ch.1) & Thèse : la collaboration est impossible. Argument : le choix de Meka (servir l'admin) le perd. \\
\hline
Illusion / Leurres & Ch. 6 & « Meka, médaille au cou, se croyait devenu l'égal des Blancs. » & Ironie dramatique. La médaille crée une fausse égalité symbolique. À opposer à la réalité Ch.8. \\
\hline
\rowcolor{popblueL}
Dignité / Perte de soi & Ch. 7, 8 & « Il avait faim. Il avait soif. Mais il avait la médaille. » (Ch.7) & Antithèse besoins vitaux / symbole vain. La médaille ne nourrit pas. Critique de l'aliénation matérielle. \\
\hline
Révolte / Lucidité & Ch. 9 & « Il la jeta dans la latrine. » (Dernière phrase) & Acte symbolique fort. Rejet de l'ordre colonial. Ouverture : révolte individuelle $\rightarrow$ conscience collective ? \\
\hline
\end{tabularx}
\end{center}
\end{popfigure}

\begin{methode}[Construire SA propre fiche (méthode « 3C »)]
1. \textbf{Choisir} 3 citations \textbf{courtes} (max 15 mots) par grand thème du programme.
2. \textbf{Contextualiser} : noter le chapitre et qui parle (narrateur / personnage).
3. \textbf{Commenter} : écrire \textbf{une phrase type} d'exploitation (lien thème + procédé + portée).
\begin{aretenir}[Règle d'or]
En dissertation, une citation sans commentaire = 0 point. La citation est une \textbf{preuve}, pas une décoration. Toujours appliquer : \textbf{Citer $\rightarrow$ Analyser $\rightarrow$ Lier à l'argument}.
\end{aretenir}
\end{methode}

\courssub{Le compte rendu : présenter sa production et sa démarche}

Le programme exige un \textbf{compte rendu} (oral ou écrit) de votre dissertation ou de vos exercices d'intégration. C'est une compétence transversale : \emph{« Rédiger et justifier une démarche »}.

\begin{methode}[Structure du compte rendu oral (3 minutes max)]
1. \textbf{Sujet et problématique :} « Le sujet était : "La médaille, symbole de reconnaissance ou outil d'aliénation ?" Ma problématique : "Dans quelle mesure la distinction honorifique offerte au colonisé masque-t-elle sa soumission réelle ?" »
2. \textbf{Plan choisi et justification :} « Plan dialectique : I. La médaille comme leurre (illusion) $\rightarrow$ II. La médaille comme instrument de domination (réalité) $\rightarrow$ III. Le rejet de la médaille comme acte de lucidité (dépassement). J'ai choisi ce plan car l'œuvre suit cette chronologie. »
3. \textbf{Point fort / Difficulté :} « Mon point fort : l'exploitation de l'ironie au Ch.6. Ma difficulté : faire des transitions fluides entre les paragraphes sur le thème de la dignité. »
4. \textbf{Axe de progression :} « Pour la prochaine fois, je préparerai une "banque de transitions types" (cf. fiche méthode Section 4). »
\end{methode}

\begin{exemplebox}[Exemple de compte rendu écrit (fiche de synthèse)]
\begin{tabularx}{\linewidth}{|l|X|}
\hline
\rowcolor{popblueL}
\textbf{Élément} & \textbf{Contenu} \\
\hline
Sujet & \textit{« Le Vieux nègre et la médaille : une satire de la colonisation ? »} \\
\hline
Problématique & Comment la satire (ironie, registre comique) sert-elle une dénonciation politique radicale ? \\
\hline
Plan & 1. Satire des rites coloniaux (cérémonie) $\rightarrow$ 2. Satire de la société villageoise (complicité) $\rightarrow$ 3. Bascule tragique : le rire laisse place à la révolte. \\
\hline
Citations clés & Ch.6 (ironie), Ch.8 (argent/prime), Ch.9 (latrine). \\
\hline
Auto-évaluation & \textcolor{popgreen}{+}\ Progression thématique respectée. \textcolor{poporange}{-}\ Introduction perfectible (accroche trop banale). \\
\hline
\end{tabularx}
\end{exemplebox}

\courssub{La remédiation collective : corriger les erreurs fréquentes}

Voici les 5 erreurs « tueuses » repérées dans vos copies d'entraînement. La remédiation consiste à les identifier, les comprendre et les éradiquer.

\begin{center}
\begin{tabularx}{\linewidth}{|l|X|X|}
\hline
\rowcolor{popblueL}
\textbf{Erreur type} & \textbf{Diagnostic} & \textbf{Remédiation (Action corrective)} \\
\hline
\textbf{1. Hors-sujet / Déploiement du cours} & Vous récitez l'histoire du roman (résumé) au lieu de répondre à la question. & \textbf{Test du « Pourquoi ? » :} Pour chaque phrase, demandez : « Est-ce que cela répond DIRECTEMENT au sujet ? » Si non $\rightarrow$ Supprimer. \\
\hline
\textbf{2. Citations « orphelines »} & Citation posée, sans introduction, sans analyse, sans lien avec l'argument. & \textbf{Règle des 3 phrases :} 1. Introduction de la citation. 2. La citation (guillemets). 3. Analyse du procédé + lien avec thèse. \\
\hline
\textbf{3. Transitions absentes ou « Mots-liens » vides} & « En premier lieu... En second lieu... Pour finir... » sans lien logique. & \textbf{Transition = Relance :} Reprendre mot-clé § précédent + Annoncer thèse § suivant + Marquer relation logique (cause, concession, opposition). \\
\hline
\textbf{4. Progression thématique rompue} & Paragraphes qui mélangent plusieurs idées, ou ordre illogique (ex: conclusion avant arguments). & \textbf{Plan détaillé obligatoire au brouillon :} 1 idée directrice par §. Vérifier l'enchaînement : Thèse $\rightarrow$ Antithèse $\rightarrow$ Synthèse. \\
\hline
\textbf{5. Langue : syntaxe défaillante / vocabulaire pauvre} & Phrases trop longues (perte du sujet), répétitions, anglicismes, fautes d'accords. & \textbf{Relecture « à l'envers » :} Lire la dernière phrase, puis l'avant-dernière... pour ne voir que la syntaxe. Utiliser \emph{le connecteur logique précis} (cependant, par conséquent, en outre). \\
\hline
\end{tabularx}
\end{center}

\begin{exoresolu}[Correction d'un paragraphe « malade »]
\textbf{Énoncé :} Corriger ce paragraphe (Sujet : \textit{La médaille est-elle un honneur pour Meka ?})
\begin{quote}
\textit{« Meka est content car il a la médaille. "Meka, médaille au cou, se croyait devenu l'égal des Blancs." C'est l'ironie. Après il va voir le chef de poste pour l'argent. Il donne de l'argent. Il est humilié. Donc la médaille n'est pas un honneur. »}
\end{quote}
\textbf{Solution détaillée (paragraphe remédié) :}
\begin{quote}
\textit{« \textbf{La médaille n'est qu'un leurre} qui masque l'humiliation sous les apparences de l'honneur. \textbf{En effet}, lors de la cérémonie, l'ironie narrative dévoile l'abîme entre la perception de Meka et la réalité : "Meka, médaille au cou, se croyait devenu l'égal des Blancs". Le verbe "se croyait" et l'adverbe "devenu" soulignent l'illusion éphémère d'une ascension sociale impossible. \textbf{Or}, cette illusion se brise immédiatement après : pour toucher la prime promise, Meka doit subir le mépris du Chef de Poste et verser des pots-de-vin. \textbf{Ainsi}, loin de conférer du pouvoir, la médaille l'expose à la prédation administrative. \textbf{Par conséquent}, ce symbole honorifique n'est que le marqueur visible de son aliénation. »}
\end{quote}
\textbf{Corrections apportées :} Idée directrice en tête (\textbf{La médaille n'est qu'un leurre}), connecteurs logiques (\textbf{En effet, Or, Ainsi, Par conséquent}), citation intégrée et analysée (ironie, verbes), progression claire (Cérémonie $\rightarrow$ Après-cérémonie $\rightarrow$ Conclusion partielle).
\end{exoresolu}

\courssub{Grille d'auto-évaluation : préparez votre Probatoire}

Utilisez cette grille \textbf{avant de rendre votre copie} (ou en examen, sur brouillon). Cochez honnêtement. Si une case « Non » ou « Partiel » reste cochée : \textbf{ne rendez pas la copie}, corrigez.

\begin{popfigure}
\legende{Grille d'auto-évaluation dissertation — Terminale Technique (à utiliser sur brouillon)}
\begin{center}
\begin{tabularx}{\linewidth}{|>{\bfseries}p{3.5cm}|X|X|X|X|}
\hline
\rowcolor{popdark}
\textcolor{white}{\textbf{Critère / Niveau}} & \textcolor{white}{\textbf{Maîtrise} (\textcolor{popgreen}{✓})} & \textcolor{white}{\textbf{Partiel} (\textcolor{poporange}{~})} & \textcolor{white}{\textbf{Non acquis} (\textcolor{poporange}{✗})} & \textcolor{white}{\textbf{Action corrective immédiate}} \\
\hline
\rowcolor{popblueL}
Pertinence : Sujet analysé, probl. claire, plan annoncé & Sujet décortiqué, probl. formulée, plan dialectique & Problématique floue ou plan binaire & Hors-sujet, pas de probl., plan narratif & Relire sujet, souligner mots-clés, rédiger probl. au brouillon \\
\hline
Cohérence : Idées directrices logiques, thèse/antithèse/synthèse & Enchaînement logique, arguments ordonnés & Quelques sauts logiques, ordre perfectible & Idées dispersées, contradictions, pas de dialectique & Numéroter arguments au brouillon, vérifier ordre thèse/antithèse \\
\hline
\rowcolor{popblueL}
Cohésion : Connecteurs variés, transitions relancées, progression thématique & Connecteurs présents mais répétitifs (et, aussi) & Pas de connecteurs, phrases juxtaposées, pas de transitions & Insérer 1 connecteur par phrase, écrire transition §n/§n+1 \\
\hline
Preuves (Œuvre) : Citations courtes, intégrées, analysées, chap. indiqués & Citations présentes mais non commentées ou trop longues & Pas de citations, ou résumé d'intrigue, citations inventées & Appliquer règle Citer-Analyser-Lier, vérifier banque de citations \\
\hline
\rowcolor{popblueL}
Langue : Syntaxe correcte, vocabulaire précis, 0 faute grave & Quelques maladresses, répétitions, 1-2 fautes & Fautes graves (accords, syntaxe), vocabulaire pauvre & Relire à l'envers, remplacer verbes faibles, vérifier accords SUJET/VERBE \\
\hline
\end{tabularx}
\end{center}
\end{popfigure}

\begin{savaistu}[Ferdinand Oyono (1929–2010) : l'écrivain au service de l'État]
Ferdinand Oyono, auteur du \emph{Vieux nègre et la médaille} (1956), de \emph{Une vie de boy} (1956) et du \emph{Chemin d'Europe} (1960), est une figure majeure des lettres camerounaises. Mais il ne fut pas seulement écrivain : il fut \textbf{ministre} (Fonction publique, puis Affaires étrangères) et \textbf{ambassadeur} du Cameroun (auprès de l'ONU, en France, au Royaume-Uni, en Algérie...). Son œuvre, écrite jeune, anticipe les défis de l'indépendance : la critique de l'assimilation, la récupération de la dignité, le rapport complexe au pouvoir. Étudier Oyono en Terminale technique, c'est lire un homme qui a \textbf{pensé} la colonisation pour mieux \textbf{construire} la nation.
\end{savaistu}

\begin{aretenir}[Synthèse de la section : Votre « Boîte à outils » Oyono pour le Bac]
\begin{enumerate}
    \item \textbf{2 extraits maîtres :} Cérémonie (Ch.6, ironie/illusion) $\leftrightarrow$ Latrine (Ch.9, révolte/lucidité).
    \item \textbf{4 thèmes, 1 citation courte par thème} (Colonisation, Illusion, Dignité, Révolte) — \textbf{apprises par cœur}.
    \item \textbf{3 procédés à citer :} Ironie dramatique, Antithèse, Registre satirique.
    \item \textbf{1 méthode :} Citer $\rightarrow$ Analyser (procédé) $\rightarrow$ Lier (argument).
    \item \textbf{1 grille :} Auto-évaluation (Pertinence, Cohérence, Cohésion, Preuves, Langue) \textbf{avant de rendre}.
\end{enumerate}
\end{aretenir}

\begin{exercice}[Entraînement final : Dissertation intégrée (Sujet type Probatoire)]
\textbf{Sujet :} \textit{« Dans \emph{Le Vieux nègre et la médaille}, la distinction honorifique révèle-t-elle la grandeur de Meka ou sa servitude ? »}
\textbf{Consignes :} 1. Analyser le sujet (mots-clés, problématique, plan dialectique). 2. Rédiger l'introduction complète. 3. Rédiger le \textbf{premier paragraphe du développement} (Thèse : la grandeur illusoire) en exploitant l'ironie du Ch.6 et une citation de votre banque. 4. Écrire la transition vers le §2. 5. Vous auto-évaluer sur la grille (critères 1 à 4).
\end{exercice}

\begin{corrige}[Exercice n°1 — Éléments de correction attendus]
\textbf{1. Analyse :} Mots-clés : « distinction honorifique » (médaille), « révélatrice » (fonction), « grandeur » vs « servitude » (antithèse). Problématique : « La médaille, par-delà son apparat honorifique, ne fonctionne-t-elle pas comme le marqueur ultime de l'aliénation de Meka ? » Plan : I. L'illusion de la grandeur (cérémonie, ironie) $\rightarrow$ II. La réalité de la servitude (extorsion, vol, dépendance) $\rightarrow$ III. Le rejet lucide (geste final, dignité retrouvée).
\textbf{2. Introduction :} Accroche (citation Oyono ou généralité sur décorations coloniales) $\rightarrow$ Présentation œuvre/auteur $\rightarrow$ Sujet amené $\rightarrow$ Problématique $\rightarrow$ Annonce de plan.
\textbf{3. §1 (Thèse) :} Idée directrice : La cérémonie construit une grandeur théâtrale. Citation : « Meka, médaille au cou, se croyait devenu l'égal des Blancs. » Analyse : Ironie dramatique, focalisation interne naïve, verbe « se croyait ». Lien : Cette grandeur est un leurre consenti.
\textbf{4. Transition :} « \textbf{Toutefois}, cette grandeur symbolique s'effondre dès le lendemain face à la réalité administrative et villageoise. » (Reprise « grandeur » + opposition « Toutefois » + annonce « réalité/servitude »).
\textbf{5. Auto-évaluation :} Vérifier cases « Maîtrise » pour Pertinence, Cohérence, Cohésion, Preuves.
\end{corrige}

\coursec{Activités d'intégration : produire une dissertation complète}

La lecture méthodique du \emph{Vieux nègre et la médaille} nous a permis d'observer comment un écrivain construit un texte cohérent : une entrée en matière, une progression organisée, une chute qui donne à réfléchir. La remédiation a corrigé les difficultés repérées dans vos introductions, vos paragraphes et vos conclusions. Il est temps maintenant de \cle{réunir toutes les compétences} acquises dans cette leçon : analyser un sujet, problématiser, rédiger une introduction, construire des paragraphes avec citations et transitions, conclure avec ouverture. L'activité d'intégration est la situation d'examen réelle : en un temps limité, vous devez produire une dissertation complète, structurée et correctement écrite. C'est exactement ce qui vous sera demandé au Probatoire et au Baccalauréat technique.

\courssub{Intégration 1 : produire une introduction et une conclusion}

\textbf{Consigne de travail.} Par binôme, à partir du sujet suivant : \emph{« La médaille récompense-t-elle toujours le mérite ? »}, rédigez en 30 minutes une introduction complète (accroche, présentation du sujet, problématique, annonce du plan) et une conclusion complète (bilan, réponse à la problématique, ouverture). Le développement n'est pas demandé à cette étape.

\begin{methode}[Travail en binôme et correction croisée]
\begin{enumerate}
\item Chaque binôme rédige d'abord au brouillon, puis met au propre introduction et conclusion.
\item Les binômes échangent leurs productions. Chaque binôme corrige la copie d'un autre à l'aide de la grille : l'accroche est-elle présente ? la problématique est-elle une vraie question ? le plan est-il annoncé en deux parties ? la conclusion répond-elle à la problématique ? l'ouverture est-elle liée au sujet sans le répéter ?
\item La correction est justifiée oralement : le correcteur explique \emph{pourquoi} un élément est réussi ou manquant.
\item Chaque binôme réécrit sa production en tenant compte des remarques.
\end{enumerate}
\end{methode}

\begin{exoresolu}[Introduction et conclusion réussies sur le sujet « La médaille récompense-t-elle toujours le mérite ? »]
Rédiger l'introduction et la conclusion de la dissertation sur ce sujet.
\tcblower
\textbf{Introduction.} Dans toutes les sociétés, on décore les hommes : médailles militaires, titres honorifiques, prix scolaires. Ces distinctions semblent couronner le travail et le courage. Pourtant, dans \emph{Le Vieux nègre et la médaille}, Ferdinand Oyono montre un vieil homme décoré non pour son mérite propre, mais pour servir la propagande coloniale. On peut alors se demander : la médaille récompense-t-elle toujours le mérite réel de celui qui la reçoit ? Nous verrons d'abord que la distinction honorifique peut sanctionner de véritables services rendus, puis nous montrerons qu'elle sert souvent des intérêts étrangers au mérite du récipiendaire.

\textbf{Conclusion.} En définitive, la médaille possède une double nature : elle peut honorer légitimement le courage et le dévouement, mais elle peut aussi être un instrument de flatterie, de récompense de la soumission, voire de manipulation, comme le découvre amèrement Meka. La valeur d'une décoration ne tient donc pas au ruban ni au métal, mais à la cause qu'elle sert. Cette réflexion invite à s'interroger plus largement sur toutes les récompenses sociales : les diplômes, les promotions et les honneurs mesurent-ils toujours la valeur réelle des hommes ?
\end{exoresolu}

\begin{attention}[Erreur fréquente : rédiger sans plan]
Beaucoup de candidats se jettent directement dans la rédaction au propre, sans plan détaillé au brouillon. C'est le risque majeur de \cle{hors-sujet} et de \cle{déséquilibre} : on découvre en cours de route que la deuxième partie est vide, que les arguments se répètent, ou que l'on a répondu à côté de la question. Une copie sans plan se reconnaît immédiatement : le correcteur le voit, et le barème le sanctionne. Le plan au brouillon n'est pas du temps perdu : c'est l'assurance de la cohérence.
\end{attention}

\courssub{Intégration 2 : produire un développement complet}

\textbf{Consigne de travail.} Individuellement, rédigez en 1 heure le développement de la dissertation précédente : deux parties (thèse / antithèse), chaque partie comportant deux paragraphes structurés (argument, exemple, citation insérée, mini-conclusion), et une transition rédigée entre les deux parties.

\begin{methode}[Réussir l'épreuve en temps limité : le plan détaillé au brouillon]
Avant toute rédaction au propre, construisez au brouillon un plan détaillé :
\begin{enumerate}
\item Notez la problématique en une phrase interrogative.
\item Formulez l'idée directrice de chaque partie en une phrase complète.
\item Sous chaque partie, notez les deux arguments, l'exemple et la citation prévus (avec la référence : auteur, œuvre).
\item Rédigez au brouillon la phrase de transition.
\item Vérifiez l'équilibre : les deux parties doivent avoir un poids comparable.
\end{enumerate}
Ce plan détaillé prend 15 à 20 minutes ; il vous fait gagner le double pendant la rédaction, car vous n'avez plus qu'à développer, sans chercher vos idées.
\end{methode}

\begin{exemplebox}[Un paragraphe de développement réussi (extrait)]
\emph{Argument :} la médaille peut récompenser un mérite authentique. \emph{Exemple et citation insérée :} dans de nombreuses armées, des soldats sont décorés pour avoir sauvé des vies au péril de la leur ; comme le rappelle l'adage, « à bon ouvrier, bonne récompense ». Au Cameroun, les lauréats du concours d'excellence reçoivent des prix qui couronnent des années de travail scolaire. \emph{Mini-conclusion :} dans ces cas, la distinction joue pleinement son rôle : elle reconnaît publiquement la valeur et encourage l'effort.
\end{exemplebox}

\begin{exemplebox}[Une transition réussie]
« Si la médaille peut ainsi couronner un mérite véritable, l'expérience de Meka, le héros d'Oyono, nous invite à nuancer ce premier constat : la distinction honorifique sert parfois des causes qui n'ont rien à voir avec la valeur de celui qui la reçoit. »
\end{exemplebox}

Remarquez la structure de la transition : la première phrase résume la partie qui se termine, la seconde annonce la partie qui commence. La transition est le pont qui garantit la \cle{progression thématique} du devoir.

\courssub{Gérer les trois heures de l'épreuve}

Au Baccalauréat technique, l'épreuve de dissertation se déroule en un temps limité. La réussite dépend autant de la gestion du temps que des connaissances. Voici la répartition conseillée pour trois heures.

\begin{popfigure}
\begin{tikzpicture}[x=1cm,y=1cm]
% Axe du temps : 3 heures = 12 cm (1 cm = 15 min)
\draw[->,thick,popdark] (0,0) -- (12.6,0) node[right]{\footnotesize temps};
% Graduations toutes les 15 min
\foreach \x/\t in {0/0 h,1/15 min,2/30 min,3/45 min,4/1 h,5/1 h 15,6/1 h 30,7/1 h 45,8/2 h,9/2 h 15,10/2 h 30,11/2 h 45,12/3 h}{
  \draw[popdark] (\x,0) -- (\x,-0.15) node[below]{\scriptsize \t};
}
% Phase 1 : analyse et plan (0 à 3 cm = 45 min)
\fill[popblueL] (0,0.35) rectangle (3,1.15);
\draw[popblue,thick] (0,0.35) rectangle (3,1.15);
\node[align=center] at (1.5,0.75) {\scriptsize \textbf{Analyse du sujet,}\\ \scriptsize \textbf{plan détaillé (45 min)}};
% Phase 2 : rédaction (3 à 11 cm = 2 h)
\fill[poporangeL] (3,0.35) rectangle (11,1.15);
\draw[poporange,thick] (3,0.35) rectangle (11,1.15);
\node at (7,0.75) {\scriptsize \textbf{Rédaction au propre (2 h)}};
% Phase 3 : relecture (11 à 12 cm = 15 min)
\fill[popgreenL] (11,0.35) rectangle (12,1.15);
\draw[popgreen,thick] (11,0.35) rectangle (12,1.15);
\node[align=center] at (11.5,1.55) {\scriptsize \textbf{Relecture}\\ \scriptsize \textbf{(15 min)}};
\draw[->,popgreen] (11.5,1.35) -- (11.5,1.18);
% Détail de la phase 1
\node[align=left,anchor=west] at (0.2,2.6) {\scriptsize \textbf{Phase 1 :} lecture du sujet (5 min), recherche d'idées (10 min),};
\node[align=left,anchor=west] at (0.2,2.25) {\scriptsize problématique et plan détaillé au brouillon (30 min).};
% Détail phase 2
\node[align=left,anchor=west] at (3.2,1.85) {\scriptsize \textbf{Phase 2 :} introduction (25 min), développement (1 h 10), conclusion (25 min).};
\end{tikzpicture}
\legende{Frise chronologique : gestion conseillée des trois heures de l'épreuve de dissertation au Baccalauréat technique. La règle d'or : ne jamais sacrifier la phase 1 (plan) ni la phase 3 (relecture).}
\end{popfigure}

\begin{attention}[Les deux pièges du temps]
\begin{itemize}
\item \textbf{Commencer à rédiger trop tôt :} le candidat qui rédige après 10 minutes de réflexion s'épuise en ratures et en repentirs ; sa copie est déséquilibrée.
\item \textbf{Oublier la relecture :} 15 minutes de relecture permettent de corriger 5 à 10 fautes d'accord ou d'orthographe, ce qui peut représenter un point entier sur la note finale. Arrêtez de rédiger à l'heure prévue, même si la dernière phrase est courte.
\end{itemize}
\end{attention}

\courssub{La relecture finale}

La relecture n'est pas une simple relecture « pour voir » : c'est un contrôle méthodique, point par point.

\begin{methode}[Les cinq contrôles de la relecture]
\begin{enumerate}
\item \textbf{Orthographe lexicale :} relisez mot à mot les termes difficiles (participe passé, mots invariables).
\item \textbf{Accords :} vérifiez chaque sujet-verbe (surtout quand le sujet est éloigné du verbe) et chaque accord du participe passé avec \emph{avoir} (le COD est-il avant ?).
\item \textbf{Ponctuation :} chaque phrase commence par une majuscule et finit par un point ; les deux-points introduisent les citations ; pas de virgule entre le sujet et le verbe.
\item \textbf{Cohérence :} les connecteurs sont-ils variés et corrects ? Les paragraphes sont-ils bien séparés (alinéa) ?
\item \textbf{Lisibilité :} ratures propres (un simple trait), écriture régulière, marge respectée. Une copie illisible fatigue le correcteur et dessert le candidat.
\end{enumerate}
\end{methode}

\begin{exemplebox}[Barème indicatif de l'épreuve de dissertation]
\begin{center}
\begin{tabularx}{\linewidth}{|l|X|c|}\hline
\rowcolor{popblueL} \textbf{Critère} & \textbf{Ce que le correcteur évalue} & \textbf{Poids} \\ \hline
Respect du sujet & Sujet compris, problématique pertinente, pas de hors-sujet & 25 \% \\ \hline
Organisation & Introduction, deux parties équilibrées, transitions, conclusion & 25 \% \\ \hline
Argumentation & Arguments solides, exemples précis, citations insérées & 25 \% \\ \hline
Langue & Orthographe, accords, ponctuation, richesse du vocabulaire & 25 \% \\ \hline
\end{tabularx}
\end{center}
Chaque critère pèse le quart de la note : négliger la langue ou le plan, c'est abandonner 25 \% des points avant même d'être lu.
\end{exemplebox}

\begin{savaistu}[Les meilleures copies ne sont pas les plus longues]
Les correcteurs d'examen le confirment chaque année : les copies les mieux notées ne sont pas les plus longues, mais les mieux structurées et les plus cohérentes. Une dissertation de trois pages avec une problématique claire, deux parties équilibrées, des exemples précis et une langue soignée obtient toujours une meilleure note qu'une copie de huit pages confuse et répétitive. Au Baccalauréat technique, la qualité de l'organisation prime sur la quantité d'encre : mieux vaut quatre pages limpides que huit pages brouillonnes.
\end{savaistu}

\courssub{Intégration 3 : produire à nouveau une introduction et une conclusion}

Après la correction croisée et la remédiation, chaque élève reprend \emph{individuellement} un nouveau sujet, par exemple : \emph{« Faut-il toujours dire la vérité ? »} ou \emph{« Le travail est-il la seule source de la dignité de l'homme ? »}. En 30 minutes, rédigez l'introduction et la conclusion.

\begin{methode}[Mesurer son progrès]
Comparez cette seconde production à la première (Intégration 1) à l'aide de la grille de correction. Le progrès se mesure sur quatre points :
\begin{enumerate}
\item la problématique est formulée plus rapidement et plus précisément ;
\item l'annonce du plan est explicite et équilibrée ;
\item la conclusion répond réellement à la question posée ;
\item le nombre de fautes de langue a diminué.
\end{enumerate}
Notez vos deux productions : l'écart entre les deux est la preuve concrète de votre progression.
\end{methode}

\courssub{Justifier sa démarche : l'oral de l'écrit}

Savoir rédiger ne suffit pas : un bon élève sait aussi \cle{expliquer ses choix}. À l'oral, devant la classe, soyez capable de répondre à trois questions :

\begin{itemize}
\item \textbf{Pourquoi ce plan ?} « J'ai choisi un plan thèse-antithèse parce que le sujet oppose deux points de vue ; la synthèse se trouve dans ma conclusion. »
\item \textbf{Pourquoi cette citation ?} « J'ai cité \emph{Le Vieux nègre et la médaille} dans ma deuxième partie parce que l'exemple de Meka illustre exactement l'argument de la médaille-instrument. »
\item \textbf{Pourquoi cette ouverture ?} « Mon ouverture élargit la réflexion aux récompenses sociales en général, sans quitter le thème de la distinction et du mérite. »
\end{itemize}

Cette justification orale entraîne la métacognition : celui qui sait dire pourquoi il écrit ainsi écrira mieux la fois suivante.

\begin{popfigure}
\begin{tikzpicture}[
  box/.style={draw=popdark,thick,rounded corners=3pt,align=center,minimum width=3.2cm,minimum height=1cm,font=\scriptsize},
  fleche/.style={->,thick,popdark}
]
\node[box,fill=popblueL, align=center] (intro) at (0,0){\textbf{INTRODUCTION}\\ accroche -- sujet\\ problématique -- plan};
\node[box,fill=poporangeL, align=center] (p1) at (-2.4,-2.2){\textbf{PARTIE 1 (thèse)}\\ § 1 : argument + exemple\\ + citation\\ § 2 : argument + exemple};
\node[box,fill=poporangeL, align=center] (p2) at (2.4,-2.2){\textbf{PARTIE 2 (antithèse)}\\ § 1 : argument + exemple\\ + citation\\ § 2 : argument + exemple};
\node[box,fill=popgreenL, align=center] (trans) at (0,-3.6){\textbf{TRANSITION}\\ bilan de la partie 1 + annonce de la partie 2};
\node[box,fill=poppurpleL, align=center] (concl) at (0,-5.4){\textbf{CONCLUSION}\\ bilan -- réponse à la problématique\\ ouverture};
\draw[fleche] (intro) -- (p1);
\draw[fleche] (intro) -- (p2);
\draw[fleche] (p1.south) |- (trans.west);
\draw[fleche] (p2.south) |- (trans.east);
\draw[fleche] (trans) -- (concl);
\end{tikzpicture}
\legende{Architecture complète de la dissertation dialectique : une introduction qui annonce le plan, deux parties équilibrées reliées par une transition, une conclusion qui répond à la problématique et ouvre le débat.}
\end{popfigure}

\begin{exercice}[Dissertation complète en conditions d'examen]
Sujet : \emph{« Les honneurs font-ils le bonheur de l'homme ? »} En vous appuyant sur vos lectures, notamment \emph{Le Vieux nègre et la médaille}, et sur votre expérience, rédigez une dissertation complète en trois heures, en respectant la répartition du temps étudiée. Votre devoir sera évalué selon le barème indicatif (25 \% par critère).
\end{exercice}

\begin{corrige}[Exercice : éléments de correction attendus]
\textbf{Respect du sujet :} la problématique doit interroger le lien entre honneurs et bonheur (ex. : « Les distinctions honorifiques procurent-elles un bonheur véritable ou une satisfaction illusoire ? »). \textbf{Organisation :} partie 1 — les honneurs comblent l'amour-propre et reconnaissent le mérite (exemples : récompenses scolaires, sportives) ; partie 2 — les honneurs peuvent tromper et rendre malheureux (exemple : Meka, humilié le jour même de sa décoration ; la médaille ne l'a pas enrichi ni respecté). \textbf{Argumentation :} au moins deux citations correctement insérées, exemples précis et variés. \textbf{Langue :} accords et ponctuation soignés, connecteurs variés, alinéas visibles. La conclusion doit trancher (les honneurs ne font le bonheur que s'ils couronnent un mérite réel et s'accompagnent de respect) et proposer une ouverture (la dignité vaut-elle plus que la gloire ?).
\end{corrige}

\begin{aretenir}[L'essentiel de l'intégration]
\begin{itemize}
\item Trois heures = 45 min d'analyse et de plan détaillé, 2 h de rédaction, 15 min de relecture.
\item Jamais de rédaction au propre sans plan au brouillon : c'est la première cause de hors-sujet.
\item La dissertation complète = introduction + deux parties équilibrées reliées par une transition + conclusion avec ouverture.
\item La relecture méthodique (orthographe, accords, ponctuation, lisibilité) peut sauver un point entier.
\item Chaque critère du barème pèse 25 \% : soignez le fond ET la forme.
\item Sachez justifier oralement vos choix de plan, de citations et d'ouverture : comprendre sa propre démarche, c'est progresser durablement.
\end{itemize}
\end{aretenir}

\newpage

\coursec{Activité d'intégration}

\courssub{Objectifs de l'activité}
\begin{itemize}
    \item Mobiliser l'ensemble des savoir-faire de la leçon : analyse du sujet, construction de la problématique, rédaction de l'introduction, du développement (avec insertion et commentaire de citations) et de la conclusion.
    \item Appliquer les notions de \cle{progression thématique} pour assurer la cohésion et la cohérence du texte.
    \item Exploiter un corpus documentaire hétérogène (données statistiques, extraits littéraires, témoignage) pour étayer une argumentation.
    \item S'auto-évaluer et évaluer un pair à l'aide d'une grille de critères, puis participer à une remédiation collective.
\end{itemize}

\courssub{Situation}
Le lycée technique de Bafoussam organise, dans le cadre de la semaine de l'excellence technique, un concours inter-établissements de dissertation. Le thème imposé est :
\begin{center}
    \textbf{« L'enseignement technique peut-il résoudre le chômage des jeunes au Cameroun ? »}
\end{center}

Chaque candidat dispose d'un dossier documentaire de quatre pièces à exploiter obligatoirement.

\begin{attention}[Nature du dossier documentaire]
Les données chiffrées du Document 1 et le témoignage du Document 4 sont des \textbf{supports pédagogiques fictifs} construits pour l'entraînement : ils illustrent des ordres de grandeur plausibles mais ne constituent pas des données officielles. Les Documents 2 et 3 sont des \textbf{adaptations libres} inspirées de l'œuvre de Mongo Beti (\emph{Le Vieux nègre et la médaille}, 1956) : à l'examen, on ne cite que des passages réellement lus et vérifiés.
\end{attention}

\textbf{Document 1 : Données chiffrées (support pédagogique inspiré des enquêtes emploi).}
\begin{itemize}
    \item Taux de chômage au sens du BIT (15-35 ans) : \textbf{13,2 \%} (urbain : 18,5 \% ; rural : 8,1 \%).
    \item Taux de sous-emploi visible (heures de travail insuffisantes) : \textbf{22,4 \%} des jeunes occupés.
    \item Part des jeunes diplômés de l'enseignement technique et professionnel parmi les chômeurs : \textbf{28 \%}.
    \item Principale cause déclarée par les employeurs : « inadéquation entre les compétences acquises et les besoins réels de l'entreprise » (62 \% des réponses).
    \item Secteurs porteurs identifiés : BTP, agro-industrie, maintenance industrielle, numérique, énergies renouvelables.
\end{itemize}

\textbf{Document 2 : Extrait adapté de \emph{Le Vieux nègre et la médaille} (Mongo Beti).}
\begin{quote}
« Meka, la médaille au cou, se sentait devenir quelqu'un d'autre. Il n'était plus le vieux cultivateur courbé sous le poids des ans et des impôts ; il était un décoré, un homme que l'Administration honorait. [...] Mais quand le bruit de la fête fut retombé, quand les officiels furent repartis dans leurs grosses voitures, Meka retrouva sa case, son champ, ses dettes. La médaille brillait toujours sur sa poitrine, mais elle ne mangeait pas, elle ne soignait pas la toux de sa femme, elle ne payait pas l'impôt. »
\end{quote}
\emph{Thème exploitable : l'illusion de la reconnaissance symbolique (diplôme, médaille) face à la réalité matérielle et à l'utilité pratique.}

\textbf{Document 3 : Extrait adapté de \emph{Le Vieux nègre et la médaille} (Mongo Beti).}
\begin{quote}
« Le catéchiste avait beau jeu de prêcher le travail, l'ordre, la propreté. Ses ouailles l'écoutaient en hochant la tête, mais leurs ventres restaient vides. [...] L'école du Blanc apprenait à lire, à écrire, à compter ; elle ne savait pas apprendre à vivre sur cette terre, à en tirer le pain quotidien sans en être chassé. »
\end{quote}
\emph{Thème exploitable : la formation déconnectée du réel (purement théorique) face à une formation ancrée dans la production concrète.}

\textbf{Document 4 : Témoignage fictif de Paul K., 29 ans, Brevet de Technicien Supérieur (BTS) Maintenance Industrielle, Douala.}
\begin{quote}
« Après mon BTS en 2018, j'ai cherché du travail pendant 14 mois. Les entreprises voulaient de l'expérience. J'ai accepté un stage non rémunéré dans une usine de transformation de bois à Bonabéri. J'ai appris sur le tas : hydraulique, pneumatique, automatismes. En 2020, avec deux anciens camarades, on a monté "TechnoMaintenance Sarl". Aujourd'hui, on emploie 5 jeunes techniciens en CDI et on sous-traite pour trois usines. La technique, ce n'est pas le diplôme qui la fait, c'est la compétence qu'on prouve chaque jour sur le terrain. Mais sans le BTS, je n'aurais pas eu la base pour comprendre les schémas électriques. »
\end{quote}

\courssub{Consignes}
Le travail se déroule en cinq étapes chronométrées (durée totale indicative : 3 h 30).

\begin{enumerate}
    \item \textbf{Analyse et problématisation (brouillon, 30 min).} Identifiez les mots-clés, les termes du débat, les tensions implicites. Formulez une problématique unique et pertinente. Ébauchez un plan détaillé en deux parties principales (chacune avec deux sous-parties argumentées).
    \item \textbf{Rédaction de l'introduction (45 min).} Rédigez l'introduction complète (accroche, présentation du sujet, définitions utiles, problématique, annonce du plan) en soignant la progression thématique.
    \item \textbf{Production du développement (1 h 15).} Rédigez les deux parties du développement.
        \begin{itemize}
            \item Chaque sous-partie = un paragraphe argumenté (idée directrice, explication, illustration par les documents, commentaire).
            \item Insérez \textbf{au moins deux citations} (une issue des Documents 2 ou 3, une issue du Document 1 ou 4) en les intégrant syntaxiquement et en les commentant.
            \item Rédigez \textbf{une transition explicite} entre la Partie I et la Partie II.
        \end{itemize}
    \item \textbf{Rédaction de la conclusion (30 min).} Faites le bilan synthétique, répondez explicitement à la problématique, proposez une ouverture pertinente (autre dimension, perspective future, comparaison).
    \item \textbf{Correction croisée et remédiation (45 min).} Échangez votre copie avec un voisin. Utilisez la \textbf{grille d'auto-évaluation} fournie (voir Ressources). Notez les forces et les faiblesses. Restitution orale : un binôme présente sa démarche et ses difficultés ; le professeur mène la remédiation collective sur les erreurs récurrentes (hors-sujet, citations « posées » sans commentaire, transitions manquantes, progression thématique rompue).
\end{enumerate}

\courssub{Ressources à mobiliser}
\begin{methode}[Structure type de la dissertation (rappel)]
\begin{enumerate}
    \item \textbf{Introduction} : accroche $\rightarrow$ définitions et mise en contexte $\rightarrow$ sujet posé $\rightarrow$ problématique $\rightarrow$ annonce du plan.
    \item \textbf{Développement} : Partie I (thèse, arguments en faveur) $\rightarrow$ \textbf{transition} $\rightarrow$ Partie II (antithèse, limites, nuances). Chaque paragraphe : idée directrice + argumentation + illustration (citation, chiffre, exemple) + commentaire + mini-bilan ou lien avec la suite.
    \item \textbf{Conclusion} : bilan des deux parties $\rightarrow$ réponse nuancée à la problématique $\rightarrow$ ouverture.
\end{enumerate}
\end{methode}

\begin{propriete}[Progression thématique : cohésion et cohérence]
\begin{itemize}
    \item \textbf{Progression linéaire (en chaîne)} : le thème de la phrase $n$ devient le point de départ de la phrase $n+1$ (ex. : \emph{L'enseignement technique forme des techniciens. Ces techniciens maîtrisent des gestes professionnels. Ces gestes répondent aux besoins des entreprises.}).
    \item \textbf{Progression à thème constant} : le même sujet grammatical (ou sa reprise pronominale) est maintenu sur plusieurs phrases (ex. : \emph{Le technicien... Il... Il...}).
    \item \textbf{Progression à thème dérivé} : déclinaison d'un hyper-thème en sous-thèmes (ex. : \emph{L'État... Le ministère... Les lycées techniques... Les ateliers...}).
    \item \textbf{Cohérence} : connecteurs logiques bien choisis (\emph{donc, toutefois, en effet, par conséquent}), champs lexicaux homogènes, concordance des temps respectée.
\end{itemize}
\end{propriete}

\begin{aretenir}[Insertion et commentaire de citation]
Une citation ne parle jamais d'elle-même. Trois étapes obligatoires :
\begin{enumerate}
    \item \textbf{Introduire} : « Comme l'illustre Mongo Beti... », « Selon les données du Document 1... », « Paul K. témoigne : ... ».
    \item \textbf{Citer} : guillemets, exactitude des mots, référence précise (Doc n°).
    \item \textbf{Commenter} : expliquer le sens, relier la citation à l'argument, montrer pourquoi elle prouve le point défendu (verbes utiles : \emph{révèle, montre, confirme, contredit, nuance}).
\end{enumerate}
\end{aretenir}

\begin{exemplebox}[Grille d'auto-évaluation et d'évaluation par les pairs]
\begin{center}
\begin{tabularx}{\linewidth}{|l|X|c|}\hline
\rowcolor{popblueL}
\textbf{Critère} & \textbf{Indicateurs de réussite} & \textbf{Oui / Partiel / Non} \\ \hline
Analyse du sujet / Problématique & Mots-clés définis, tension réelle dégagée, question unique & \\ \hline
Introduction & Accroche pertinente, définitions claires, plan annoncé distinctement & \\ \hline
Développement : argumentation & Idées directrices claires, arguments logiques, exemples variés (documents + savoirs personnels) & \\ \hline
Développement : citations & Au moins 2 citations insérées \textbf{et} commentées (pas « posées ») & \\ \hline
Développement : transitions & Transition I/II rédigée (rappel + annonce), liens internes fluides & \\ \hline
Progression thématique & Enchaînement des phrases et des paragraphes fluide (linéaire, à thème constant ou dérivé) & \\ \hline
Conclusion & Bilan synthétique, réponse explicite, ouverture pertinente & \\ \hline
Langue / orthographe & Syntaxe correcte, vocabulaire précis, accords, ponctuation & \\ \hline
\end{tabularx}
\end{center}
\end{exemplebox}

\courssub{Démarche et corrigé commenté}
\begin{exoresolu}[Production d'une dissertation modèle : « L'enseignement technique peut-il résoudre le chômage des jeunes au Cameroun ? »]
\textbf{ÉTAPE 1 : Analyse du sujet et problématique (brouillon)}

\begin{itemize}
    \item \textbf{Mots-clés :} « enseignement technique » (formation professionnelle, compétences pratiques, CAP, Probatoire et Baccalauréat technique, BTS, écoles d'ingénieurs) ; « résoudre » (éradiquer, réduire significativement, être la solution principale) ; « chômage des jeunes » (insertion, sous-emploi, inadéquation formation-emploi) ; « Cameroun » (contexte : poids du secteur informel, PME fragiles, forte croissance démographique).
    \item \textbf{Tensions et débat :} l'enseignement technique est présenté comme \emph{la} solution dans les discours officiels $\leftrightarrow$ réalités : chômage des diplômés techniques (Doc 1 : 28 \%), inadéquation des compétences (62 \%), illusion du diplôme (Doc 2), nécessité de l'expérience de terrain (Doc 4).
    \item \textbf{Problématique choisie :} \emph{Si l'enseignement technique constitue un levier indispensable pour l'insertion professionnelle par l'acquisition de compétences opérationnelles, ne risque-t-il pas de rester inefficace face au chômage structurel s'il n'est pas adossé à une politique globale de création d'emplois et d'adaptation continue aux besoins réels du tissu productif ?}
    \item \textbf{Plan :}
    \begin{itemize}
        \item \textbf{I. L'enseignement technique : un levier majeur pour l'employabilité et l'autonomisation des jeunes.}
        \begin{itemize}
            \item A. L'adéquation formation-emploi : une réponse à la demande de compétences techniques (Doc 1, Doc 4).
            \item B. La formation comme tremplin vers l'entrepreneuriat et la création d'emplois (Doc 4).
        \end{itemize}
        \item \textbf{II. Les limites structurelles qui entravent son efficacité face au chômage de masse.}
        \begin{itemize}
            \item A. Le risque d'une formation déconnectée du réel : l'illusion du diplôme (Doc 2, Doc 3, Doc 1).
            \item B. L'insuffisance de l'offre économique et de l'accompagnement post-formation (Doc 1, Doc 4).
        \end{itemize}
    \end{itemize}
\end{itemize}

\textbf{ÉTAPE 2 : Introduction (rédaction)}

Dans ses rapports sur l'emploi, l'Organisation internationale du Travail répète que les compétences sont le passeport pour l'emploi ; pourtant, l'Afrique subsaharienne demeure confrontée à des taux de chômage et de sous-emploi des jeunes parmi les plus élevés du monde. Au Cameroun, la question de l'insertion professionnelle des 15-35 ans, qui représentent une part considérable de la population, est devenue un enjeu national. L'enseignement technique, qu'il soit secondaire (CAP, Probatoire et Baccalauréat technique) ou supérieur (BTS, écoles d'ingénieurs), est présenté par les pouvoirs publics comme la \cle{voie royale} pour résorber ce fléau grâce à l'acquisition de « compétences immédiatement opérationnelles ». \textbf{Pourtant, les données du Document 1 révèlent que 28 \% des chômeurs sont issus de cette filière et que 62 \% des employeurs dénoncent une inadéquation des compétences.} Faut-il y voir l'échec du modèle ou les symptômes d'un mal plus profond ? \textbf{Si l'enseignement technique constitue un levier indispensable pour l'insertion professionnelle par l'acquisition de compétences opérationnelles, ne risque-t-il pas de rester inefficace face au chômage structurel s'il n'est pas adossé à une politique globale de création d'emplois et d'adaptation continue aux besoins réels du tissu productif ?} Nous montrerons d'abord que l'enseignement technique offre une réponse concrète à la demande de qualifications (I), avant d'analyser les obstacles structurels — pédagogiques, économiques et institutionnels — qui limitent sa portée (II).

\textbf{ÉTAPE 3 : Développement (rédaction avec citations et transition)}

\textbf{PARTIE I : L'ENSEIGNEMENT TECHNIQUE, UN LEVIER MAJEUR POUR L'EMPLOYABILITÉ ET L'AUTONOMISATION}

\textbf{A. L'adéquation formation-emploi : une réponse à la demande de compétences techniques.}
L'atout premier de la voie technique réside dans sa pédagogie de l'action : alternance entre théorie et pratique, stages en entreprise, travaux sur plateaux techniques. Contrairement à l'enseignement général, parfois accusé de former des diplômés sans métier précis, elle vise l'\cle{opérationnalité} immédiate. Les secteurs porteurs identifiés dans le Document 1 — BTP, agro-industrie, maintenance, numérique — requièrent des techniciens capables de lire un plan, de programmer un automate, d'entretenir une chaîne de froid. \textbf{Comme le montrent les données du Document 1, les « secteurs porteurs identifiés » correspondent exactement aux spécialités enseignées dans les lycées techniques : BTP, maintenance industrielle, agro-alimentaire.} Cette corrélation n'est pas fortuite ; elle prouve que lorsque la carte des formations est alignée sur la carte économique du pays, le technicien trouve sa place. La progression thématique adoptée ici (secteurs $\rightarrow$ spécialités $\rightarrow$ technicien $\rightarrow$ place) assure la cohérence de l'argumentation.

\textbf{B. La formation comme tremplin vers l'entrepreneuriat et la création d'emplois.}
Au-delà du salariat, l'enseignement technique forge l'esprit d'initiative. La maîtrise d'un métier technique permet de créer sa propre structure de services, transformant le demandeur d'emploi en pourvoyeur d'emplois. \textbf{Paul K., dans son témoignage (Document 4), l'illustre parfaitement : « En 2020, avec deux anciens camarades, on a monté "TechnoMaintenance Sarl". Aujourd'hui, on emploie 5 jeunes techniciens en CDI. »} Cette citation n'est pas anecdotique ; elle démontre l'effet \cle{multiplicateur} de la formation technique : un diplômé formé, c'est potentiellement plusieurs emplois créés. Le propos de Paul K. — « ce n'est pas le diplôme qui la fait, c'est la compétence qu'on prouve chaque jour » — apporte toutefois une nuance : le diplôme est la \emph{condition nécessaire} (la base pour « comprendre les schémas électriques »), mais l'expérience de terrain en est la \emph{condition de réussite} entrepreneuriale.

\textbf{TRANSITION}
\emph{Toutefois, si le technicien formé trouve plus aisément sa place que le généraliste, les chiffres du Document 1 (28 \% de chômeurs issus de la filière technique, 22,4 \% de sous-emploi) interdisent tout triomphalisme. Le diplôme technique ne saurait être une médaille magique qui dispenserait de lutter contre les causes structurelles du chômage. C'est ce que suggère, avec une ironie lucide, Mongo Beti dans \emph{Le Vieux nègre et la médaille}.}

\textbf{PARTIE II : LES LIMITES STRUCTURELLES QUI ENTRAVENT SON EFFICACITÉ FACE AU CHÔMAGE DE MASSE}

\textbf{A. Le risque d'une formation déconnectée du réel : l'illusion du diplôme.}
Le danger majeur est la \cle{fétichisation du diplôme} : croire que le titre suffit, comme Meka croit que la médaille change sa condition. \textbf{Dans le Document 2, le narrateur observe : « La médaille brillait toujours sur sa poitrine, mais elle ne mangeait pas, elle ne soignait pas la toux de sa femme, elle ne payait pas l'impôt. »} Cette image puissante s'applique au diplôme technique obtenu dans des ateliers vétustes, avec des programmes obsolètes et sans stage réel : il brille sur le CV mais ne « nourrit » pas son titulaire. Le Document 3 renforce cette critique : « L'école du Blanc apprenait à lire, à écrire, à compter ; elle ne savait pas apprendre à vivre sur cette terre, à en tirer le pain quotidien. » Transposée à notre contexte, cette remarque vise toute formation purement scolaire, déconnectée du tissu productif local : elle produit des diplômés « inemployables », incapables de tirer « le pain quotidien » de leur compétence. L'inadéquation signalée par 62 \% des employeurs (Doc 1) trouve ici sa racine pédagogique.

\textbf{B. L'insuffisance de l'offre économique et de l'accompagnement post-formation.}
Même compétent, le jeune technicien se heurte à un marché de l'emploi étroit. L'économie camerounaise ne crée pas assez d'emplois formels pour absorber le flux annuel des sortants du système éducatif, et le secteur informel, qui occupe l'essentiel de la main-d'œuvre, n'intègre pas toujours les qualifications techniques certifiées. \textbf{Le parcours de Paul K. (Doc 4) est révélateur : « J'ai cherché du travail pendant 14 mois. Les entreprises voulaient de l'expérience. J'ai accepté un stage non rémunéré... »} Cette citation montre que la \cle{barrière à l'entrée} n'est pas d'abord la compétence, mais l'expérience exigée et la capacité d'absorption du marché. Sans politique volontariste d'incitation à l'embauche (allègements de charges, accès à la commande publique) et sans mécanismes de financement de l'entrepreneuriat des jeunes, l'enseignement technique risque de former pour le chômage : il ne « résout » pas le chômage, il ne fait que \emph{qualifier} les chômeurs.

\textbf{ÉTAPE 4 : Conclusion (rédaction)}

\textbf{Bilan.} Au terme de cette analyse, l'enseignement technique apparaît comme une \cle{condition nécessaire mais non suffisante} pour résoudre le chômage des jeunes au Cameroun. Nécessaire, car il fournit le capital humain technique sans lequel aucune industrialisation, aucune modernisation de l'agriculture, aucune maintenance des infrastructures n'est possible (Partie I). Non suffisante, car il ne commande ni la croissance économique, ni la politique de l'emploi, ni la mise à jour continue des programmes de formation (Partie II).

\textbf{Réponse à la problématique.} L'enseignement technique ne « résout » pas le chômage à lui seul ; il en est un \emph{instrument puissant} dont l'efficacité dépend d'une chaîne complète : \emph{orientation $\rightarrow$ formation de qualité (stages, plateaux techniques) $\rightarrow$ certification reconnue $\rightarrow$ insertion accompagnée (stage rémunéré, incubation) $\rightarrow$ demande réelle d'emplois (politique industrielle, soutien aux PME)}. Briser un seul maillon — stages non rémunérés généralisés, programmes obsolètes — rend l'outil inopérant.

\textbf{Ouverture.} Le développement de la \cle{formation par alternance} et la modernisation des lycées techniques vont dans le bon sens. Mais l'enjeu ultime dépasse l'école : c'est la \textbf{transformation structurelle de l'économie camerounaise} — industrialisation, formalisation progressive du secteur informel — qui décidera si le technicien camerounais de demain sera un Meka décoré mais démuni, ou un Paul K. qui « prouve sa compétence chaque jour sur le terrain » en créant une richesse partagée.
\end{exoresolu}

\courssub{Bilan de l'activité}
Cette activité d'intégration a permis de mettre en évidence plusieurs acquis majeurs et points de vigilance pour l'examen du Probatoire et du Baccalauréat technique :

\begin{itemize}
    \item \textbf{Maîtrise de la méthode de la dissertation :} les élèves qui ont suivi l'enchaînement « analyse $\rightarrow$ problématique $\rightarrow$ plan $\rightarrow$ rédaction » ont produit des copies structurées et lisibles, avec des transitions explicites. Ceux qui ont commencé par rédiger l'introduction sans brouillon d'analyse ont souvent abouti à des problématiques floues ou à des plans déséquilibrés.
    \item \textbf{Progression thématique :} l'exercice a montré la différence entre un texte « qui tient la route » (progression linéaire ou à thème constant bien gérée) et un texte « décousu » (sauts de thème brutaux, répétition de tournures comme « il y a », « c'est »). La révision par les pairs s'est révélée très efficace sur ce point : lire le texte d'un autre oblige à repérer les ruptures de cohérence.
    \item \textbf{Insertion des citations :} c'est le point le plus perfectible. Beaucoup d'élèves ont « posé » la citation (alinéa, guillemets, point) sans la \emph{commenter}. Le corrigé modèle montre qu'une citation doit être \textbf{préparée} (phrase introductrice), \textbf{intégrée} (syntaxe) et \textbf{exploitée} (verbe de commentaire + lien avec l'argument). La grille d'évaluation sanctionne cet écart.
    \item \textbf{Exploitation du corpus :} l'obligation d'utiliser les quatre documents a forcé la confrontation des sources (statistiques, littérature, témoignage). Les meilleures copies ont croisé les Documents 1 et 4 (chiffres et vécu), puis les Documents 2 et 3 avec le Document 1 (illusion littéraire et réalité statistique).
    \item \textbf{Remédiation collective :} la restitution orale a permis de nommer les erreurs types : « introduction sans accroche », « plan en liste de courses », « conclusion qui répète l'introduction », « ouverture hors sujet ». Le professeur a pu réexpliquer \emph{in situ} la différence entre une \emph{transition} (rappel de ce qui précède + annonce de ce qui suit) et un simple connecteur (« de plus », « par ailleurs »).
\end{itemize}

\begin{attention}[Pièges à éviter pour l'examen]
\begin{itemize}
    \item Ne pas confondre \textbf{sujet} et \textbf{problématique}. Le sujet est l'énoncé donné ; la problématique est la question tensionnée que \emph{vous} posez à partir de lui.
    \item Ne pas écrire « Je vais dire que... » dans l'annonce du plan. Préférer : « Nous analyserons d'abord... puis nous montrerons... ».
    \item Une citation \textbf{non commentée} ne rapporte aucun point en argumentation : elle alourdit le texte sans le faire avancer.
    \item La \textbf{transition} est une phrase (ou deux) à part entière, pas un simple mot placé au début d'un paragraphe.
    \item L'\textbf{ouverture} ne doit pas être une question creuse (« Et si on faisait mieux ? ») mais une perspective réelle : économique, sociale, historique, comparative ou prospective.
    \item Ne citer que des passages réellement lus : une citation inventée ou déformée est sanctionnée. En cas de doute, paraphrasez en nommant l'œuvre et l'auteur.
\end{itemize}
\end{attention}

\newpage

\coursec{Méthodes et exercices résolus}

\courssub{Méthode 1 : Analyser le sujet et construire la problématique}

\begin{methode}[Analyser un sujet de dissertation]
\begin{enumerate}
\item \textbf{Lire le sujet plusieurs fois} et souligner les \cle{mots-clés} (thème, notion, auteur, œuvre).
\item \textbf{Définir chaque mot-clé} : donner son sens précis, ses synonymes, son champ lexical.
\item \textbf{Dégager l'enjeu} : pourquoi cette question est-elle posée ? Quel débat sous-entend-elle ?
\item \textbf{Transformer le sujet en question} : c'est la \cle{problématique}, qui doit appeler une réponse argumentée (et non un simple oui ou non).
\item \textbf{Mobiliser ses connaissances} : œuvres lues en classe, exemples de la vie courante, citations mémorisées.
\item \textbf{Choisir un plan} : thématique (idées), dialectique (thèse / antithèse / synthèse) ou analytique (constat / causes / conséquences).
\end{enumerate}
\textbf{Réflexe :} la problématique est une question qui contient les mots-clés du sujet et annonce le plan.
\end{methode}

\begin{definition}[Problématique]
La \cle{problématique} est la question directrice de la dissertation : elle naît de l'analyse des mots-clés du sujet, formule le débat à trancher et commande l'ensemble du plan. Elle se pose généralement sous forme interrogative, à la fin de la présentation du sujet dans l'introduction.
\end{definition}

\begin{exoresolu}[Analyse du sujet et problématique]
\textbf{Énoncé (type Bac) :} Sujet : \emph{Le roman africain ne sert qu'à dénoncer les injustices de la colonisation.} Partagez-vous ce point de vue ? Vous appuierez votre réflexion sur les œuvres étudiées en classe.
\par 1) Relevez les mots-clés. 2) Formulez une problématique. 3) Proposez un plan adapté.
\tcblower
\textbf{Solution détaillée :}
\par \textbf{1) Mots-clés :} \emph{roman africain} (genre littéraire produit par des écrivains d'Afrique) ; \emph{ne sert qu'à} (restriction, réduction à une seule fonction) ; \emph{dénoncer} (fonction critique, engagée) ; \emph{injustices de la colonisation} (contexte historique).
\par \textbf{2) Problématique :} Le roman africain se réduit-il à la dénonciation des injustices coloniales, ou remplit-il d'autres fonctions ?
\par \textbf{3) Plan dialectique :}
\begin{itemize}
\item \textbf{Thèse :} le roman africain est bien une littérature de dénonciation (ex. \emph{Le Vieux nègre et la médaille} de Ferdinand Oyono : satire de l'administration coloniale et de la mission religieuse).
\item \textbf{Antithèse :} mais il a d'autres fonctions : valorisation de la culture et des traditions africaines, peinture de la société, divertissement, éducation.
\item \textbf{Synthèse :} le roman africain est à la fois une arme de combat et un miroir de la société africaine.
\end{itemize}
\textbf{Conclusion :} le plan dialectique convient car le sujet contient une restriction (\emph{ne sert qu'à}) qu'il faut discuter.
\end{exoresolu}

\courssub{Méthode 2 : Assurer la progression thématique (cohésion et cohérence)}

\begin{definition}[Cohérence et cohésion]
La \cle{cohérence} concerne le fond : toutes les idées du texte répondent à la problématique et s'organisent logiquement. La \cle{cohésion} concerne la forme : les phrases sont reliées entre elles par des procédés linguistiques (reprises, pronoms, connecteurs). La \cle{progression thématique} désigne la manière dont le thème (ce dont on parle) est repris et enrichi de phrase en phrase par le propos (l'information nouvelle).
\end{definition}

\begin{methode}[Garantir cohésion et cohérence]
\begin{enumerate}
\item \textbf{Cohérence} : toutes les idées doivent répondre à la problématique ; chaque paragraphe développe une seule idée.
\item \textbf{Cohésion} : relier les phrases entre elles par des \cle{reprises nominales} (synonymes, hyperonymes : \emph{le colon} $\to$ \emph{l'administrateur}), des \cle{pronoms} (il, celui-ci) et des \cle{connecteurs logiques}.
\item Utiliser les connecteurs adaptés : addition (\emph{d'abord, ensuite, de plus}), opposition (\emph{cependant, néanmoins, en revanche}), cause (\emph{car, parce que, c'est pourquoi}), conséquence (\emph{donc, ainsi, par conséquent}), illustration (\emph{ainsi, par exemple}).
\item Vérifier la \cle{progression thématique} : le thème (ce dont on parle) doit être repris d'une phrase à l'autre ; le propos (ce qu'on en dit) apporte l'information nouvelle.
\item Relire en vérifiant qu'aucune phrase n'est « orpheline » (sans lien avec la précédente).
\end{enumerate}
\end{methode}

\begin{exoresolu}[Améliorer la cohésion d'un paragraphe]
\textbf{Énoncé :} Voici un paragraphe mal construit. Réécrivez-le en assurant sa cohésion : « Meka reçoit une médaille. La médaille est donnée par l'administration coloniale. Les villageois admirent Meka. Meka est humilié par les Blancs. La médaille ne protège pas. »
\tcblower
\textbf{Solution détaillée :}
\par On identifie le thème récurrent : Meka et sa médaille. On relie ensuite les phrases par des reprises et des connecteurs :
\par « Meka, le héros du \emph{Vieux nègre et la médaille}, reçoit une médaille des mains de l'administration coloniale. Cette récompense suscite d'abord l'admiration des villageois, qui voient en lui un homme honoré. Cependant, le vieillard est bientôt humilié par les mêmes Blancs qui l'ont décoré. Ainsi, la médaille ne le protège en rien : elle apparaît comme un instrument de domination. »
\par \textbf{Justification :} reprise nominale (\emph{le héros}, \emph{le vieillard}), pronom (\emph{lui}), connecteurs (\emph{d'abord, cependant, ainsi}) : le paragraphe est désormais cohérent et cohésif.
\end{exoresolu}

\courssub{Méthode 3 : Rédiger l'introduction}

\begin{methode}[Les 4 étapes de l'introduction]
\begin{enumerate}
\item \textbf{L'amorce} (ou accroche) : une phrase générale qui situe le sujet dans un contexte littéraire, historique ou social. Jamais de définition de dictionnaire.
\item \textbf{La présentation du sujet} : reformuler le sujet avec ses propres mots, en définissant les termes importants.
\item \textbf{La problématique} : poser clairement la question directrice.
\item \textbf{L'annonce du plan} : présenter les deux ou trois grandes parties du développement (sans les numéroter).
\end{enumerate}
\textbf{Longueur :} 8 à 12 lignes. \textbf{Réflexe :} l'introduction se rédige au brouillon en dernier, quand le plan est fixé.
\end{methode}

\begin{exoresolu}[Rédiger une introduction complète]
\textbf{Énoncé :} Rédigez l'introduction de la dissertation sur le sujet : \emph{La littérature africaine doit-elle être engagée ?}
\tcblower
\textbf{Solution détaillée :}
\par \textbf{Amorce :} Depuis les indépendances, la littérature africaine occupe une place croissante sur la scène mondiale.
\par \textbf{Présentation du sujet :} Née dans un contexte de domination coloniale, elle a souvent servi d'arme pour dénoncer l'oppression et revendiquer la dignité des peuples. Mais être « engagée », c'est-à-dire prendre parti dans les débats de la société, est-ce là sa seule mission ?
\par \textbf{Problématique :} La littérature africaine doit-elle nécessairement être engagée, ou peut-elle poursuivre d'autres finalités ?
\par \textbf{Annonce du plan :} Nous verrons d'abord que l'engagement est au cœur de la littérature africaine ; nous montrerons ensuite qu'elle remplit aussi des fonctions esthétiques et culturelles ; enfin, nous dégagerons une position nuancée.
\par \textbf{Conclusion :} l'introduction contient les quatre étapes obligatoires, dans l'ordre, et annonce un plan dialectique.
\end{exoresolu}

\courssub{Méthode 4 : Rédiger un paragraphe avec citation et transition}

\begin{methode}[Construire un paragraphe argumentatif]
\begin{enumerate}
\item \textbf{Phrase d'idée} : énoncer l'argument du paragraphe en une phrase claire.
\item \textbf{Explication} : développer l'idée en deux ou trois phrases.
\item \textbf{Exemple et citation} : illustrer par un exemple précis tiré d'une œuvre ; la \cle{citation} est introduite (selon + auteur, verbe de parole), entre guillemets, puis \emph{commentée} (jamais laissée seule).
\item \textbf{Phrase de bilan} : conclure le paragraphe en rattachant l'idée à la problématique.
\item \textbf{Transition} : à la fin de chaque grande partie, faire le bilan de ce qui vient d'être dit et annoncer la partie suivante (\emph{Si..., il n'en reste pas moins que...}).
\end{enumerate}
\textbf{Attention :} une citation courte, exacte, entre guillemets, avec le nom de l'auteur et le titre de l'œuvre en italique.
\end{methode}

\begin{exoresolu}[Paragraphe avec citation et transition]
\textbf{Énoncé :} Rédigez un paragraphe montrant que \emph{Le Vieux nègre et la médaille} de Ferdinand Oyono dénonce l'hypocrisie coloniale, puis la transition vers la partie suivante (les autres fonctions du roman africain).
\tcblower
\textbf{Solution détaillée :}
\par \textbf{Paragraphe :} Le roman de Ferdinand Oyono illustre d'abord la fonction de dénonciation de la littérature africaine. En effet, l'auteur y ridiculise l'administration coloniale et le clergé, qui prétendent honorer les Africains tout en les méprisant. Meka, décoré pour avoir « donné » ses terres et ses fils à la France, est arrêté et battu dès la nuit tombée. Cette médaille, symbole de reconnaissance, se révèle donc un leurre : elle masque la violence réelle de la colonisation. Par l'ironie et la satire, Oyono démasque l'hypocrisie du discours colonial.
\par \textbf{Transition :} Si le roman africain est ainsi une arme de dénonciation, il ne saurait se réduire à cette seule fonction : il est aussi un lieu de mémoire et de célébration de la culture, comme nous allons le montrer.
\par \textbf{Conclusion :} le paragraphe respecte la structure idée / explication / exemple commenté / bilan, et la transition relie logiquement les deux parties.
\end{exoresolu}

\courssub{Méthode 5 : Rédiger la conclusion}

\begin{methode}[Les 3 étapes de la conclusion]
\begin{enumerate}
\item \textbf{Le bilan} : répondre explicitement à la problématique en résumant les grandes étapes du raisonnement (sans répéter mot à mot).
\item \textbf{La prise de position} : affirmer clairement son point de vue personnel et argumenté.
\item \textbf{L'ouverture} : élargir la réflexion vers une question voisine, un autre genre, un enjeu contemporain (sans nouveau développement).
\end{enumerate}
\textbf{Interdits :} aucune idée nouvelle, aucune citation nouvelle, pas de « je » trop appuyé (préférer \emph{nous} ou les tournures impersonnelles). Longueur : 5 à 8 lignes.
\end{methode}

\begin{exoresolu}[Rédiger une conclusion]
\textbf{Énoncé :} Rédigez la conclusion de la dissertation : \emph{La littérature africaine doit-elle être engagée ?}
\tcblower
\textbf{Solution détaillée :}
\par \textbf{Bilan :} En définitive, la littérature africaine est née de l'engagement : des œuvres comme \emph{Le Vieux nègre et la médaille} ont dénoncé les injustices coloniales et redonné voix aux opprimés. Mais nous avons vu qu'elle célèbre aussi les traditions, peint la société et procure du plaisir esthétique.
\par \textbf{Prise de position :} L'engagement apparaît donc non comme une obligation, mais comme une dimension essentielle parmi d'autres : un écrivain africain reste libre de choisir son sujet, sans pour autant renier l'histoire de son continent.
\par \textbf{Ouverture :} Cette question rejoint celle, plus large, du rôle de l'art dans les sociétés africaines contemporaines, confrontées à de nouveaux défis comme la mondialisation.
\par \textbf{Conclusion :} la conclusion répond à la problématique, prend position et s'ouvre sans introduire d'argument nouveau.
\end{exoresolu}

\courssub{Méthode 6 : Lecture méthodique du \emph{Vieux nègre et la médaille} (n\textsuperscript{o} 2 et 3)}

\begin{aretenir}[Repères sur l'œuvre]
\emph{Le Vieux nègre et la médaille} (1956) est un roman de l'écrivain camerounais \cle{Ferdinand Oyono}. Le héros, \cle{Meka}, vieux villageois de Doum, reçoit une médaille de l'administration coloniale pour avoir « donné » ses terres à la mission catholique et ses deux fils, morts pour la France pendant la guerre. Arrêté puis battu la nuit même de la cérémonie pour s'être égaré dans le quartier des Blancs, il comprend que cette récompense est un leurre : la médaille ne change rien au mépris que les colons portent aux Africains.
\end{aretenir}

\begin{methode}[Exploiter l'œuvre dans une dissertation]
\begin{enumerate}
\item \textbf{Dégager les thèmes majeurs} : l'hypocrisie coloniale, l'illusion de la récompense, le choc des cultures, la solidarité villageoise, la révolte naissante.
\item \textbf{Repérer les procédés d'écriture} : l'\cle{ironie} et la \cle{satire} (ridiculiser pour dénoncer), le récit à la troisième personne centré sur Meka, le comique de situation.
\item \textbf{Sélectionner des exemples précis} : la cérémonie de remise de la médaille, l'arrestation nocturne de Meka, les réactions des villageois de Doum.
\item \textbf{Rattacher chaque exemple à un argument} : un exemple ne vaut que s'il illustre une idée au service de la problématique.
\end{enumerate}
\end{methode}

\begin{exemplebox}[Exemples exploitables tirés de l'œuvre]
\begin{itemize}
\item \textbf{Dénonciation de l'hypocrisie :} le commandant et le père Vandermayer célèbrent Meka le jour, et le font brutaliser la nuit : l'écart entre le discours et les actes révèle le mensonge colonial.
\item \textbf{Peinture de la société :} le roman décrit la vie du village de Doum, les fêtes, les croyances et la solidarité des anciens : il est aussi un témoignage sur la société camerounaise traditionnelle.
\item \textbf{Prise de conscience :} la médaille perdue puis retrouvée symbolise le cheminement de Meka, de la fierté naïve à la lucidité.
\end{itemize}
\end{exemplebox}

\courssub{Méthode 7 : Produire une dissertation complète (activité d'intégration)}

\begin{methode}[Organiser son travail le jour de l'épreuve]
\begin{enumerate}
\item \textbf{Gérer le temps} (environ 3 h) : 30 min d'analyse du sujet et de plan au brouillon, 2 h de rédaction, 30 min de relecture.
\item \textbf{Au brouillon :} mots-clés, problématique, plan détaillé (idées + exemples + citations), rédaction de l'introduction et de la conclusion.
\item \textbf{Sur la copie :} introduction, développement (un paragraphe = une idée + un exemple ; transitions visibles entre les parties), conclusion.
\item \textbf{Matérialiser le plan :} sauter une ligne entre introduction, développement et conclusion ; alinéa à chaque paragraphe.
\item \textbf{Relire} : orthographe, accords, ponctuation, cohésion (connecteurs), exactitude des citations.
\end{enumerate}
\end{methode}

\begin{exoresolu}[Dissertation complète : squelette rédigé]
\textbf{Énoncé (type Bac) :} \emph{Le roman, miroir de la société : qu'en pensez-vous ?} Vous rédigerez une dissertation complète en vous appuyant sur les œuvres étudiées.
\tcblower
\textbf{Solution détaillée (plan et éléments rédigés) :}
\par \textbf{Problématique :} Dans quelle mesure le roman reflète-t-il la société de son époque ?
\par \textbf{Plan thématique en trois parties :}
\begin{itemize}
\item \textbf{I. Le roman peint la réalité sociale :} \emph{Le Vieux nègre et la médaille} montre la société coloniale camerounaise (village, administration, mission) ; le réalisme du genre romanesque.
\item \textbf{II. Le roman juge et transforme cette réalité :} la satire d'Oyono déforme pour mieux dénoncer ; le romancier n'est pas un photographe mais un critique.
\item \textbf{III. Le roman dépasse aussi le réel :} imagination, idéal, rêve ; le lecteur trouve dans le roman un moyen d'évasion et de réflexion universelle.
\end{itemize}
\par \textbf{Introduction (rédigée) :} Le roman est souvent présenté comme un « miroir que l'on promène le long d'un chemin ». Cette image suggère que la fiction reflète fidèlement la société. Or le romancier sélectionne, transforme et juge le réel. Dans quelle mesure le roman est-il alors le miroir de la société ? Nous montrerons qu'il peint la réalité sociale, qu'il la juge en la transformant, et qu'il la dépasse par l'imaginaire.
\par \textbf{Conclusion (rédigée) :} Le roman est bien un miroir de la société, mais un miroir actif : il reflète pour mieux éclairer, dénoncer et faire rêver. C'est ce pouvoir qui explique sa place centrale dans la littérature africaine comme universelle.
\par \textbf{Conclusion de l'exercice :} une dissertation complète exige un plan équilibré, des exemples précis et une rédaction soignée dans le temps imparti.
\end{exoresolu}

\begin{exercice}[Entraînement à la dissertation]
Sujet : \emph{La littérature n'a d'autre but que de divertir le lecteur.} Discutez cette opinion en vous appuyant sur les œuvres étudiées en classe, notamment \emph{Le Vieux nègre et la médaille}.
\par 1) Relevez les mots-clés et formulez la problématique. 2) Construisez un plan dialectique détaillé. 3) Rédigez l'introduction et la conclusion. 4) Rédigez un paragraphe du développement avec une citation commentée et une transition.
\end{exercice}

\begin{corrige}[Exercice : Entraînement à la dissertation]
\textbf{1) Mots-clés et problématique :} \emph{littérature} (ensemble des œuvres écrites) ; \emph{n'a d'autre but que} (restriction à discuter) ; \emph{divertir} (distraire, procurer du plaisir). Problématique : la littérature se réduit-elle au divertissement, ou poursuit-elle d'autres finalités ?
\par \textbf{2) Plan dialectique :} Thèse : la littérature divertit (plaisir du récit, humour, évasion). Antithèse : elle instruit et dénonce (la satire d'Oyono éveille les consciences). Synthèse : la littérature plaît et instruit à la fois, selon la formule classique \emph{plaire et instruire}.
\par \textbf{3) Introduction :} amorce (place du livre dans nos sociétés), présentation du sujet (débat entre plaisir et utilité), problématique, annonce du plan. Conclusion : bilan (la littérature divertit mais aussi enseigne et combat), prise de position, ouverture (le rôle des nouveaux médias dans le divertissement).
\par \textbf{4) Paragraphe :} idée (Oyono instruit en amusant) ; explication (la satire fait rire pour mieux faire réfléchir) ; exemple commenté (la cérémonie de la médaille, à la fois comique et pathétique) ; bilan ; transition vers la partie suivante.
\end{corrige}

\begin{attention}[Les 5 erreurs qui coûtent des points au Bac]
\begin{enumerate}
\item \textbf{Hors sujet :} ne pas traiter les mots-clés du sujet ou réciter un plan appris par cœur sans l'adapter. Toujours rattacher chaque partie à la problématique.
\item \textbf{Introduction incomplète :} oublier l'annonce du plan ou la problématique ; commencer par « De tout temps, les hommes... » (amorce vide et banale).
\item \textbf{Citation non commentée ou fausse :} une citation doit être introduite, exacte, entre guillemets, puis expliquée. Mieux vaut paraphraser que d'inventer une citation.
\item \textbf{Absence de transitions :} enchaîner les parties sans lien logique rend le devoir décousu ; chaque transition doit faire le bilan et annoncer la suite.
\item \textbf{Négliger la relecture :} fautes d'accord (sujet-verbe, participe passé), oubli des connecteurs, conclusion qui introduit une idée nouvelle ou qui ne répond pas à la problématique.
\end{enumerate}
\end{attention}

\begin{aretenir}[L'essentiel de la leçon]
Une dissertation réussie repose sur quatre piliers : une \cle{problématique} claire issue des mots-clés du sujet ; un \cle{plan} adapté (dialectique, thématique ou analytique) ; une \cle{progression thématique} assurée par la cohésion (reprises, pronoms, connecteurs) et la cohérence (une idée par paragraphe) ; enfin une \cle{introduction} en quatre étapes, un \cle{développement} en paragraphes structurés avec citations commentées et transitions, et une \cle{conclusion} en trois étapes (bilan, prise de position, ouverture).
\end{aretenir}

\newpage

\coursec{Exercices}

\coursec{Leçon 20 : Rédiger une dissertation}

\courssub{Je vérifie mes connaissances}

\begin{aretenir}[Structure canonique de la dissertation]
La dissertation comprend trois parties principales : \textbf{l'introduction}, \textbf{le développement} et \textbf{la conclusion}. L'introduction suit un ordre immuable en quatre étapes :
\begin{enumerate}
\item \textbf{L'amorce (accroche) :} citation, fait d'actualité, définition large, statistique, référence culturelle... pour capter l'attention et introduire le thème.
\item \textbf{La définition des termes et la contextualisation :} définition précise des mots-clés du sujet, délimitation du cadre (spatial, temporel, disciplinaire).
\item \textbf{La problématique :} question unique, ouverte, tensionnelle (opposition de deux thèses), issue de l'analyse du sujet. Elle guide tout le devoir.
\item \textbf{L'annonce du plan :} présentation des grandes parties (I, II, III ou A, B) en lien direct avec la problématique, sans en dévoiler le détail argumentatif.
\end{enumerate}
\end{aretenir}

\begin{exercice}[QCM : Structure de la dissertation]
Parmi les propositions suivantes, cochez celle qui correspond à l'ordre canonique des parties d'une introduction de dissertation.
\begin{enumerate}
\item \textbf{A.} Accroche -- Annonce du plan -- Problématique -- Définition des termes
\item \textbf{B.} Accroche -- Définition des termes -- Problématique -- Annonce du plan
\item \textbf{C.} Problématique -- Accroche -- Définition des termes -- Annonce du plan
\item \textbf{D.} Définition des termes -- Accroche -- Annonce du plan -- Problématique
\end{enumerate}
\end{exercice}

\begin{aretenir}[La progression thématique : cohésion et cohérence]
La \cle{progression thématique} organise l'enchaînement des thèmes (sujets des phrases) et des rhèmes (informations nouvelles) pour assurer la fluidité du texte. Trois types principaux :
\begin{itemize}
\item \textbf{Progression à thème constant (parallèle) :} même thème repris en début de phrases successives. Ex : \emph{Meka est fier. Meka attend. Meka reçoit la médaille.}
\item \textbf{Progression linéaire (en chaîne) :} le rhème d'une phrase devient le thème de la suivante. Ex : \emph{Meka reçoit une médaille. Cette médaille est un piège. Ce piège le perdra.}
\item \textbf{Progression à thème dérivé :} les thèmes sont des aspects différents d'un même hyper-thème. Ex : \emph{La cérémonie commence. L'administrateur arrive. La foule acclame.}
\end{itemize}
La \cle{cohésion} repose sur les marqueurs linguistiques (connecteurs, anaphores, reprises lexicales). La \cle{cohérence} repose sur la logique des idées et la progression thématique.
\end{aretenir}

\begin{exercice}[Vrai / Faux justifié : Progression thématique]
Indiquez si les affirmations suivantes sont vraies ou fausses. Justifiez brièvement votre réponse en mobilisant la notion de \cle{progression thématique}.
\begin{enumerate}
\item La progression thématique consiste à répéter le même mot-clé au début de chaque phrase pour assurer la cohésion.
\item Un texte cohérent peut manquer de cohésion si les connecteurs logiques sont absents.
\item La progression linéaire (ou en chaîne) relie le thème d'une phrase au rhème de la phrase précédente.
\item L'utilisation de pronoms relatifs ou personnels suffit à garantir la progression thématique d'un paragraphe argumentatif.
\end{enumerate}
\end{exercice}

\begin{exercice}[Définitions et repérage : Le Vieux Nègre et la Médaille]
Donnez une définition précise des notions suivantes en vous appuyant sur l'œuvre de Ferdinand Oyono :
\begin{enumerate}
\item La \cle{colonialisation} telle qu'elle transparaît dans le destin de Meka.
\item L'\cle{ironie} dramatique (ou situationnelle) dans la scène de la remise de la médaille.
\item Le \cle{paradoxe} de la « réussite » de Meka aux yeux de l'administration coloniale vs aux yeux de sa communauté.
\end{enumerate}
\end{exercice}

\begin{methode}[Choisir les connecteurs pour une transition]
Une transition efficace opère en trois temps : \textbf{1. Bilan} (reprise du thème de la partie achevée), \textbf{2. Pivot} (connecteur logique marquant la relation : opposition, approfondissement, cause...), \textbf{3. Annonce} (introduction du thème de la partie suivante). Le tableau ci-dessous répertorie les principaux connecteurs selon la relation logique.
\end{methode}

\begin{exercice}[Application directe : Connecteurs et transitions]
Complétez le tableau suivant en associant chaque relation logique au(x) connecteur(s) approprié(s) pour une transition entre deux paragraphes de développement.
\begin{center}
\begin{tabularx}{\linewidth}{|l|X|}\hline
\textbf{Relation logique} & \textbf{Connecteurs / Tournures de transition} \\ \hline
Addition / Renforcement & \\ \hline
Opposition / Concession & \\ \hline
Cause / Conséquence & \\ \hline
Illustration / Exemple & \\ \hline
Conclusion partielle / Synthèse & \\ \hline
\end{tabularx}
\end{center}
\end{exercice}

\courssub{Rédiger l'introduction}

\begin{methode}[Rédiger une introduction en 4 étapes (Méthode pas à pas)]
\begin{enumerate}
\item \textbf{Analyser le sujet :} Souligner les mots-clés, définir le champ notionnel, repérer la tension (opposition, évolution, évaluation).
\item \textbf{Rédiger l'amorce (2-3 lignes) :} Partir du général vers le particulier. Éviter les lieux communs (\emph{Depuis la nuit des temps...}). Préférer une citation d'auteur africain, un fait social camerounais, une statistique.
\item \textbf{Définir et contextualiser (3-4 lignes) :} Définir les termes clés (dictionnaire + sens contextuel). Situer le débat (au Cameroun, aujourd'hui, dans le cadre de l'œuvre étudiée).
\item \textbf{Formuler la problématique (1 phrase) :} Question unique, contenant une opposition (ex : \emph{Si l'école est une voie royale, n'est-elle pas aussi une impasse pour beaucoup ?}).
\item \textbf{Annoncer le plan (2-3 lignes) :} \emph{Nous montrerons d'abord que... (I). Nous verrons ensuite que... (II). Enfin, nous analyserons... (III).} Utiliser le futur ou l'infinitif.
\end{enumerate}
\end{methode}

\begin{exercice}[Rédiger une introduction complète]
Sujet type Baccalauréat technique : \emph{« L'école est-elle la seule voie de la réussite au Cameroun ? »}
Rédigez une introduction complète (environ 15--20 lignes) respectant scrupuleusement les quatre étapes : \cle{amorce} (accroche), \cle{définition des termes} / contextualisation, \cle{problématique}, \cle{annonce du plan} (en deux ou trois parties équilibrées). Soignez la progression thématique entre les phrases.
\end{exercice}

\courssub{Rédiger le développement : paragraphes, citations et transitions}

\begin{methode}[Construire un paragraphe argumentatif (Structure A.E.C.)]
\begin{enumerate}
\item \textbf{Argument (Phrase thèse) :} Une phrase affirmative, claire, qui répond directement à la problématique et annonce l'idée directrice du paragraphe.
\item \textbf{Explication / Preuve :} Développer l'argument (raisonnement, faits, données). \textbf{Insérer une citation} (courte, < 3 lignes) : l'introduire par deux-points, l'encadrer de guillemets français « ... », l'intégrer syntaxiquement (ex : \emph{Meka exprime sa confusion : « Il ne savait s'il devait rire ou pleurer ».}).
\item \textbf{Commentaire (Analyse) :} Expliquer \emph{pourquoi} et \emph{comment} la citation prouve l'argument. Analyser le sens, le style, l'ironie, le vocabulaire. Ne jamais laisser une citation « parler toute seule ».
\item \textbf{Transition (mini-synthèse + ouverture) :} Refermer le paragraphe et annoncer le suivant.
\end{enumerate}
\end{methode}

\begin{exercice}[Insertion et commentaire de citations -- Le Vieux Nègre et la Médaille]
À partir de l'extrait suivant (chapitre 3, arrivée de l'administrateur), rédigez un paragraphe argumentatif d'une quinzaine de lignes répondant à la question : \emph{« En quoi la cérémonie de la médaille révèle-t-elle l'aliénation de Meka ? »}
Votre paragraphe doit :
\begin{itemize}
\item Comporter une phrase d'ouverture (thèse).
\item Intégrer \textbf{deux citations courtes} (moins de 3 lignes chacune) prélevées dans l'extrait ou l'œuvre, insérées syntaxiquement (guillemets, deux-points, intégration dans la phrase).
\item Commenter chaque citation (analyse du sens, du style, de l'ironie).
\item Utiliser des connecteurs logiques variés pour assurer la progression thématique.
\end{itemize}
\textit{Extrait indicatif :} « Meka, tout tremblant d'émotion, s'avança... L'administrateur lui épingla la médaille sur la poitrine... Meka ne savait s'il devait rire ou pleurer. »
\end{exercice}

\begin{methode}[Rédiger une transition entre deux grandes parties]
\begin{enumerate}
\item \textbf{Reprendre (Renvoi) :} Résumer l'apport de la partie achevée (mots-clés, thèse principale).
\item \textbf{Pivoter (Connecteur) :} Utiliser un connecteur de relation logique (Cependant, Par ailleurs, En outre, Dès lors, Certes... mais).
\item \textbf{Annoncer (Ouverture) :} Introduire le thème de la partie suivante en lien avec la problématique.
\end{enumerate}
Exemple : \emph{« Au terme de cette première partie, il apparaît que la tradition structure l'identité (Bilan). \textbf{Cependant} (Opposition), force est de constater que certaines coutumes entravent les droits humains (Annonce). C'est ce que nous analyserons maintenant. »}
\end{methode}

\begin{exercice}[Rédiger une transition entre deux paragraphes]
Voici la fin d'un paragraphe (Partie I) et le début du suivant (Partie II) d'une dissertation sur le sujet : \emph{« La tradition est-elle un frein au progrès ? »}
\begin{quote}
\textbf{Fin Partie I :} ... Ainsi, les rites initiatiques comme le \textit{ngondo} ou le \textit{laakam} structurent l'identité collective et transmettent des valeurs de solidarité indispensables à la cohésion sociale.
\textbf{Début Partie II :} ... Cependant, certaines pratiques traditionnelles, comme l'excision ou le mariage forcé des mineures, portent atteinte aux droits fondamentaux de la personne humaine.
\end{quote}
\begin{enumerate}
\item Rédigez la \textbf{transition} (2 à 3 phrases) qui relie logiquement ces deux paragraphes.
\item Justifiez votre choix de connecteurs et expliquez comment la transition opère le \cle{renvoi} (reprise du thème) et l'\cle{annonce} (ouverture sur le nouveau thème).
\end{enumerate}
\end{exercice}

\courssub{Rédiger la conclusion}

\begin{methode}[Les trois temps de la conclusion]
\begin{enumerate}
\item \textbf{Bilan synthétique (3-4 lignes) :} Reprendre les thèses principales de chaque grande partie (I, II, III) sans copier-coller l'annonce du plan. Utiliser : \emph{En définitive, l'analyse a montré que... (I). Elle a également mis en lumière que... (II). Enfin, il est apparu que... (III).}
\item \textbf{Réponse nuancée à la problématique (2-3 lignes) :} Répondre directement à la question posée. Refuser le « Oui/Non » binaire. Utiliser : \emph{Ainsi, si [thèse I], toutefois [thèse II], en sorte que [synthèse].}
\item \textbf{Ouverture pertinente (2-3 lignes) :} Élargir la réflexion : autre sujet voisin, citation d'auteur (Oyono, Beti, Diome...), actualité camerounaise (réforme éducative, migration, code de la famille), perspective philosophique. Éviter l'ouverture « fourre-tout » sans lien.
\end{enumerate}
\end{methode}

\begin{exercice}[Rédiger une conclusion : Bilan, Réponse, Ouverture]
Sujet type Probatoire : \emph{« Faut-il préserver les traditions face à la modernité ? »}
On donne le plan suivant du développement :
\begin{itemize}
\item \textbf{I.} La tradition, socle identitaire et mémoire collective (exemples : langues, rites, organisation sociale).
\item \textbf{II.} Les dangers d'un traditionalisme figé : frein à l'émancipation, aux droits de l'homme, au développement technique.
\item \textbf{III.} La nécessaire dialectique : une tradition vivante, réinventée par la modernité (exemples : littérature africaine d'expression française, droit coutumier modernisé).
\end{itemize}
Rédigez la \textbf{conclusion complète} (15--20 lignes) structurée en trois temps : \cle{Bilan synthétique} (reprise des grandes lignes du développement), \cle{Réponse nuancée} à la problématique, \cle{Ouverture} pertinente (autre sujet, citation, perspective actuelle au Cameroun).
\end{exercice}

\courssub{Lecture méthodique du Vieux nègre et la médaille (n° 2 et 3) et remédiation}

\begin{aretenir}[Repères pour la lecture méthodique n°2 et 3 -- Le Vieux Nègre et la Médaille]
\begin{itemize}
\item \textbf{Lecture n°2 (Chap. 3 à 5 - La cérémonie et l'après) :} Focus sur la mise en scène du pouvoir colonial (discours de l'administrateur, symboles : drapeau, médaille, uniforme), la passivité contrainte de Meka, l'ironie du décalage entre le faste apparent et la misère réelle (la case qui tombe, la faim). Étudier le style : phrases courtes, style indirect libre, ironie narrative.
\item \textbf{Lecture n°3 (Chap. 6 à la fin - La chute) :} La médaille ne protège pas de l'arbitraire (arrestation de Kelara, spoliation des terres). La mort de Meka (abandonné, la médaille sur la poitrine) scelle l'inanité de l'honneur colonial. Thèmes : aliénation, complicité des élites (le fils missionnaire), résistance impossible ? Le titre : oxymore « Vieux Nègre » (déshumanisation) / « Médaille » (fausse reconnaissance).
\end{itemize}
\end{aretenir}

\courssub{Activités d'intégration : produire une dissertation complète}

\begin{exercice}[Sujet type Baccalauréat Technique -- Dissertation complète]
\textbf{Sujet :} \emph{« La médaille de Meka dans \textit{Le Vieux Nègre et la Médaille} : honneur ou piège ? »}
Produisez une \textbf{dissertation complète} (environ 1,5 à 2 pages) en suivant la méthode apprise.
\textbf{Consignes :}
\begin{enumerate}
\item \textbf{Analyse du sujet} (brouillon) : Définissez les termes « honneur », « piège », « médaille ». Identifiez la tension (double lecture). Formulez la problématique.
\item \textbf{Introduction} : Rédigez-la en bonne et due forme (4 étapes).
\item \textbf{Développement} en deux parties équilibrées (A et B), chacune composée de deux paragraphes argumentatifs.
\begin{itemize}
\item \textbf{Partie I :} La médaille comme reconnaissance apparente / honneur social (citations : discours de l'administrateur, fierté de la foule, chants).
\item \textbf{Partie II :} La médaille comme instrument d'aliénation / piège colonial (citations : ironie du sort de Meka, inutilité de la médaille face à la misère, récupération par l'administration, fin tragique).
\end{itemize}
\item \textbf{Transitions} : Soignez la transition I $\to$ II et les transitions intra-parties.
\item \textbf{Conclusion} : Bilan, réponse nuancée (honneur vide, piège réel), ouverture sur la littérature engagée (Oyono, Beti, Sembène).
\item \textbf{Relecture} : Vérifiez la cohérence (progression thématique), la correction linguistique (accords, ponctuation), la présence des citations.
\end{enumerate}
\end{exercice}

\begin{exercice}[Exercice de remédiation : Corriger une introduction d'élève]
Voici une introduction produite par un élève de Terminale sur le sujet : \emph{« Le travail est-il une valeur suprême ? »} Elle contient \textbf{cinq erreurs majeures} (structure, fond, forme).
\begin{quote}
« Le travail c'est quand on fait quelque chose pour gagner de l'argent. Depuis toujours les hommes travaillent. Dans le monde d'aujourd'hui le chômage est un gros problème. Est-ce que le travail est une valeur suprême ? Je vais dire que oui c'est important mais pas que. D'abord on voit que le travail permet de vivre. Ensuite que le travail peut être aliénant. Enfin la famille et la santé sont plus importantes. »
\end{quote}
\begin{enumerate}
\item Identifiez et numérotez les 5 erreurs (ex : absence d'amorce, définition floue, problématique mal formulée, annonce de plan confuse, niveau de langue familier, absence de progression thématique...).
\item Proposez une \textbf{réécriture corrigée} de cette introduction (15 lignes), en appliquant la méthode des 4 étapes et en soignant la progression thématique.
\end{enumerate}
\end{exercice}

\begin{exercice}[Sujet type Probatoire -- Production intégrée : Développement et Conclusion]
\textbf{Sujet :} \emph{« L'exil est-il une solution pour la jeunesse africaine ? »}
\textbf{Consignes :}
\begin{enumerate}
\item Proposez un \textbf{plan détaillé} (thèses / antithèses / synthèse ou causes / conséquences / perspectives) avec les arguments principaux et une citation ou référence littéraire par grand paragraphe (ex : \textit{L'Étranger} de Camus, \textit{Le Ventre de l'Atlantique} de Fatou Diome, articles de presse sur la migration).
\item Rédigez \textbf{uniquement le développement} (deux ou trois parties, 4 à 6 paragraphes au total) en insérant les citations choisies et en assurant les transitions.
\item Rédigez \textbf{la conclusion} (bilan, réponse, ouverture) à partir de votre développement.
\end{enumerate}
\textit{Note : Cet exercice évalue votre capacité à produire un texte cohérent, cohésif, argumenté et correct linguistiquement, dans les conditions de l'examen (durée conseillée : 2h30 pour l'ensemble).}
\end{exercice}

\newpage

\coursec{Corrigés des exercices}

\coursec{Rédiger une dissertation}

\begin{corrige}[Exercice 1 — QCM : Structure de la dissertation]
\textbf{Réponse correcte : B.} Amorce -- Définition des termes -- Problématique -- Annonce du plan.

\textbf{Justification :} L'introduction de dissertation suit un ordre immuable en quatre étapes. L'\cle{amorce} vient en premier pour capter l'attention du lecteur et introduire le thème général. Ensuite seulement on définit les \cle{termes du sujet} et on contextualise le débat : il serait illogique de définir les mots-clés après avoir posé la problématique, puisque celle-ci repose sur ces définitions (ce qui élimine A et C). La \cle{problématique} ne peut être formulée qu'une fois le sujet analysé et délimité ; elle ne peut donc pas ouvrir l'introduction (ce qui élimine C). Enfin, l'\cle{annonce du plan} clôt toujours l'introduction : elle découle de la problématique et ne peut la précéder (ce qui élimine A et D).

\textbf{Points de vigilance :} à l'examen, toute inversion de cet ordre est sanctionnée ; l'annonce du plan ne doit jamais dévoiler le contenu argumentatif des parties, seulement leur orientation.
\end{corrige}

\begin{corrige}[Exercice 2 — Vrai / Faux justifié : Progression thématique]
\begin{enumerate}
\item \textbf{Faux.} La \cle{progression thématique} ne consiste pas à répéter mécaniquement le même mot en tête de phrase. Même dans la progression à thème constant, le thème peut être repris par un pronom, un synonyme ou une périphrase (\emph{Meka... Le vieil homme... Il...}). De plus, il existe d'autres types de progression (linéaire, à thème dérivé) où le thème change à chaque phrase. La répétition pure et simple relèverait de la redondance, qui nuit au style.
\item \textbf{Vrai.} La \cle{cohérence} (logique des idées) et la \cle{cohésion} (marqueurs linguistiques) sont deux notions distinctes. Un texte peut être logiquement cohérent dans son enchaînement d'idées mais manquer de cohésion de surface si les connecteurs, anaphores et reprises lexicales sont absents : le lecteur doit alors reconstituer seul les liens, ce qui rend la lecture laborieuse. À l'examen, les deux sont exigées.
\item \textbf{Vrai.} C'est la définition exacte de la \cle{progression linéaire} (ou en chaîne) : le \cle{rhème} (information nouvelle) d'une phrase devient le \cle{thème} (point de départ) de la phrase suivante. Exemple : \emph{Meka reçoit une médaille. Cette médaille est un piège.} « Cette médaille », rhème de la première phrase, devient thème de la seconde.
\item \textbf{Faux.} Les pronoms (personnels, relatifs) sont des outils de \cle{cohésion} grammaticale, mais ils ne suffisent pas à garantir la progression thématique. Celle-ci suppose une organisation réfléchie de l'information : choix des thèmes successifs, lien logique entre les idées, sous-parties clairement ordonnées. Un paragraphe peut être truffé de pronoms et rester incohérent si les idées s'y succèdent sans logique.
\end{enumerate}
\textbf{Point de vigilance :} ne jamais confondre cohésion (niveau linguistique, visible) et cohérence (niveau sémantique et logique) : la justification exige de citer cette distinction.
\end{corrige}

\begin{corrige}[Exercice 3 — Définitions et repérage : Le Vieux Nègre et la Médaille]
\begin{enumerate}
\item \textbf{La colonisation} dans le destin de Meka désigne la domination politique, économique et culturelle exercée par l'administration française sur les populations camerounaises. Elle transparaît dans la vie de Meka comme une spoliation totale : ses terres ont été données à la mission, ses deux fils sont morts à la guerre pour la France, ses huttes ont été incendiées. La médaille qu'on lui décerne est la récompense dérisoire de cette dépossession : elle symbolise une « reconnaissance » qui ne restitue rien et masque la violence du système colonial.
\item \textbf{L'ironie dramatique (situationnelle)} dans la scène de la remise de la médaille réside dans le décalage entre ce que croit Meka et ce que le lecteur comprend. Meka, « tout tremblant d'émotion », croit recevoir un honneur véritable ; le lecteur, lui, perçoit que cette médaille est un instrument de propagande coloniale, une récompense vide décernée à un homme ruiné. L'ironie naît de cette double lecture : le faste de la cérémonie (drapeau, discours, uniformes) contraste avec la misère réelle du récipiendaire (sa case qui tombe en ruine, sa faim).
\item \textbf{Le paradoxe de la « réussite » de Meka} : aux yeux de l'administration coloniale, Meka est un « bon indigène », un modèle de loyalisme récompensé publiquement — c'est une réussite. Mais aux yeux de sa communauté et en réalité, il n'a rien gagné : il reste pauvre, sans terres, sans fils, et la médaille ne le protégera de rien (ni de l'arrestation, ni de la mort). Le paradoxe est donc que l'honneur suprême accordé par le colonisateur coïncide avec l'aliénation la plus complète : être célébré par ceux-là mêmes qui vous ont tout pris.
\end{enumerate}
\textbf{Point de vigilance :} chaque définition doit être illustrée par un élément précis de l'œuvre (terres, fils, cérémonie, mort) ; une définition abstraite sans ancrage textuel est insuffisante.
\end{corrige}

\begin{corrige}[Exercice 4 — Application directe : Connecteurs et transitions]
\begin{center}
\begin{tabularx}{\linewidth}{|l|X|}\hline
\rowcolor{popblueL}\textbf{Relation logique} & \textbf{Connecteurs / Tournures de transition} \\ \hline
Addition / Renforcement & De plus ; en outre ; par ailleurs ; également ; non seulement... mais encore ; qui plus est \\ \hline
Opposition / Concession & Cependant ; néanmoins ; toutefois ; en revanche ; certes... mais ; il n'en reste pas moins que \\ \hline
Cause / Conséquence & C'est pourquoi ; par conséquent ; dès lors ; ainsi ; de sorte que ; c'est la raison pour laquelle \\ \hline
Illustration / Exemple & Ainsi ; c'est le cas de ; comme le montre ; à titre d'exemple ; notamment ; c'est ainsi que \\ \hline
Conclusion partielle / Synthèse & En définitive ; au terme de cette analyse ; en somme ; finalement ; il apparaît donc que \\ \hline
\end{tabularx}
\end{center}
\textbf{Points de vigilance :}
\begin{itemize}
\item « Ainsi » peut exprimer soit la conséquence, soit l'illustration : son emploi doit être éclairci par le contexte.
\item Une transition complète ne se réduit pas à un connecteur : elle exige la structure en trois temps (bilan -- pivot -- annonce). Le connecteur n'est que le pivot.
\item Éviter les connecteurs familiers (« donc », « mais » en tête de phrase, « en gros ») dans une copie d'examen.
\end{itemize}
\end{corrige}

\begin{corrige}[Exercice 5 — Rédiger une introduction complète]
\textbf{Sujet :} « L'école est-elle la seule voie de la réussite au Cameroun ? »

\textbf{Corrigé-type (introduction rédigée) :}

\emph{[Amorce]} Au lendemain des indépendances, l'école est apparue aux Africains comme le passeport vers la modernité et l'émancipation sociale. Au Cameroun, nombre de familles consentent encore d'immenses sacrifices pour scolariser leurs enfants, convaincues que le diplôme ouvre toutes les portes.

\emph{[Définition et contextualisation]} Par « école », il faut entendre ici l'institution scolaire formelle, du primaire à l'université, sanctionnée par des diplômes. La « réussite » désigne quant à elle l'accomplissement social et économique d'un individu : emploi stable, statut social reconnu, bien-être matériel. Or, dans le Cameroun contemporain, le chômage des diplômés côtoie la prospérité d'artisans, de commerçants ou d'agriculteurs autodidactes, ce qui brouille l'équation traditionnelle entre école et succès.

\emph{[Problématique]} Dès lors, si l'école demeure une voie privilégiée d'ascension sociale, peut-on encore la considérer comme la seule voie de la réussite au Cameroun ?

\emph{[Annonce du plan]} Nous montrerons d'abord que l'école reste un instrument irremplaçable de formation et de promotion sociale. Nous verrons ensuite qu'elle ne garantit plus à elle seule la réussite, d'autres parcours — technique, artisanal, entrepreneurial — y conduisant également.

\textbf{Barème indicatif (sur 5 points) :} amorce pertinente et non banale (1 pt) ; définitions précises des deux termes-clés et contextualisation camerounaise (1,5 pt) ; problématique unique et tensionnelle (1,5 pt) ; annonce de plan claire en deux parties équilibrées (1 pt).

\textbf{Points de vigilance :} interdiction du lieu commun « Depuis la nuit des temps... » ; la problématique doit être une seule question, formulée interrogativement ; la progression thématique doit relier chaque étape à la précédente (ici : école $\to$ réussite $\to$ tension $\to$ question).
\end{corrige}

\begin{corrige}[Exercice 6 — Insertion et commentaire de citations]
\textbf{Question :} « En quoi la cérémonie de la médaille révèle-t-elle l'aliénation de Meka ? »

\textbf{Paragraphe corrigé-type (structure A.E.C.) :}

\emph{[Argument]} La cérémonie de remise de la médaille constitue le sommet de l'aliénation de Meka, car le vieil homme y célèbre, sans le comprendre, le système même qui l'a dépossédé de tout.

\emph{[Preuve 1 + commentaire]} En effet, Meka vit cet instant dans un égarement total : « Meka, tout tremblant d'émotion, s'avança ». Le syntagme « tout tremblant d'émotion » révèle une émotion subie, presque animale, qui prive le personnage de toute distance critique. L'adverbe « tout » intensifie cette perte de maîtrise : Meka n'est plus sujet de sa vie, mais objet de la mise en scène coloniale.

\emph{[Preuve 2 + commentaire]} De surcroît, le narrateur souligne sa confusion intérieure : « Meka ne savait s'il devait rire ou pleurer ». Cette hésitation entre deux réactions opposées traduit l'incompréhension profonde du personnage : il pressent confusément que cet honneur est ambigu, sans parvenir à en saisir la nature de piège. Le style indirect libre, qui fond la voix de Meka dans celle du narrateur, fait du lecteur le seul témoin lucide de cette aliénation — c'est là tout le procédé de l'ironie oyonienne.

\emph{[Transition]} Ainsi, la médaille, loin d'élever Meka, consacre sa soumission consentie. Reste à examiner comment cet honneur apparent se retourne, dans la suite du récit, en instrument de sa perte.

\textbf{Barème indicatif (sur 6 points) :} phrase-thèse claire répondant à la question (1 pt) ; deux citations courtes correctement insérées (guillemets, deux-points, intégration syntaxique) (2 pts) ; commentaire analytique de chaque citation (sens, style, ironie) (2 pts) ; connecteurs variés et progression thématique (1 pt).

\textbf{Points de vigilance :} une citation ne doit jamais « parler toute seule » : chaque citation est suivie d'une analyse ; les citations doivent être courtes (moins de 3 lignes) et exactes ; le commentaire doit viser les procédés (adverbe, style indirect libre, ironie), pas seulement paraphraser.
\end{corrige}

\begin{corrige}[Exercice 7 — Rédiger une transition entre deux paragraphes]
\textbf{1. Transition rédigée (corrigé-type) :}

« Ainsi, en structurant l'identité collective et en transmettant des valeurs de solidarité, la tradition apparaît comme un pilier irremplaçable de la cohésion sociale. \textbf{Cependant}, cette fonction positive ne doit pas faire oublier que certaines pratiques traditionnelles, loin d'épanouir l'individu, violent ses droits les plus fondamentaux. C'est à l'examen de ces dérives que nous devons maintenant nous consacrer. »

\textbf{2. Justification :}
\begin{itemize}
\item \textbf{Le renvoi (bilan) :} la première phrase reprend les mots-clés de la Partie I (« structurant l'identité collective », « valeurs de solidarité », « cohésion sociale ») et en résume la thèse : la tradition est un socle positif. Le connecteur d'ouverture « Ainsi » marque la conclusion partielle.
\item \textbf{Le pivot :} le connecteur d'opposition « Cependant » signale le basculement logique entre les deux parties : on passe d'une thèse favorable à une thèse critique, ce qui correspond exactement à la tension de la problématique (« frein au progrès ? »).
\item \textbf{L'annonce :} la deuxième et la troisième phrases introduisent le nouveau thème (les atteintes aux droits humains : excision, mariage forcé) et annoncent explicitement la Partie II (« C'est à l'examen de ces dérives que... »), en maintenant le lien avec la problématique.
\end{itemize}
\textbf{Barème indicatif (sur 4 points) :} bilan reprenant les mots-clés de la Partie I (1 pt) ; connecteur d'opposition pertinent (1 pt) ; annonce explicite du thème de la Partie II (1 pt) ; correction linguistique et fluidité (1 pt).

\textbf{Point de vigilance :} une transition ne doit ni répéter mot pour mot la fin du paragraphe précédent, ni dévoiler les arguments détaillés de la partie suivante ; elle reste à un niveau général.
\end{corrige}

\begin{corrige}[Exercice 8 — Rédiger une conclusion : Bilan, Réponse, Ouverture]
\textbf{Sujet :} « Faut-il préserver les traditions face à la modernité ? »

\textbf{Conclusion corrigée-type :}

\emph{[Bilan synthétique]} En définitive, notre analyse a d'abord montré que la tradition constitue le socle identitaire des communautés : langues, rites et organisation sociale transmettent une mémoire collective sans laquelle aucun peuple ne peut se penser. Elle a ensuite mis en lumière les dangers d'un traditionalisme figé, qui freine l'émancipation individuelle, bafoue parfois les droits humains et retarde le développement. Enfin, il est apparu qu'une tradition vivante, capable de se réinventer au contact de la modernité, offre une voie dépassant cette opposition.

\emph{[Réponse nuancée]} Ainsi, s'il faut indéniablement préserver les traditions en tant qu'héritage et repère identitaire, cette préservation ne saurait signifier leur conservation intégrale et immobile : seules les traditions réinterprétées, débarrassées de leurs aspects attentatoires à la dignité humaine, méritent d'être transmises. La vraie question n'est donc pas de choisir entre tradition et modernité, mais d'opérer un tri critique et créateur.

\emph{[Ouverture]} C'est d'ailleurs le pari de toute la littérature africaine d'expression française, de Ferdinand Oyono à Mongo Beti : écrire en langue héritée du colonisateur pour dire l'âme des traditions camerounaises. À l'heure où le Cameroun débat de la place des langues nationales dans l'enseignement, cette dialectique entre héritage et innovation demeure plus actuelle que jamais.

\textbf{Barème indicatif (sur 5 points) :} bilan reprenant les trois parties sans copier l'annonce du plan (1,5 pt) ; réponse nuancée, non binaire, à la problématique (1,5 pt) ; ouverture pertinente et reliée au sujet (1 pt) ; correction linguistique et cohésion (1 pt).

\textbf{Points de vigilance :} ne jamais répondre simplement « oui » ou « non » ; l'ouverture doit avoir un lien logique avec le sujet (pas d'ouverture « fourre-tout ») ; la conclusion n'introduit aucun argument nouveau.
\end{corrige}

\begin{corrige}[Exercice 9 — Sujet type Baccalauréat : Dissertation complète]
\textbf{Sujet :} « La médaille de Meka : honneur ou piège ? »

\textbf{1. Analyse du sujet (brouillon corrigé) :}
\begin{itemize}
\item \textbf{« Honneur » :} reconnaissance publique, distinction qui élève socialement celui qui la reçoit.
\item \textbf{« Piège » :} apparence trompeuse qui dissimule un danger ou une manipulation.
\item \textbf{« Médaille » :} symbole matériel de la récompense coloniale ; objet ambivalent.
\item \textbf{Tension :} la médaille est simultanément lue comme gloire (par Meka, la foule, l'administration) et comme instrument de domination (par le lecteur lucide).
\item \textbf{Problématique :} La médaille décernée à Meka constitue-t-elle une authentique reconnaissance de ses mérites, ou bien l'instrument raffiné de son aliénation coloniale ?
\end{itemize}

\textbf{2. Introduction (corrigé-type) :}

Dans les colonies, le pouvoir ne s'exerce pas seulement par la force : il séduit, il décore, il flatte. \emph{Le Vieux Nègre et la Médaille} (1956) de Ferdinand Oyono met en scène Meka, vieux Camerounais qui a tout sacrifié à la France — ses terres données à la mission, ses deux fils morts à la guerre — et qui reçoit en récompense une médaille. La médaille, insigne honorifique censé couronner un mérite, peut être définie comme une reconnaissance publique ; mais elle peut aussi fonctionner comme un leurre, un « piège » qui masque la domination sous les apparences de la gratitude. Dès lors, la médaille de Meka constitue-t-elle un honneur authentique ou l'instrument de son aliénation ? Nous montrerons d'abord que la médaille revêt toutes les apparences d'un honneur social ; nous verrons ensuite qu'elle n'est en réalité qu'un piège colonial.

\textbf{3. Développement (corrigé-type) :}

\textbf{Partie I — Les apparences de l'honneur.}

\emph{Paragraphe 1.} La cérémonie de remise donne à la médaille toute la solennité d'une consécration. Le discours de l'administrateur, le drapeau, les uniformes et la foule assemblée composent une mise en scène du prestige : Meka est publiquement distingué, épinglé, célébré. Cette solennité n'est pas négligeable : aux yeux de la communauté réunie, le vieil homme accède à un statut exceptionnel, celui de « grand homme » reconnu par le Blanc. La fierté de la foule et les chants qui accompagnent la cérémonie confirment que l'honneur est socialement réel, partagé, visible.

\emph{Paragraphe 2.} De plus, Meka lui-même vit cet honneur avec une sincérité poignante : « Meka, tout tremblant d'émotion, s'avança ». Son émotion atteste qu'il croit à la valeur de la distinction ; pendant un instant, le vieil homme dépossédé se sent exister aux yeux de tous. La médaille comble ainsi un besoin profond de reconnaissance, et l'on ne peut nier qu'elle procure à Meka une forme de bonheur, fût-il illusoire.

\emph{Transition.} Au terme de cette première partie, il apparaît que la médaille brille de tous les attributs de l'honneur : solennité, reconnaissance collective, fierté intime. Cependant, ces apparences résistent-elles à l'épreuve des faits ? C'est ce que nous devons examiner, car l'éclat de la cérémonie masque une réalité bien plus sombre.

\textbf{Partie II — La réalité du piège.}

\emph{Paragraphe 3.} En réalité, la médaille ne restitue rien de ce que la colonisation a pris à Meka. L'ironie du roman tient à ce décalage : on décore un homme dont la case tombe en ruine et qui a faim. La médaille est un honneur vide, sans contenu matériel ni droit attaché : elle ne rend ni les terres spoliées, ni les fils sacrifiés. Oyono souligne cette vacuité par l'ironie narrative : « Meka ne savait s'il devait rire ou pleurer », formule où l'hésitation du personnage révèle, à son insu, l'ambiguïté foncière de la récompense.

\emph{Paragraphe 4.} Pis encore, la médaille ne protège de rien : instrument de propagande, elle sert l'administration en exhibant un « bon indigène » modèle, mais elle n'oppose aucun bouclier à l'arbitraire colonial, comme le montrent l'arrestation et la violence qui frappent ensuite la communauté. La fin tragique du roman scelle cette vérité : Meka meurt abandonné, la médaille sur la poitrine. L'insigne de l'honneur devient ainsi le symbole dérisoire d'une reconnaissance qui n'a jamais existé que du côté du colonisateur. Le titre même du roman, oxymore associant « vieux nègre » (déshumanisation) et « médaille » (fausse gloire), condensait d'emblée cette duplicité.

\textbf{4. Conclusion (corrigé-type) :}

En définitive, notre analyse a montré que la médaille de Meka revêt toutes les apparences de l'honneur : solennité de la cérémonie, fierté de la foule, émotion sincère du récipiendaire. Elle a également établi que cet honneur est un leurre : récompense vide qui ne restitue rien, instrument de propagande qui ne protège de rien, jusqu'à la mort abandonnée du vieil homme. Ainsi, si la médaille est un honneur dans sa forme, elle est un piège dans son essence : elle récompense la soumission et masque la spoliation. Cette lucidité critique fait de \emph{Le Vieux Nègre et la Médaille} une œuvre engagée, dans la lignée de Mongo Beti et de Sembène Ousmane : la littérature africaine y assume sa fonction de dévoilement des mystifications coloniales, fonction qui demeure nécessaire chaque fois que le pouvoir décore au lieu de rendre justice.

\textbf{Barème indicatif (sur 20 points) :} introduction en 4 étapes (4 pts) ; problématique claire (2 pts) ; développement équilibré en deux parties, arguments étayés (6 pts) ; citations exactes, insérées et commentées (3 pts) ; transitions soignées (2 pts) ; conclusion en trois temps (2 pts) ; correction linguistique (1 pt).

\textbf{Points de vigilance :} citations exactes et courtes ; chaque citation commentée ; équilibre des deux parties ; aucune idée nouvelle en conclusion ; relecture des accords et de la ponctuation.
\end{corrige}

\begin{corrige}[Exercice 10 — Remédiation : Corriger une introduction d'élève]
\textbf{1. Les cinq erreurs majeures :}
\begin{enumerate}
\item \textbf{Absence d'amorce véritable :} la copie s'ouvre brutalement sur une définition, sans accroche (citation, fait social, constat large) destinée à capter l'attention et à introduire le thème.
\item \textbf{Définition floue et réductrice :} « Le travail c'est quand on fait quelque chose pour gagner de l'argent » réduit le travail au salariat et à l'argent ; or le travail peut être non rémunéré (travail domestique, bénévolat, travail champêtre familial). La définition doit être précise et englobante.
\item \textbf{Niveau de langue familier et impropre à l'écrit :} « c'est quand on », « un gros problème », « Je vais dire que oui c'est important mais pas que » relèvent de l'oral ; une dissertation exige un registre soutenu.
\item \textbf{Problématique mal formulée et réponse anticipée :} la question est simplement recopiée (« Est-ce que le travail est une valeur suprême ? ») sans reformulation tensionnelle, et l'élève y répond immédiatement (« oui c'est important mais pas que »), ce qui supprime toute suspense argumentative.
\item \textbf{Annonce de plan confuse et non articulée à la problématique :} « D'abord... Ensuite... Enfin... » énumère les parties sans les relier à la question, et la troisième partie (« la famille et la santé sont plus importantes ») introduit une réponse déjà tranchée ; de plus, la \cle{progression thématique} est absente : les phrases se juxtaposent sans reprise ni lien.
\end{enumerate}

\textbf{2. Réécriture corrigée (corrigé-type) :}

\emph{[Amorce]} « L'homme est né pour travailler », affirme un adage largement partagé. Au Cameroun comme ailleurs, le travail occupe une place centrale dans la vie des individus : il structure les journées, fonde les statuts sociaux et conditionne la survie matérielle des familles.

\emph{[Définition et contextualisation]} Par « travail », il faut entendre toute activité humaine, manuelle ou intellectuelle, rémunérée ou non, visant la production d'un bien ou d'un service. La notion de « valeur suprême » désigne, elle, un principe placé au-dessus de tous les autres, auquel tout devrait être subordonné. Or, dans nos sociétés marquées par le chômage des jeunes et par la recherche d'un équilibre de vie, la place accordée au travail fait débat.

\emph{[Problématique]} Dès lors, si le travail apparaît comme une condition essentielle de la dignité et de l'insertion sociale, peut-on pour autant l'élever au rang de valeur suprême, au-dessus de la santé, de la famille ou de la liberté ?

\emph{[Annonce du plan]} Nous montrerons d'abord que le travail constitue une valeur fondamentale, source de subsistance et d'épanouissement. Nous verrons ensuite qu'il peut devenir aliénant lorsqu'il est érigé en absolu. Enfin, nous analyserons la nécessité de subordonner le travail à des valeurs plus hautes, comme la santé et la vie familiale.

\textbf{Barème indicatif (sur 6 points) :} identification correcte des 5 erreurs (2,5 pts) ; réécriture respectant les 4 étapes (2,5 pts) ; registre soutenu et progression thématique (1 pt).

\textbf{Point de vigilance :} dans un exercice de remédiation, il faut non seulement corriger, mais justifier chaque correction par la règle de méthode correspondante (amorce, définition, problématique, annonce).
\end{corrige}

\begin{corrige}[Exercice 11 — Sujet type Probatoire : Production intégrée]
\textbf{Sujet :} « L'exil est-il une solution pour la jeunesse africaine ? »

\textbf{1. Plan détaillé (corrigé-type) :}
\begin{itemize}
\item \textbf{I. L'exil comme espoir légitime (thèse).} \emph{Arg. 1 :} fuir le chômage et la précarité pour trouver un emploi et soutenir sa famille (transferts d'argent). Référence : \emph{Le Ventre de l'Atlantique} de Fatou Diome (le rêve européen du jeune migrant). \emph{Arg. 2 :} accéder à des formations et à des opportunités absentes du pays d'origine.
\item \textbf{II. Les désillusions et les dangers de l'exil (antithèse).} \emph{Arg. 1 :} les périls du voyage (désert, Méditerranée) et la mort de milliers de migrants. \emph{Arg. 2 :} le déracinement, la discrimination, les emplois déqualifiés ; Diome montre que « le ventre de l'Atlantique » dévore les rêves. Référence possible : l'étranger de Camus comme figure de la solitude existentielle (à manier avec prudence).
\item \textbf{III. Perspectives : réussir chez soi ou circuler autrement (synthèse).} \emph{Arg. 1 :} l'entrepreneuriat, l'agriculture moderne et le numérique offrent des voies de réussite au pays. \emph{Arg. 2 :} la migration choisie, formée et temporaire (études, retour au pays) plutôt que l'exil subi et définitif.
\end{itemize}

\textbf{2. Développement rédigé (corrigé-type, extraits structurants) :}

\emph{[I, §1]} Il est indéniable que l'exil répond d'abord à une réalité économique implacable. Face au chômage massif des jeunes diplômés, partir apparaît comme la seule issue pour nombre d'Africains. Fatou Diome le montre dans \emph{Le Ventre de l'Atlantique} : pour le jeune héros, l'Europe incarne la terre promise où le travail serait enfin récompensé. Les transferts d'argent des migrants, qui soutiennent des familles entières, confirment que l'exil peut être une solution concrète à la pauvreté.

\emph{[I, §2]} En outre, l'exil ouvre l'accès à des formations et à des technologies souvent indisponibles sur le continent : étudier à l'étranger, c'est acquérir des compétences rares, valorisables ensuite partout.

\emph{[Transition I $\to$ II]} Ainsi, l'exil se présente d'abord comme un chemin d'émancipation économique et intellectuelle. Cependant, ce rêve affiché résiste-t-il à l'épreuve du réel ? L'expérience vécue des migrants invite à nuancer sévèrement cette première image.

\emph{[II, §1]} En effet, le voyage lui-même est devenu un chemin de croix : désert du Sahara, embarcations de fortune, Méditerranée transformée en cimetière. Des milliers de jeunes Africains paient de leur vie un espoir mal informé.

\emph{[II, §2]} De surcroît, ceux qui parviennent à destination découvrent souvent une réalité amère : emplois déqualifiés, précarité, discrimination, solitude. Diome le souligne avec ironie : le « ventre de l'Atlantique », loin de nourrir, dévore les illusions. Le migrant, déraciné de sa communauté, devient un étranger perpétuel, partageant la condition de solitude que Camus prêtait à Meursault. L'exil subi n'est donc pas une solution, mais souvent un nouveau problème.

\emph{[Transition II $\to$ III]} Au terme de cette analyse, l'exil apparaît comme un pari risqué, aux gains incertains et aux pertes certaines. Dès lors, ne convient-il pas d'interroger les ressources du continent lui-même ?

\emph{[III, §1]} En réalité, des voies de réussite existent en Afrique : entrepreneuriat, agriculture modernisée, économie numérique permettent à une jeunesse formée de bâtir chez elle. Les États, comme le Cameroun avec ses programmes d'appui à l'auto-emploi, multiplient les dispositifs en ce sens.

\emph{[III, §2]} Enfin, la mobilité elle-même peut être repensée : une migration choisie, préparée, temporaire — études ou formation suivies d'un retour — transforme l'exil subi en circulation maîtrisée, au service du développement du pays d'origine.

\textbf{3. Conclusion rédigée (corrigé-type) :}

\emph{[Bilan]} En définitive, notre réflexion a montré que l'exil naît d'un espoir légitime : échapper au chômage, accéder à la formation, soutenir les siens. Elle a ensuite établi que cet espoir se heurte à des désillusions tragiques : périls du voyage, déracinement, précarité. Enfin, il est apparu que la réussite au pays et la migration maîtrisée offrent des perspectives plus sûres.

\emph{[Réponse nuancée]} Ainsi, si l'exil peut être une solution individuelle dans des cas précis et préparés, il ne saurait constituer une réponse collective au destin de la jeunesse africaine : la vraie solution réside dans la création, sur le continent, des conditions de l'espérance.

\emph{[Ouverture]} Comme l'écrit Fatou Diome, « celui qui part ne revient jamais tout à fait » : reste à bâtir des sociétés où partir soit un choix libre, non une fatalité — tâche première des politiques de jeunesse au Cameroun et en Afrique.

\textbf{Barème indicatif (sur 20 points) :} plan cohérent et détaillé avec références (4 pts) ; développement argumenté, citations insérées et commentées (8 pts) ; transitions efficaces (3 pts) ; conclusion en trois temps (3 pts) ; correction linguistique et cohésion (2 pts).

\textbf{Points de vigilance :} chaque grand paragraphe doit comporter au moins une référence ou citation commentée ; les citations doivent être exactes et courtes ; gérer le temps (2h30) : ne pas sacrifier la conclusion ; vérifier la progression thématique entre tous les paragraphes.
\end{corrige}

\newpage

\coursec{Fiche bilan}

\begin{popfigure}
\begin{center}\tcbox[colback=poporange,colframe=poporange,coltext=white,arc=9pt,boxsep=4pt,left=6pt,right=6pt]{\bfseries\small Dissertation Terminale Tech}\end{center}
\vspace{-2pt}
\begin{tcbraster}[raster columns=3,raster equal height=rows,raster column skip=6pt,raster row skip=6pt,enhanced,blankest]
\begin{tcolorbox}[enhanced,colback=poporange,colframe=poporange,coltext=white,boxrule=0.9pt,arc=3pt,left=4pt,right=4pt,top=3pt,bottom=3pt,boxsep=1pt,halign=left]\bfseries\scriptsize 1. Analyse du sujet \& Problématique\end{tcolorbox}
\begin{tcolorbox}[enhanced,colback=popgreen,colframe=popgreen,coltext=white,boxrule=0.9pt,arc=3pt,left=4pt,right=4pt,top=3pt,bottom=3pt,boxsep=1pt,halign=left]\bfseries\scriptsize 2. Introduction (AMI : Accroche, Mise au point, Intérêt, Plan)\end{tcolorbox}
\begin{tcolorbox}[enhanced,colback=poppurple,colframe=poppurple,coltext=white,boxrule=0.9pt,arc=3pt,left=4pt,right=4pt,top=3pt,bottom=3pt,boxsep=1pt,halign=left]\bfseries\scriptsize 3. Développement Paragraphes AER Citations \& Transitions\end{tcolorbox}
\begin{tcolorbox}[enhanced,colback=poppink,colframe=poppink,coltext=white,boxrule=0.9pt,arc=3pt,left=4pt,right=4pt,top=3pt,bottom=3pt,boxsep=1pt,halign=left]\bfseries\scriptsize 4. Conclusion (Bilan, Ouverture)\end{tcolorbox}
\begin{tcolorbox}[enhanced,colback=popgold,colframe=popgold,coltext=white,boxrule=0.9pt,arc=3pt,left=4pt,right=4pt,top=3pt,bottom=3pt,boxsep=1pt,halign=left]\bfseries\scriptsize 5. Progression thématique Cohésion / Cohérence\end{tcolorbox}
\begin{tcolorbox}[enhanced,colback=popblue,colframe=popblue,coltext=white,boxrule=0.9pt,arc=3pt,left=4pt,right=4pt,top=3pt,bottom=3pt,boxsep=1pt,halign=left]\bfseries\scriptsize 6. Intégration Entraînement complet\end{tcolorbox}
\end{tcbraster}

\legende{Carte mentale de la leçon 20 : Les 6 piliers de la dissertation au Bac Technique.}
\end{popfigure}

\begin{bilan}

\courssub{Définitions clés}

\begin{definition}[Dissertation]
Exercice d'argumentation écrite où le candidat répond à une question (sujet) en développant une thèse structurée, étayée d'exemples et de citations, dans un français correct et cohérent.
\end{definition}

\begin{definition}[Problématique]
Question centrale, issue de l'analyse du sujet, à laquelle la dissertation doit répondre. Elle guide le plan et évite le hors-sujet.
\end{definition}

\begin{definition}[Progression thématique]
Enchaînement logique des idées (thèmes) d'un paragraphe à l'autre et à l'intérieur du paragraphe. Elle assure la \cle{cohérence} (liens logiques/sémantiques) et la \cle{cohésion} (connecteurs, reprises, pronoms, champs lexicaux).
\end{definition}

\begin{definition}[Paragraphe argumenté (AER)]
Unité de base du développement : \textbf{A}ffirmation (idée directrice) -- \textbf{E}xplication/Preuve (arguments, citations) -- \textbf{R}eprise/Transition (lien vers le suivant).
\end{definition}

\courssub{Structures types à connaître par cœur (Méthodes express)}

\begin{center}
\begin{tabularx}{\linewidth}{|>{\bfseries}l|X|X|}\hline
\textbf{Étape} & \textbf{Structure / Contenu} & \textbf{Connecteurs / Repères utiles} \\\hline
\textbf{Introduction (AMI)} & \textbf{A}ccroche (citation, fait, définition) $\rightarrow$ \textbf{M}ise au point (définitions, délimitation) $\rightarrow$ \textbf{I}ntérêt / Problématique $\rightarrow$ \textbf{Plan} annoncé (2 ou 3 parties) & \emph{« Il convient de se demander... » ; « Nous verrons d'abord..., ensuite..., enfin... »} \\\hline
\textbf{Paragraphe (AER)} & \textbf{A}ffirmation (phrase thèse) $\rightarrow$ \textbf{E}xplication (arguments, explications) $\rightarrow$ \textbf{E}xemple/Citation (intégrée) $\rightarrow$ \textbf{R}eprise / \textbf{T}ransition & \emph{En effet, par exemple, ainsi, par conséquent, par ailleurs, de plus} \\\hline
\textbf{Transition} & 1. Reprise de l'idée précédente $\rightarrow$ 2. Opposition / Complément / Nuance $\rightarrow$ 3. Annonce de l'idée suivante & \emph{Si... toutefois... ; Après avoir vu..., il faut examiner... ; D'un côté... de l'autre...} \\\hline
\textbf{Conclusion} & \textbf{B}ilan (réponse synthétique à la problématique) $\rightarrow$ \textbf{O}uverture (perspective, autre question, citation finale) & \emph{Au total, en définitive, force est de constater que... ; Cette étude invite à...} \\\hline
\end{tabularx}
\end{center}

\courssub{Progression thématique : Outils de cohésion \& cohérence}

\begin{itemize}
    \item \textbf{Connecteurs logiques :} addition (\emph{de plus, en outre}), conséquence (\emph{donc, par conséquent}), opposition (\emph{mais, cependant}), concession (\emph{bien que, certes}), illustration (\emph{par exemple, notamment}).
    \item \textbf{Reprises / Chaînage :} pronoms (lequel, dont, y, en), synonymes, antonymes, hyperonymes, champs lexicaux.
    \item \textbf{Parcours thématique :} thème constant (même sujet), thème dérivé (aspects différents), thème nouveau (relancé par transition).
\end{itemize}

\courssub{Intégration : \emph{Le Vieux Nègre et la Médaille} (F. Eboussi Boulaga)}
\begin{itemize}
    \item \textbf{Lecture n°2 :} Analyse des arguments sur la colonisation, l'aliénation, la dignité. Repérage de la thèse de l'auteur.
    \item \textbf{Lecture n°3 :} Étude du style (ironie, parabole, oralité), de la structure narrative et de la portée universelle.
    \item \textbf{Application :} Utiliser cet ouvrage comme \cle{référence culturelle} ou \cle{citation} dans une dissertation sur : « L'honneur », « La résistance », « Le regard de l'Autre », « Histoire et mémoire ».
\end{itemize}

\end{bilan}

\begin{aretenir}[Checklist avant le Bac -- Dissertation]
$\square$ 1. J'analyse le sujet : mots-clés, type (thèse, discussion, analyse), portée (générale / littérature).
$\square$ 2. Je formule une \textbf{problématique} claire, débattue, qui structure tout le devoir.
$\square$ 3. Je construis un \textbf{plan dialectique} (Thèse / Antithèse / Synthèse) ou thématique (2-3 parties équilibrées).
$\square$ 4. Mon \textbf{introduction} suit le modèle AMI (Accroche, Mise au point, Intérêt/Problématique, Plan annoncé).
$\square$ 5. Chaque paragraphe de développement suit \textbf{AER} : Idée directrice + Arguments + Exemple/Citation intégrée.
$\square$ 6. J'insère les citations correctement : « citation » (Auteur, Œuvre) ou Auteur écrit : « citation ».
$\square$ 7. Je rédige des \textbf{transitions} explicites entre parties et paragraphes (Reprise $\rightarrow$ Pivot $\rightarrow$ Annonce).
$\square$ 8. J'assure la \textbf{progression thématique} : connecteurs, reprises, champs lexicaux, pas de saut logique.
$\square$ 9. Ma \textbf{conclusion} fait le Bilan (réponse à la problématique) + Ouverture (perspective nouvelle).
$\square$ 10. Je relis : orthographe, grammaire, syntaxe, ponctuation, lisibilité, respect du temps imparti.
$\square$ 11. J'ai mobilisé au moins une référence du programme (\emph{Le Vieux Nègre...}, autres œuvres étudiées).
$\square$ 12. Ma copie est soignée : paragraphes sautés, marges respectées, écriture lisible.
\end{aretenir}

\findecours

\end{document}
